% Étude de fonctions numériques d'une variable réelle — Mathématiques Tle D — Cours signé OnBuch+
% Programme officiel MINESEC. Compiler avec : tectonic M03-etude-de-fonctions-numeriques.tex
\newcommand{\DOCMATIERE}{Mathématiques}
\newcommand{\DOCNIVEAU}{Tle D}
\newcommand{\DOCMODULE}{Relations et opérations fondamentales dans l'ensemble des nombres réels}
\newcommand{\DOCLECON}{03}
\newcommand{\DOCTITRE}{Étude de fonctions numériques d'une variable réelle}
% OnBuch+ — préambule « cours signés » ENRICHI (compilé avec tectonic / XeLaTeX)
% Système visuel « Pop » d'OnBuch+ : fond crème (jamais blanc), bordures encre
% épaisses, ombres « dures » décalées, titres Archivo Black en capitales.
% Version MATHÉMATIQUES : graphes (pgfplots/tikz), boîtes pédagogiques
% (définition, propriété, méthode, attention, le savais-tu, exercice résolu,
% exercice, TP, bilan), page de garde et signature OnBuch+ ancrée en bas.
%
% Macros attendues dans main.tex AVANT \input{preamble} :
%   \DOCMATIERE  \DOCNIVEAU  \DOCMODULE  \DOCLECON (numéro)  \DOCTITRE
\documentclass[11pt]{article}

\usepackage{fontspec}
\usepackage[french]{babel}
\usepackage[a4paper,margin=2cm,top=3.4cm,bottom=2.8cm,headheight=1.4cm,footskip=1.3cm]{geometry}
\usepackage{amsmath,amssymb,amsfonts}
\usepackage{mathtools} % \xmapsto, \coloneqq... souvent utilisés par les modèles
\usepackage{cancel} % simplifications barrées
% \abs{z} (module, valeur absolue) et \norm{u} (norme), utilisés par les modèles
\providecommand{\abs}{}\let\abs\relax\DeclarePairedDelimiter{\abs}{\lvert}{\rvert}
\providecommand{\norm}{}\let\norm\relax\DeclarePairedDelimiter{\norm}{\lVert}{\rVert}
\usepackage[table]{xcolor} % [table] : fournit aussi \rowcolors (alternance de lignes)
\usepackage{pagecolor}
\usepackage{tikz}
\usetikzlibrary{arrows.meta,positioning,calc,shapes.geometric,shapes.misc,shapes.arrows,decorations.pathmorphing,decorations.markings,patterns,shadows,mindmap,trees,fit,backgrounds,matrix,intersections}
\usepackage{pgfplots}
\usepackage{pgf-pie} % diagrammes circulaires \pie (statistiques)
\pgfplotsset{compat=1.18}
\usepgfplotslibrary{fillbetween,groupplots} % aires entre courbes (calcul intégral)
\usepackage{enumitem}
\usepackage{fancyhdr}
% [most] charge toutes les bibliothèques tcolorbox (listings, minted, raster,
% documentation...) et gonfle inutilement la mémoire TeX sur un document long
% (~300 boîtes) : on ne charge que ce qui est réellement utilisé.
\usepackage[skins,breakable]{tcolorbox}
\usepackage{graphicx}
\usepackage{colortbl}
\usepackage{booktabs}
\usepackage{tabularx}
\usepackage{diagbox} % case d'en-tête diagonale des tableaux à double entrée
% Colonnes extensibles centrée / gauche / droite (C, L, R), souvent utilisées par les modèles
\newcolumntype{C}{>{\centering\arraybackslash}X}
\newcolumntype{L}{>{\raggedright\arraybackslash}X}
\newcolumntype{R}{>{\raggedleft\arraybackslash}X}
\usepackage{array}
\usepackage{multicol}
\usepackage{siunitx}
% parse-numbers=false : accepte aussi \SI{2,5 \times 10^{-3}}{...}
\sisetup{locale=FR,output-decimal-marker={,},group-minimum-digits=4,parse-numbers=false}
\DeclareSIUnit\ohm{\ensuremath{\symup{\Omega}}} % U+2126 absent de la police
\usepackage{qrcode}
% \degree : commande fréquemment utilisée par les modèles pour un angle ou une
% température en dehors de \SI{}{} (non fournie par les paquets chargés).
\providecommand{\degree}{^{\circ}}
% \qed (fin de démonstration), utilisé par les modèles sans amsthm
\providecommand{\qed}{\hfill$\square$}
% \barre : double barre des valeurs interdites dans un tableau de variations (tkz-tab)
\providecommand{\barre}{\ensuremath{\|}}
% environnement proof (amsthm non chargé) utilisé par les modèles
\ifdefined\proof\else\newenvironment{proof}[1][Démonstration]{\par\noindent\textit{#1.}\ }{\qed\par}\fi
\usepackage{lastpage}
\usepackage{hyperref}
\hypersetup{hidelinks}

% ============================================================
% Palette « Pop » OnBuch+ — mêmes hex que la classe P (plus_ui.dart)
% ============================================================
\definecolor{poporange}{HTML}{F59321}
\definecolor{popdark}{HTML}{E07A0C}
\definecolor{popcream}{HTML}{FFF6E8}
\definecolor{popink}{HTML}{1C1714}
\definecolor{poppurple}{HTML}{7A5AE0}
\definecolor{popgreen}{HTML}{2E9E6A}
\definecolor{poppink}{HTML}{E0457B}
\definecolor{popblue}{HTML}{2F7DE1}
\definecolor{popgold}{HTML}{C98A12}
\definecolor{popmuted}{HTML}{8A7F76}
\definecolor{popline}{HTML}{EADFD2}
\definecolor{popgray}{HTML}{8A7F76}
\colorlet{popgrey}{popgray}
% Alias tolérants (le modèle nomme parfois les couleurs différemment)
\colorlet{popwhite}{white}
\colorlet{popblack}{popink}
\colorlet{popred}{poppink}
\colorlet{popyellow}{popgold}
% Teintes claires pour fonds de tableaux / zones
\colorlet{popblueL}{popblue!10!white}
\colorlet{poporangeL}{poporange!14!white}
\colorlet{popgreenL}{popgreen!12!white}
\colorlet{poppurpleL}{poppurple!12!white}
\colorlet{poppinkL}{poppink!12!white}
\colorlet{popgoldL}{popgold!14!white}

\pagecolor{popcream}

% ============================================================
% Typographie
% ============================================================
\defaultfontfeatures{Ligatures=TeX}
\setmainfont{PlusJakartaSans-Medium.ttf}[Path=../../../fonts/,BoldFont=PlusJakartaSans-Bold.ttf]
% Maths en sans-serif (Fira Math) avec chiffres et lettres droites de Plus
% Jakarta, pour que les formules se fondent dans le texte.
\usepackage[math-style=ISO,bold-style=ISO]{unicode-math}
\setmathfont{FiraMath-Regular.otf}[Path=../../../fonts/]
\setmathfont{PlusJakartaSans-Medium.ttf}[Path=../../../fonts/,range={up/num,up/{Latin,latin},\mathup}]
\usepackage{icomma} % virgule décimale sans espace en mode math
% Caractères mathématiques Unicode absents des polices de texte (Plus Jakarta,
% Archivo Black) : rendus par leur équivalent mathématique, en texte comme en
% formule — sinon ils s'affichent en carré vide (ex. « Arithmétique dans ℤ »).
\usepackage{newunicodechar}
\newunicodechar{ℝ}{\ensuremath{\mathbb{R}}}
\newunicodechar{ℤ}{\ensuremath{\mathbb{Z}}}
\newunicodechar{ℂ}{\ensuremath{\mathbb{C}}}
\newunicodechar{ℕ}{\ensuremath{\mathbb{N}}}
\newunicodechar{ℚ}{\ensuremath{\mathbb{Q}}}
\newunicodechar{𝔻}{\ensuremath{\mathbb{D}}}
\newunicodechar{α}{\ensuremath{\alpha}}
\newunicodechar{β}{\ensuremath{\beta}}
\newunicodechar{γ}{\ensuremath{\gamma}}
\newunicodechar{δ}{\ensuremath{\delta}}
\newunicodechar{ε}{\ensuremath{\varepsilon}}
\newunicodechar{θ}{\ensuremath{\theta}}
\newunicodechar{λ}{\ensuremath{\lambda}}
\newunicodechar{μ}{\ensuremath{\mu}}
\newunicodechar{σ}{\ensuremath{\sigma}}
\newunicodechar{τ}{\ensuremath{\tau}}
\newunicodechar{φ}{\ensuremath{\varphi}}
\newunicodechar{ψ}{\ensuremath{\psi}}
\newunicodechar{ω}{\ensuremath{\omega}}
\newunicodechar{Σ}{\ensuremath{\Sigma}}
% majuscules produites par \MakeUppercase dans les titres (α-aminés -> Α)
\newunicodechar{Α}{\ensuremath{\alpha}}
\newunicodechar{Β}{\ensuremath{\beta}}
\newunicodechar{Γ}{\ensuremath{\Gamma}}
\newunicodechar{Δ}{\ensuremath{\Delta}}
\newunicodechar{Ε}{\ensuremath{\varepsilon}}
\newunicodechar{Θ}{\ensuremath{\Theta}}
\newunicodechar{Λ}{\ensuremath{\Lambda}}
\newunicodechar{Μ}{\ensuremath{\mu}}
\newunicodechar{Τ}{\ensuremath{\tau}}
\newunicodechar{Φ}{\ensuremath{\Phi}}
\newunicodechar{Ψ}{\ensuremath{\Psi}}
\newunicodechar{Ω}{\ensuremath{\Omega}}
\newunicodechar{↦}{\ensuremath{\mapsto}}
\newunicodechar{∈}{\ensuremath{\in}}
\newunicodechar{∉}{\ensuremath{\notin}}
\newunicodechar{∧}{\ensuremath{\wedge}}
\newunicodechar{⇔}{\ensuremath{\Leftrightarrow}}
\newunicodechar{⟺}{\ensuremath{\Longleftrightarrow}}
\newunicodechar{⇒}{\ensuremath{\Rightarrow}}
\newunicodechar{⟹}{\ensuremath{\Longrightarrow}}
\newunicodechar{∘}{\ensuremath{\circ}}
\newunicodechar{∪}{\ensuremath{\cup}}
\newunicodechar{∩}{\ensuremath{\cap}}
\newunicodechar{≡}{\ensuremath{\equiv}}
\newunicodechar{⊂}{\ensuremath{\subset}}
\newunicodechar{⊥}{\ensuremath{\perp}}
\newunicodechar{∃}{\ensuremath{\exists}}
\newunicodechar{∀}{\ensuremath{\forall}}
\newunicodechar{∛}{\ensuremath{\sqrt[3]{\,}}}
\newunicodechar{∜}{\ensuremath{\sqrt[4]{\,}}}
\newunicodechar{⃗}{\ensuremath{{}^{\to}}}
\newunicodechar{ⁿ}{\ifmmode^{n}\else\textsuperscript{\ensuremath{n}}\fi}
\newunicodechar{ˣ}{\ifmmode^{x}\else\textsuperscript{\ensuremath{x}}\fi}
\newunicodechar{ᵇ}{\ifmmode^{b}\else\textsuperscript{\ensuremath{b}}\fi}
\newunicodechar{ʳ}{\ifmmode^{r}\else\textsuperscript{\ensuremath{r}}\fi}
\newunicodechar{ᵘ}{\ifmmode^{u}\else\textsuperscript{\ensuremath{u}}\fi}
\newunicodechar{ʸ}{\ifmmode^{y}\else\textsuperscript{\ensuremath{y}}\fi}
\newunicodechar{ᵃ}{\ifmmode^{a}\else\textsuperscript{\ensuremath{a}}\fi}
\newunicodechar{ᵉ}{\ifmmode^{e}\else\textsuperscript{\ensuremath{e}}\fi}
\newunicodechar{ᵏ}{\ifmmode^{k}\else\textsuperscript{\ensuremath{k}}\fi}
\newunicodechar{ᵖ}{\ifmmode^{p}\else\textsuperscript{\ensuremath{p}}\fi}
\newunicodechar{ᴺ}{\ifmmode^{N}\else\textsuperscript{\ensuremath{N}}\fi}
\newunicodechar{ᶜ}{\ifmmode^{c}\else\textsuperscript{\ensuremath{c}}\fi}
\newunicodechar{ʰ}{\ifmmode^{h}\else\textsuperscript{\ensuremath{h}}\fi}
\newunicodechar{ᵠ}{\ifmmode^{q}\else\textsuperscript{\ensuremath{q}}\fi}
\newunicodechar{ₙ}{\ifmmode_{n}\else\textsubscript{\ensuremath{n}}\fi}
\newunicodechar{ₐ}{\ifmmode_{a}\else\textsubscript{\ensuremath{a}}\fi}
\newunicodechar{ᵢ}{\ifmmode_{i}\else\textsubscript{\ensuremath{i}}\fi}
\newunicodechar{ᵣ}{\ifmmode_{r}\else\textsubscript{\ensuremath{r}}\fi}
\newunicodechar{ₖ}{\ifmmode_{k}\else\textsubscript{\ensuremath{k}}\fi}
\newunicodechar{ᵦ}{\ifmmode_{\beta}\else\textsubscript{\ensuremath{\beta}}\fi}
\newunicodechar{ₓ}{\ifmmode_{x}\else\textsubscript{\ensuremath{x}}\fi}
\newfontfamily\popdisplay{ArchivoBlack-Regular.ttf}[Path=../../../fonts/]
\newfontfamily\poplogo{SpaceGrotesk-Bold.ttf}[Path=../../../fonts/]
\newfontfamily\popsemi{PlusJakartaSans-SemiBold.ttf}[Path=../../../fonts/]
\newfontfamily\popxbold{PlusJakartaSans-ExtraBold.ttf}[Path=../../../fonts/]
\newfontfamily\popxboldls{PlusJakartaSans-ExtraBold.ttf}[Path=../../../fonts/,LetterSpace=3.0]

\setlist[enumerate]{leftmargin=1.7em,itemsep=5pt,topsep=4pt}
\setlist[itemize]{leftmargin=1.7em,itemsep=4pt,topsep=4pt,label={\color{poporange}\textbullet}}
\setlength{\parindent}{0pt}
\setlength{\parskip}{5pt}
\renewcommand{\arraystretch}{1.25}

% pgfplots : style Pop par défaut
\pgfplotsset{
  every axis/.append style={axis line style={popink,line width=1pt},tick style={popink},
    grid style={popline,line width=0.5pt},label style={font=\small},tick label style={font=\footnotesize}},
  popcurve/.style={poporange,line width=1.8pt,no markers,smooth},
}
% Même style disponible dans un \draw TikZ simple (hors pgfplots)
\tikzset{popcurve/.style={poporange,line width=1.8pt}}

% ============================================================
% Pill / squiggle
% ============================================================
\newcommand{\poppill}[3]{%
  \tikz[baseline=-0.55ex]\node[draw=popink,line width=1.2pt,rounded corners=2.6pt,%
    fill=#2,inner xsep=6.5pt,inner ysep=3pt]{%
    \popxboldls\fontsize{7}{7}\selectfont\color{#3}\MakeUppercase{#1}};%
}
\newcommand{\popsquiggle}[1]{%
  \tikz[baseline=-0.4ex]{
    \draw[#1,line width=2.6pt,line cap=round]
      plot [smooth] coordinates {(0,0.09) (0.34,0.01) (0.68,0.21) (1.02,0.01) (1.36,0.10)};
  }%
}
% Mot-clé surligné (marqueur orange sous le texte)
\newcommand{\cle}[1]{{\popxbold\color{popdark}#1}}

% ============================================================
% En-tête / pied
% ============================================================
\pagestyle{fancy}
\fancyhf{}
\renewcommand{\headrulewidth}{0pt}
\renewcommand{\footrulewidth}{0pt}
\lhead{\raisebox{-0.15ex}{\poplogo\fontsize{13}{13}\selectfont\color{popink}OnBuch\color{poporange}+}}
\rhead{\poppill{\DOCMATIERE\ \textbullet\ \DOCNIVEAU\ \textbullet\ Leçon \DOCLECON}{white}{popink}}
\lfoot{\popxbold\fontsize{7.6}{7.6}\selectfont\color{popmuted}\MakeUppercase{Cours signé OnBuch+}\ \textbullet\ {\popsemi\DOCTITRE}}
\rfoot{\tikz[baseline=-0.6ex]\node[rounded corners=8.5pt,fill=popink,minimum height=17pt,minimum width=17pt,inner xsep=5pt,inner sep=0]{\popxbold\fontsize{8}{8}\selectfont\color{popcream}\thepage\,/\,\pageref*{LastPage}};}

% ============================================================
% Titres
% ============================================================
\newcommand{\chaptitre}[2]{%
  \begin{tcolorbox}[enhanced,colback=poporange,colframe=popink,boxrule=2.6pt,arc=11pt,
      drop shadow=popink,left=15pt,right=15pt,top=12pt,bottom=11pt,before skip=3pt,after skip=18pt]
    \poppill{#2}{popink}{popcream}\par\vspace{9pt}
    {\raggedright\hyphenpenalty=10000\exhyphenpenalty=10000\popdisplay\fontsize{22}{24}\selectfont\color{popcream}\MakeUppercase{#1}\par}
  \end{tcolorbox}
}
% Section numérotée : pastille orange avec le numéro + titre Archivo
\newcounter{popsec}
\newcommand{\coursec}[1]{%
  \stepcounter{popsec}\setcounter{popsub}{0}%
  \par\vspace{16pt}\noindent
  \begin{minipage}{\linewidth}
  \tikz[baseline=-0.7ex]\node[draw=popink,line width=1.6pt,fill=poporange,rounded corners=4pt,minimum size=22pt,inner sep=0]{\popdisplay\fontsize{12}{12}\selectfont\color{popcream}\thepopsec};%
  \hspace{8pt}{\popdisplay\fontsize{15}{17}\selectfont\color{popink}\MakeUppercase{#1}}\par
  \vspace{4pt}\noindent{\color{poporange}\rule{36pt}{3.6pt}}
  \end{minipage}\par\nopagebreak\vspace{8pt}
}
\newcounter{popsub}
\newcommand{\courssub}[1]{%
  \stepcounter{popsub}%
  \par\vspace{9pt}\noindent{\popxbold\fontsize{11.5}{13}\selectfont\color{popink}{\color{poporange}\thepopsec.\thepopsub}\hspace{6pt}#1}\par\nopagebreak\vspace{4pt}
}

% ============================================================
% Boîtes pédagogiques — même recette : fond blanc, bordure épaisse colorée,
% ombre dure encre, badge plein. Titre optionnel [#1].
% ============================================================
\tcbset{popbase/.style={breakable,enhanced,colback=white,boxrule=2.3pt,arc=9pt,
  shadow={3pt}{-3pt}{0pt}{popink},left=14pt,right=14pt,top=11pt,bottom=11pt,before skip=11pt,after skip=15pt,
  fontupper=\popsemi\fontsize{10.6}{14.5}\selectfont\color{popink},
  fontlower=\popsemi\fontsize{10.6}{14.5}\selectfont\color{popink}}}
% Coche dessinée (le glyphe ✓ n'existe pas dans les polices du document)
\newcommand{\popcheck}[1]{\tikz[baseline=-0.5ex]\draw[#1,line width=1.6pt,line cap=round,line join=round](-0.13,0.0)--(-0.04,-0.09)--(0.13,0.1);}
\newcommand{\popboxhead}[3]{\poppill{#1}{#2}{white}\ifx\relax#3\relax\else\hspace{6pt}{\popxbold\fontsize{10.6}{12}\selectfont\color{popink}#3}\fi\par\vspace{7pt}}

\newtcolorbox{aretenir}[1][]{popbase,colframe=popgreen,before upper={\popboxhead{À retenir}{popgreen}{#1}}}
\newtcolorbox{exemplebox}[1][]{popbase,colframe=poporange,before upper={\popboxhead{Exemple}{poporange}{#1}}}
\newtcolorbox{definition}[1][]{popbase,colframe=popblue,before upper={\popboxhead{Définition}{popblue}{#1}}}
\newtcolorbox{propriete}[1][]{popbase,colframe=poppurple,before upper={\popboxhead{Propriété}{poppurple}{#1}}}
\newtcolorbox{methode}[1][]{popbase,colframe=popgold,before upper={\popboxhead{Méthode}{popgold}{#1}}}
\newtcolorbox{attention}[1][]{popbase,colframe=poppink,before upper={\popboxhead{Attention}{poppink}{#1}}}
\newtcolorbox{experience}[1][]{popbase,colframe=popblue,colback=popblueL,before upper={\popboxhead{Expérience}{popblue}{#1}}}
% « Le savais-tu ? » : carte violette pleine (contexte, culture, Cameroun)
\newtcolorbox{savaistu}[1][]{popbase,colback=poppurple,colframe=popink,
  fontupper=\popsemi\fontsize{10.4}{14.2}\selectfont\color{white},
  before upper={\poppill{Le savais-tu\,?}{white}{poppurple}\ifx\relax#1\relax\else\hspace{6pt}{\popxbold\color{white}#1}\fi\par\vspace{7pt}}}
% Exercice résolu : énoncé en haut, \tcblower puis la solution sur fond crème
\newcounter{popexo}\newcounter{popexr}
\newtcolorbox{exoresolu}[1][]{popbase,colframe=poporange,colbacklower=popcream,
  segmentation style={popink,line width=1.2pt,dashed},
  before upper={\stepcounter{popexr}\popboxhead{Exercice résolu \thepopexr}{poporange}{#1}},
  before lower={\poppill{Solution}{popink}{popcream}\par\vspace{6pt}}}
% Exercice (non résolu en place) — la correction est dans la partie Corrigés
\newtcolorbox{exercice}[1][]{popbase,colframe=popink,
  before upper={\stepcounter{popexo}\popboxhead{Exercice \thepopexo}{popink}{#1}}}
% Corrigé
\newtcolorbox{corrige}[1][]{popbase,colframe=popgreen,colback=popgreenL,
  before upper={\popboxhead{Corrigé}{popgreen}{#1}}}
% Objectifs / prérequis / bilan
\newtcolorbox{objectifs}{popbase,colframe=popink,colback=poporangeL,
  before upper={\popboxhead{Objectifs de la leçon}{popink}{}}}
\newtcolorbox{prerequis}{popbase,colframe=popink,colback=popblueL,
  before upper={\popboxhead{Prérequis}{popblue}{}}}
\newtcolorbox{bilan}{popbase,colframe=popink,colback=white,boxrule=2.8pt,
  before upper={\popboxhead{Fiche bilan}{poporange}{L'essentiel à connaître pour l'examen}}}
% Figure encadrée avec légende
\newtcolorbox{popfigure}[1][]{enhanced,breakable=false,colback=white,colframe=popline,boxrule=1.4pt,arc=8pt,
  left=10pt,right=10pt,top=10pt,bottom=8pt,before skip=10pt,after skip=12pt,halign=center,
  lower separated=false,halign lower=center,
  fontlower=\popsemi\fontsize{9}{11.5}\selectfont\color{popmuted},
  #1}
\newcommand{\legende}[1]{\par\vspace{6pt}{\popsemi\fontsize{9}{11.5}\selectfont\color{popmuted}#1}}

% ============================================================
% Signature OnBuch+
% ============================================================
\newcommand{\poplogoinline}[1]{{\poplogo\fontsize{#1}{#1}\selectfont\color{popink}OnBuch\color{poporange}+}}
\newcommand{\signatureonbuch}{%
  \begin{tcolorbox}[enhanced,colback=popink,colframe=popink,boxrule=2.2pt,arc=10pt,
      drop shadow=poporange,left=14pt,right=14pt,top=10pt,bottom=10pt,before skip=0pt,after skip=0pt]
    \tikz[baseline=-0.7ex]\node[circle,fill=popgreen,draw=popcream,line width=1.3pt,minimum size=22pt,inner sep=0]{\popcheck{white}};%
    \hspace{9pt}\begin{minipage}[c]{0.62\linewidth}
      {\popxbold\fontsize{12}{13}\selectfont\color{popcream}Signé\ \textbullet\ {\poplogo OnBuch\color{poporange}+}}\par\vspace{2pt}
      {\raggedright\popsemi\fontsize{8}{10.5}\selectfont\color{popline}\MakeUppercase{Équipe pédagogique OnBuch+}\\\MakeUppercase{Conforme au programme officiel MINESEC}\par}
    \end{minipage}\hfill
    {\poplogo\fontsize{22}{22}\selectfont\color{popcream}OnBuch\color{poporange}+}
  \end{tcolorbox}%
}

% ============================================================
% Page de garde
% #1 titre · #2 matière · #3 niveau · #4 module · #5 n° leçon · #6 durée ·
% #7 description / objectifs (texte ou itemize)
% ============================================================
\newcommand{\pagegarde}[7]{%
  \thispagestyle{empty}
  \newgeometry{a4paper,margin=1.7cm,top=1.6cm,bottom=1.4cm}%
  \begin{tikzpicture}[remember picture,overlay]
    \fill[poporange!18] ([xshift=-3.2cm,yshift=-3.6cm]current page.north east) circle (4.6cm);
    \fill[poppurple!14] ([xshift=2.4cm,yshift=7.6cm]current page.south west) circle (3.8cm);
  \end{tikzpicture}%
  {\poplogo\fontsize{30}{30}\selectfont\color{popink}OnBuch\color{poporange}+}\hfill
  \raisebox{0.6ex}{\poppill{Cours signé \textbullet\ Premium}{popink}{popcream}}\par
  \vspace{1pt}\popsquiggle{poppurple}\par
  \vspace{8pt}
  \begin{tcolorbox}[enhanced,colback=poporange,colframe=popink,boxrule=2.8pt,arc=13pt,
      drop shadow=popink,left=18pt,right=18pt,top=12pt,bottom=13pt,after skip=10pt]
    \poppill{#2 \textbullet\ #3}{popink}{popcream}\hspace{5pt}\poppill{#4}{white}{popink}\par\vspace{8pt}
    {\popdisplay\fontsize{14}{14}\selectfont\color{popink}LEÇON #5}\par\vspace{4pt}
    {\raggedright\hyphenpenalty=10000\exhyphenpenalty=10000\popdisplay\fontsize{25}{27}\selectfont\color{popcream}\MakeUppercase{#1}\par}
    \vspace{8pt}
    {\popxbold\fontsize{9.5}{9.5}\selectfont\color{popink}\MakeUppercase{Durée conseillée : #6}}
  \end{tcolorbox}
  \begin{tcolorbox}[enhanced,colback=white,colframe=popink,boxrule=2.2pt,arc=10pt,
      drop shadow=popink,left=16pt,right=16pt,top=10pt,bottom=10pt,after skip=8pt]
    \poppill{Dans cette leçon}{poppurple}{white}\par\vspace{5pt}
    \popsemi\fontsize{9.3}{12.6}\selectfont\color{popink}#7
  \end{tcolorbox}
  \noindent\begin{minipage}[t]{0.487\linewidth}
    \begin{tcolorbox}[enhanced,colback=popgreen,colframe=popink,boxrule=2.2pt,arc=10pt,
        drop shadow=popink,left=11pt,right=11pt,top=10pt,bottom=10pt,height=3.1cm,valign=center]
      \tcbox[colback=white,colframe=popink,boxrule=1.3pt,arc=5pt,boxsep=0pt,nobeforeafter,
        left=3pt,right=3pt,top=3pt,bottom=3pt,valign=center]{\qrcode[height=1.9cm]{https://play.google.com/store/apps/details?id=com.luvvix.onbuch&hl=fr}}%
      \hspace{8pt}\begin{minipage}[c]{0.5\linewidth}
        {\popdisplay\fontsize{11}{12.5}\selectfont\color{white}\MakeUppercase{Télécharge OnBuch}}\par\vspace{4pt}
        {\raggedright\popsemi\fontsize{8.4}{11}\selectfont\color{white}Gratuit \textbullet\ Google Play\\Tuteur IA, annales, concours, résultats\par}
      \end{minipage}
    \end{tcolorbox}
  \end{minipage}\hfill
  \begin{minipage}[t]{0.487\linewidth}
    \begin{tcolorbox}[enhanced,colback=poppurple,colframe=popink,boxrule=2.2pt,arc=10pt,
        drop shadow=popink,left=11pt,right=11pt,top=10pt,bottom=10pt,height=3.1cm,valign=center]
      \tikz[baseline=-0.7ex]\node[circle,fill=white,draw=popink,line width=1.3pt,minimum size=1.6cm,inner sep=0]{\popdisplay\fontsize{24}{24}\selectfont\color{poppurple}+};%
      \hspace{8pt}\begin{minipage}[c]{0.6\linewidth}
        {\popdisplay\fontsize{11}{12.5}\selectfont\color{white}\MakeUppercase{Passe à OnBuch+}}\par\vspace{4pt}
        {\raggedright\popsemi\fontsize{8.4}{11}\selectfont\color{white}Tous les cours signés, Tuteur IA illimité, fiches PDF — onglet OnBuch+ de l'app\par}
      \end{minipage}
    \end{tcolorbox}
  \end{minipage}
  \vspace{8pt}
  \popcontenu
  \vfill
  \signatureonbuch
  \restoregeometry
  \newpage
}

% Rangée « Ce cours contient » (page de garde)
\newcommand{\poptile}[3]{% #1 couleur #2 gros texte #3 légende
  \begin{tcolorbox}[enhanced,colback=#1,colframe=popink,boxrule=2pt,arc=9pt,shadow={2.5pt}{-2.5pt}{0pt}{popink},
      left=6pt,right=6pt,top=9pt,bottom=9pt,halign=center,valign=center,height=1.75cm,width=\linewidth,nobeforeafter]
    {\popdisplay\fontsize{12.5}{13}\selectfont\color{white}\MakeUppercase{#2}}\par\vspace{3pt}
    {\centering\popsemi\fontsize{7.8}{9.5}\selectfont\color{white}#3\par}
  \end{tcolorbox}}
\newcommand{\popcontenu}{%
  \par\noindent{\popxboldls\fontsize{7.5}{7.5}\selectfont\color{popmuted}CE COURS SIGNÉ CONTIENT}\par\vspace{7pt}
  \noindent\begin{minipage}[t]{0.235\linewidth}\poptile{popblue}{Cours}{complet, illustré, pas à pas}\end{minipage}\hfill
  \begin{minipage}[t]{0.235\linewidth}\poptile{poporange}{Méthodes}{et exercices résolus}\end{minipage}\hfill
  \begin{minipage}[t]{0.235\linewidth}\poptile{poppink}{Exercices}{type examen + corrigés}\end{minipage}\hfill
  \begin{minipage}[t]{0.235\linewidth}\poptile{popgreen}{Fiche bilan}{l'essentiel à retenir}\end{minipage}\par}

% Sommaire
\newtcolorbox{sommaire}{enhanced,colback=white,colframe=popink,boxrule=2.2pt,arc=10pt,
  drop shadow=popink,left=18pt,right=18pt,top=13pt,bottom=13pt,before skip=6pt,after skip=20pt,
  before upper={\poppill{Sommaire}{poppurple}{white}\par\vspace{9pt}\popsemi\fontsize{10.6}{14.5}\selectfont\color{popink}}}

% Signature de fin de cours — poussée tout en bas de la dernière page
\newcommand{\findecours}{%
  \par\vfill\nopagebreak
  \begin{center}\popsquiggle{poporange}\end{center}\vspace{4pt}
  \signatureonbuch
}
\AtBeginDocument{\def\llbracket{\mathopen{[\mkern-3.5mu[}}\def\rrbracket{\mathclose{]\mkern-3.5mu]}}} % intervalles d'entiers (glyphe ⟦ absent de la police)
% Glyphes absents de Fira Math : redéfinis à partir de glyphes disponibles
\AtBeginDocument{%
  \def\setminus{\mathbin{\backslash}}\let\smallsetminus\setminus
  \def\vdots{\mathord{\vbox{\baselineskip3.2pt\lineskiplimit0pt\kern4pt\hbox{.}\hbox{.}\hbox{.}}}}%
  \def\ddots{\mathinner{\mkern1mu\raise6.5pt\hbox{.}\mkern2mu\raise3.6pt\hbox{.}\mkern2mu\raise0.7pt\hbox{.}\mkern1mu}}%
  \def\triangle{\mathord{\scalebox{0.9}{$\symup{\Delta}$}}}%
  \def\nabla{\mathord{\raisebox{\depth}{\scalebox{1}[-1]{$\symup{\Delta}$}}}}%
  \def\checkmark{\popcheck{popdark}}%
  \def\star{\mathbin{\ast}}\def\bigstar{\mathord{\ast}}%
  \def\diamond{\mathbin{\scalebox{0.6}{$\Diamond$}}}%
}

%% FIGFIX-BEGIN v4 — correctifs globaux des figures (docs/onbuch-generation/tools/figfix)
\usepackage{adjustbox}\usepackage{etoolbox}
% toute figure TikZ/pgfplots plus large que la ligne est réduite à la largeur disponible
\BeforeBeginEnvironment{tikzpicture}{\begin{adjustbox}{max width=\linewidth}}
\AfterEndEnvironment{tikzpicture}{\end{adjustbox}}
% légendes pgfplots lisibles par-dessus les courbes
\pgfplotsset{every axis/.append style={legend style={fill=white,fill opacity=0.92,text opacity=1,draw=black!15,inner sep=3pt}}}
% étiquettes d'arêtes (syntaxe edge["..."]) avec fond blanc
\tikzset{every edge quotes/.append style={fill=white,inner sep=1.2pt}}
%% FIGFIX-END
\begin{document}

\pagegarde{Étude de fonctions numériques d'une variable réelle}{Mathématiques}{Tle D}{Relations et opérations fondamentales dans l'ensemble des nombres réels}{03}{14 h}{Cette leçon achève l'étude des fonctions numériques d'une variable réelle : théorème des valeurs intermédiaires, fonctions réciproques, racines n-ièmes, inégalité des accroissements finis et étude des branches infinies. Elle fournit tous les outils pour mener une étude complète de fonction et tracer sa courbe, compétence centrale de l'épreuve de mathématiques au Baccalauréat.
\vspace{4pt}
\begin{multicols}{2}\small
\begin{itemize}
\item À la fin de cette leçon, je sais déterminer l'image d'un intervalle par une fonction continue
\item À la fin de cette leçon, je sais démontrer l'existence et l'unicité d'une solution de l'équation f(x) = c grâce au théorème des valeurs intermédiaires
\item À la fin de cette leçon, je sais justifier qu'une fonction est bijective, définir sa fonction réciproque et calculer la dérivée de celle-ci
\item À la fin de cette leçon, je sais résoudre l'équation xⁿ = a et simplifier des expressions avec des puissances rationnelles
\item À la fin de cette leçon, je sais utiliser l'inégalité des accroissements finis pour encadrer f(b) - f(a) ou approcher une valeur
\item À la fin de cette leçon, je sais déterminer les branches infinies d'une fonction et les équations de ses asymptotes (verticales, horizontales, obliques)
\end{itemize}\end{multicols}}

\begin{sommaire}
\begin{multicols}{2}
\begin{enumerate}
\item Image d'un intervalle par une fonction continue et théorème des valeurs intermédiaires
\item Fonction réciproque d'une bijection
\item Racines n-ièmes et puissances rationnelles
\item Inégalité des accroissements finis
\item Limites usuelles et branches infinies
\item Étude complète de fonctions et tracé de courbes
\item Activité d'intégration
\item Méthodes et exercices résolus
\item Exercices
\item Corrigés des exercices
\item Fiche bilan
\end{enumerate}
\end{multicols}
\end{sommaire}

\begin{objectifs}
\begin{itemize}
\item À la fin de cette leçon, je sais déterminer l'image d'un intervalle par une fonction continue
\item À la fin de cette leçon, je sais démontrer l'existence et l'unicité d'une solution de l'équation f(x) = c grâce au théorème des valeurs intermédiaires
\item À la fin de cette leçon, je sais justifier qu'une fonction est bijective, définir sa fonction réciproque et calculer la dérivée de celle-ci
\item À la fin de cette leçon, je sais résoudre l'équation xⁿ = a et simplifier des expressions avec des puissances rationnelles
\item À la fin de cette leçon, je sais utiliser l'inégalité des accroissements finis pour encadrer f(b) - f(a) ou approcher une valeur
\item À la fin de cette leçon, je sais déterminer les branches infinies d'une fonction et les équations de ses asymptotes (verticales, horizontales, obliques)
\item À la fin de cette leçon, je sais mener l'étude complète d'une fonction rationnelle, irrationnelle ou trigonométrique et tracer sa courbe représentative
\item À la fin de cette leçon, je sais exploiter la symétrie des courbes de f et de sa réciproque par rapport à la droite y = x
\end{itemize}
\end{objectifs}

\begin{prerequis}
\begin{itemize}
\item Notion de continuité d'une fonction en un point et sur un intervalle (Première)
\item Dérivation : dérivées usuelles, sens de variation et signe de la dérivée (Première)
\item Calcul de limites : opérations, formes indéterminées, croissances comparées (début de Terminale)
\item Fonctions trigonométriques : valeurs remarquables, dérivées de sin, cos, tan (Première/Terminale)
\item Équations de droites et lecture graphique d'une courbe représentative
\end{itemize}
\end{prerequis}

\newpage

\coursec{Image d'un intervalle par une fonction continue et théorème des valeurs intermédiaires}

\textbf{Situation d'accroche.} À Douala, un opérateur de téléphonie mobile suit l'évolution de son parc d'abonnés. Le nombre d'abonnés, exprimé en milliers, est modélisé par
\[ f(t) = t^{3} - 3t + 5, \]
où $t$ désigne le temps écoulé en mois depuis le lancement d'une offre. Deux questions se posent à lui. D'abord : à quel moment le nombre d'abonnés atteindra-t-il \emph{exactement} \num{4000}, c'est-à-dire $f(t) = 4$ ? Ensuite : pourra-t-il \og remonter le temps \fg{}, c'est-à-dire, connaissant le nombre d'abonnés, retrouver le mois correspondant ? Par ailleurs, un ingénieur de la SONAREL doit estimer $\sqrt[3]{10}$ pour dimensionner une cuve, sans calculatrice : comment prouver que ce nombre existe et comment l'encadrer ?

\textbf{Problématique.} Comment garantir l'\emph{existence} et l'\emph{unicité} de la solution d'une équation du type $f(x) = c$ ? Comment localiser cette solution avec une précision donnée ? Les réponses reposent sur une propriété fondamentale des fonctions continues : le \cle{théorème des valeurs intermédiaires}.

\courssub{Rappel : continuité sur un intervalle}

\begin{definition}[Continuité sur un intervalle]
Soit $f$ une fonction définie sur un intervalle $I$.
\begin{itemize}
\item $f$ est \cle{continue en un point} $a$ de $I$ si $\displaystyle\lim_{x \to a} f(x) = f(a)$.
\item $f$ est \cle{continue sur} $I$ si elle est continue en tout point de $I$.
\end{itemize}
\end{definition}

Graphiquement, une fonction continue sur $I$ est une fonction dont la courbe se trace \og sans lever le crayon \fg{} sur $I$ : il n'y a ni saut, ni trou, ni rupture.

\begin{aretenir}[Fonctions continues de référence]
Sont continues sur tout intervalle de leur ensemble de définition : les fonctions \cle{polynômes} (sur $\mathbb{R}$), les fonctions \cle{rationnelles} (sur tout intervalle ne contenant pas de pôle), les fonctions $\sin$ et $\cos$ (sur $\mathbb{R}$), la fonction $\sqrt{x}$ (sur $[0\,;+\infty[$), ainsi que toute somme, produit, quotient (là où le dénominateur ne s'annule pas) et composée de fonctions continues.
\end{aretenir}

\begin{exemplebox}[Contre-exemple : la fonction partie entière]
La fonction \cle{partie entière} $E$, qui à tout réel $x$ associe le plus grand entier inférieur ou égal à $x$, est \textbf{discontinue en tout point entier}. Par exemple :
\[ \lim_{\substack{x \to 2 \\ x < 2}} E(x) = 1 \quad \text{et} \quad E(2) = 2. \]
Sa courbe présente des \og marches d'escalier \fg{} : elle saute de $1$ à $2$ en $x = 2$. C'est le prototype de fonction à laquelle le théorème des valeurs intermédiaires \textbf{ne s'applique pas}.
\end{exemplebox}

\courssub{Image d'un intervalle par une fonction continue}

\begin{definition}[Image d'un intervalle par une fonction]
Soit $f$ une fonction définie sur un intervalle $I$. L'\cle{image de $I$ par $f$}, notée $f(I)$, est l'ensemble de toutes les valeurs prises par $f$ quand $x$ décrit $I$ :
\[ f(I) = \{\, f(x) \; ; \; x \in I \,\}. \]
\end{definition}

\begin{propriete}[Image d'un intervalle par une fonction continue (admis)]
L'image d'un \textbf{intervalle} par une fonction \textbf{continue} est un \textbf{intervalle}.
\end{propriete}

Ce théorème, que nous admettons, dit une chose très intuitive : si la courbe se trace sans lever le crayon, alors $f$ ne peut pas \og sauter \fg{} des valeurs ; elle atteint toutes les valeurs intermédiaires, donc l'ensemble des valeurs prises ne comporte aucun trou.

\begin{propriete}[Image d'un segment (admis)]
Soit $f$ une fonction continue sur un \textbf{segment} $[a\,;b]$ (intervalle fermé et borné). Alors :
\begin{itemize}
\item $f$ est \cle{bornée} sur $[a\,;b]$ et \cle{atteint ses bornes} : il existe $x_{m}$ et $x_{M}$ dans $[a\,;b]$ tels que
\[ m = f(x_{m}) = \min_{[a\,;b]} f \qquad \text{et} \qquad M = f(x_{M}) = \max_{[a\,;b]} f \,; \]
\item l'image du segment est le segment $f\big([a\,;b]\big) = [m\,;M]$.
\end{itemize}
\end{propriete}

\begin{attention}[Un intervalle ouvert n'a pas toujours pour image un intervalle fermé]
Le résultat $f([a\,;b]) = [m\,;M]$ est propre aux \textbf{segments}. Par exemple, $f(x) = \dfrac{1}{x}$ est continue sur $I = {]0\,;1]}$, mais $f(I) = [1\,;+\infty[$, qui n'est pas borné. De même, l'image de $]0\,;+\infty[$ par $x \mapsto \dfrac{1}{x}$ est $]0\,;+\infty[$ : les bornes ouvertes se traduisent par des \textbf{limites}, non par des valeurs atteintes.
\end{attention}

Lorsque $f$ est \textbf{continue et strictement monotone}, on détermine l'image d'un intervalle en lisant le tableau de variations, selon le tableau récapitulatif suivant (on note $\ell = \lim_{a^{+}} f$, $L = \lim_{b^{-}} f$, et $\lambda = \lim_{+\infty} f$).

\begin{center}
\begin{tabularx}{\linewidth}{|l|X|X|}
\hline
\rowcolor{popblueL} \textbf{Intervalle $I$} & \textbf{$f$ strictement croissante} & \textbf{$f$ strictement décroissante} \\
\hline
$[a\,;b]$ & $[f(a)\,;f(b)]$ & $[f(b)\,;f(a)]$ \\
\hline
$[a\,;b[$ & $[f(a)\,;L[$ & $]L\,;f(a)]$ \\
\hline
$]a\,;b]$ & $]\ell\,;f(b)]$ & $[f(b)\,;\ell[$ \\
\hline
$[a\,;+\infty[$ & $[f(a)\,;\lambda[$ & $]\lambda\,;f(a)]$ \\
\hline
$\mathbb{R}$ & $]\lim_{-\infty} f\,;\lim_{+\infty} f[$ & $]\lim_{+\infty} f\,;\lim_{-\infty} f[$ \\
\hline
\end{tabularx}
\end{center}

\begin{exemplebox}[Lecture d'images d'intervalles]
Soit $f(x) = x^{3} - 3x + 5$. On a $f'(x) = 3x^{2} - 3 = 3(x-1)(x+1)$, donc $f$ est strictement croissante sur $[1\,;+\infty[$, avec $f(1) = 3$ et $\lim_{+\infty} f = +\infty$. Ainsi :
\[ f\big([1\,;+\infty[\big) = [3\,;+\infty[. \]
Sur le segment $[0\,;1]$, $f$ est strictement décroissante, donc $f([0\,;1]) = [f(1)\,;f(0)] = [3\,;5]$.
\end{exemplebox}

\courssub{Théorème des valeurs intermédiaires}

Voici le théorème central de cette section. L'intuition est celle d'un automobiliste sur la route Douala--Yaoundé : s'il part de l'altitude $f(a)$ et arrive à l'altitude $f(b)$ sans quitter la route (continuité), alors il est \emph{forcément} passé par toutes les altitudes intermédiaires.

\begin{propriete}[Théorème des valeurs intermédiaires (TVI) — admis]
Soit $f$ une fonction \textbf{continue} sur un segment $[a\,;b]$. Pour tout réel $c$ compris entre $f(a)$ et $f(b)$, l'équation $f(x) = c$ admet \textbf{au moins une} solution dans $[a\,;b]$ : il existe $\alpha \in [a\,;b]$ tel que $f(\alpha) = c$.
\end{propriete}

\begin{popfigure}
\begin{tikzpicture}
\begin{axis}[width=0.85\linewidth, height=6cm, axis lines=middle, xlabel={$x$}, ylabel={$y$},
  xmin=-0.5, xmax=5.5, ymin=-0.5, ymax=5.5, xtick={1,4.6}, xticklabels={$a$,$b$},
  ytick={1,2.6,4.4}, yticklabels={$f(a)$,$c$,$f(b)$}, clip=false]
\addplot[popcurve, domain=1:4.6, samples=100] {0.75*x + 0.35*sin(deg(2.2*x)) + 0.15};
\addplot[poporange, dashed, domain=0.3:5.2] {2.6};
\addplot[popgreen, dashed] coordinates {(2.9,0) (2.9,2.6)};
\addplot[only marks, mark=*, popink] coordinates {(2.9,2.6)};
\node[popgreen, below] at (axis cs:2.9,0) {$\alpha$};
\node[popink, above right] at (axis cs:4.2,4.1) {$\mathcal{C}_{f}$};
\node[poporange, right] at (axis cs:5.2,2.6) {$y=c$};
\end{axis}
\end{tikzpicture}
\legende{La droite $y = c$, comprise entre $f(a)$ et $f(b)$, rencontre forcément la courbe continue $\mathcal{C}_{f}$ : il existe $\alpha \in [a\,;b]$ avec $f(\alpha) = c$.}
\end{popfigure}

\begin{attention}[La continuité est indispensable]
Si $f$ n'est \textbf{pas continue}, la conclusion peut tomber en défaut : la droite $y = c$ peut passer \og à travers \fg{} le saut de la courbe sans la rencontrer. Avant d'appliquer le TVI, \textbf{toujours justifier la continuité} de $f$ sur l'intervalle considéré.
\end{attention}

\begin{popfigure}
\begin{tikzpicture}
\begin{axis}[width=0.85\linewidth, height=6cm, axis lines=middle, xlabel={$x$}, ylabel={$y$},
  xmin=-0.5, xmax=5.5, ymin=-0.5, ymax=5.5, xtick={1,2.8,4.6}, xticklabels={$a$,$x_{0}$,$b$},
  ytick={1,2.5,4.3}, yticklabels={$f(a)$,$c$,$f(b)$}, clip=false]
\addplot[popcurve, domain=1:2.79, samples=60] {0.8*x + 0.2};
\addplot[popcurve, domain=2.81:4.6, samples=60] {0.8*x + 1.05};
\addplot[only marks, mark=*, popink] coordinates {(2.8,2.44)};
\addplot[only marks, mark=o, popink] coordinates {(2.8,3.29)};
\addplot[poporange, dashed, domain=0.3:5.2] {2.5};
\node[poporange, right] at (axis cs:5.2,2.5) {$y=c$};
\node[popink, above right] at (axis cs:3.6,4.1) {$\mathcal{C}_{f}$};
\end{axis}
\end{tikzpicture}
\legende{Contre-exemple : $f$ est discontinue en $x_{0}$ (saut). Bien que $c$ soit compris entre $f(a)$ et $f(b)$, la droite $y = c$ ne rencontre pas la courbe : l'équation $f(x) = c$ n'a aucune solution.}
\end{popfigure}

\begin{experience}[Idée de la démonstration : la dichotomie]
Pourquoi le TVI est-il vrai ? Discutons-en en classe sur un cas concret : supposons $f$ continue sur $[a\,;b]$, $f(a) < c < f(b)$, et construisons une suite d'intervalles emboîtés qui \og piègent \fg{} une solution.
\begin{enumerate}
\item On pose $a_{0} = a$, $b_{0} = b$. On considère le milieu $m_{0} = \dfrac{a_{0}+b_{0}}{2}$.
\item Si $f(m_{0}) \leq c$, la solution est à droite : on garde $[a_{1}\,;b_{1}] = [m_{0}\,;b_{0}]$. Sinon, on garde $[a_{1}\,;b_{1}] = [a_{0}\,;m_{0}]$.
\item On recommence : à chaque étape, $f(a_{n}) \leq c \leq f(b_{n})$ et la longueur de l'intervalle est divisée par $2$ : $b_{n} - a_{n} = \dfrac{b-a}{2^{n}}$.
\item Les suites $(a_{n})$ croissante et $(b_{n})$ décroissante sont adjacentes ; elles convergent vers une même limite $\alpha$. Par continuité, $f(a_{n}) \to f(\alpha)$ et $f(b_{n}) \to f(\alpha)$, et l'encadrement $f(a_{n}) \leq c \leq f(b_{n})$ donne à la limite $f(\alpha) = c$.
\end{enumerate}
Cet algorithme, appelé \cle{dichotomie}, fournit en prime une méthode effective d'encadrement de la solution.
\end{experience}

\begin{popfigure}
\begin{tikzpicture}[scale=1]
\draw[popline, thick] (0,0) -- (10,0);
\foreach \x/\lab in {0/{$a_{0}$},10/{$b_{0}$}} \draw[popink, thick] (\x,0.12) -- (\x,-0.12) node[below=2pt] {\lab};
\draw[popblue, thick] (5,0.5) -- (5,-0.5);
\draw[popblue, thick] (5,0.55) -- (10,0.55); \node[popblue, above] at (7.5,0.55) {$[a_{1}\,;b_{1}]$};
\draw[popgreen, thick] (5,1.05) -- (7.5,1.05); \node[popgreen, above] at (6.25,1.05) {$[a_{2}\,;b_{2}]$};
\draw[poporange, thick] (6.25,1.55) -- (7.5,1.55); \node[poporange, above] at (6.9,1.55) {$[a_{3}\,;b_{3}]$};
\draw[poppink, thick, ->] (6.9,2.3) -- (6.9,1.85); \node[poppink, above] at (6.9,2.3) {$\alpha$};
\draw[poppink, dashed] (6.9,1.55) -- (6.9,0);
\end{tikzpicture}
\legende{Dichotomie : les intervalles emboîtés $[a_{n}\,;b_{n}]$, dont la longueur est divisée par $2$ à chaque étape, resserrent l'étau autour de la solution $\alpha$.}
\end{popfigure}

\courssub{Corollaire : le théorème de la bijection}

Le TVI donne l'\emph{existence} d'une solution, mais pas son \emph{unicité} : une courbe continue peut traverser plusieurs fois la droite $y = c$. Pour garantir l'unicité, on ajoute une hypothèse : la \textbf{stricte monotonie}.

\begin{propriete}[Corollaire du TVI : théorème de la bijection]
Soit $f$ une fonction \textbf{continue} et \textbf{strictement monotone} sur un intervalle $I$. Alors pour tout réel $c$ appartenant à l'intervalle image $f(I)$, l'équation $f(x) = c$ admet une \textbf{unique} solution dans $I$.
\end{propriete}

\begin{experience}[Démonstration guidée (exigible : savoir la rédiger)]
Soit $c \in f(I)$. Raisonnons en deux temps.
\begin{enumerate}
\item \textbf{Existence.} Puisque $c \in f(I)$, il existe par définition de l'image un réel $\alpha \in I$ tel que $f(\alpha) = c$. (Sur un segment $[a\,;b] \subset I$ avec $c$ entre $f(a)$ et $f(b)$, c'est exactement le TVI qui fournit cet $\alpha$.)
\item \textbf{Unicité.} Supposons que $f$ soit strictement croissante (le cas décroissant est analogue) et qu'il existe deux solutions $\alpha$ et $\alpha'$ dans $I$ avec $\alpha < \alpha'$. Par stricte croissance, $f(\alpha) < f(\alpha')$, ce qui contredit $f(\alpha) = f(\alpha') = c$. Donc $\alpha = \alpha'$ : la solution est unique.
\end{enumerate}
\textbf{Conclusion :} l'existence vient du TVI (continuité), l'unicité vient de la stricte monotonie. Retiens bien cette répartition des rôles.
\end{experience}

\begin{methode}[Résoudre $f(x) = c$ : existence et unicité]
\begin{enumerate}
\item Justifier que $f$ est \textbf{continue} sur $I$ (fonction polynôme, rationnelle, composée, etc.).
\item Étudier la \textbf{stricte monotonie} de $f$ sur $I$ (signe de $f'$).
\item Dresser le \textbf{tableau de variations complet} avec les limites aux bornes, et en déduire l'intervalle image $f(I)$.
\item Vérifier que $c \in f(I)$.
\item Conclure : \og D'après le corollaire du théorème des valeurs intermédiaires, l'équation $f(x) = c$ admet une unique solution $\alpha$ dans $I$. \fg{}
\end{enumerate}
\end{methode}

\begin{attention}[Les deux pièges classiques du Bac]
\begin{itemize}
\item \textbf{Piège 1 :} appliquer le TVI sans avoir vérifié (et rédigé !) la continuité de $f$. Sans continuité, le théorème ne s'applique pas : voir le contre-exemple de la fonction en escalier.
\item \textbf{Piège 2 :} conclure à l'\textbf{unicité} avec le seul TVI. Le TVI seul ne donne que l'existence d'\emph{au moins une} solution. L'unicité exige la \textbf{stricte monotonie}. Si $f$ n'est monotone que par morceaux, applique le corollaire \emph{sur chaque intervalle de monotonie} et compte les solutions.
\end{itemize}
\end{attention}

\courssub{Localisation d'une solution : balayage et dichotomie}

Une fois l'existence et l'unicité de $\alpha$ établies, il reste à la \emph{localiser}. Deux techniques complémentaires :

\begin{itemize}
\item \textbf{Balayage :} on calcule $f$ en des points régulièrement espacés (par exemple de $0{,}1$ en $0{,}1$) et l'on repère un \textbf{changement de signe} de $f(x) - c$ : si $f$ est continue et $f(x_{1}) - c$ et $f(x_{2}) - c$ sont de signes contraires, alors $\alpha \in [x_{1}\,;x_{2}]$.
\item \textbf{Dichotomie :} on part d'un encadrement $a_{0} < \alpha < b_{0}$ avec $f(a_{0}) - c$ et $f(b_{0}) - c$ de signes contraires ; on teste le milieu et l'on garde la moitié d'intervalle où le signe change. Après $n$ étapes, l'amplitude est $\dfrac{b_{0}-a_{0}}{2^{n}}$. Pour obtenir un encadrement à $10^{-2}$ près à partir d'un intervalle d'amplitude $1$, il faut $n$ tel que $2^{n} \geq 100$, soit $n = 7$ étapes.
\end{itemize}

\begin{exoresolu}[Le problème de l'opérateur de Douala]
\textbf{Énoncé.} Soit $f(x) = x^{3} - 3x + 5$ définie sur $\mathbb{R}$. Montrer que l'équation $f(x) = 4$ admet une unique solution $\alpha$ dans $\mathbb{R}$, puis donner un encadrement de $\alpha$ d'amplitude $0{,}1$.
\tcblower
\textbf{Solution.} \textbf{1. Continuité :} $f$ est une fonction polynôme, donc continue sur $\mathbb{R}$.

\textbf{2. Variations :} $f'(x) = 3x^{2} - 3 = 3(x-1)(x+1)$. On a $f'(x) > 0$ sur $]{-\infty}\,;-1[$ et $]1\,;+\infty[$, et $f'(x) < 0$ sur $]-1\,;1[$. Valeurs remarquables : $f(-1) = 7$, $f(1) = 3$. Limites : $\lim_{-\infty} f = -\infty$ et $\lim_{+\infty} f = +\infty$.

\textbf{3. Application du corollaire sur chaque intervalle de monotonie :}
\begin{itemize}
\item Sur $]{-\infty}\,;-1]$ : $f$ est continue, strictement croissante, et $f\big(]{-\infty}\,;-1]\big) = ]{-\infty}\,;7]$. Or $4 \in ]{-\infty}\,;7]$, donc l'équation $f(x) = 4$ y admet une unique solution $\alpha$.
\item Sur $[-1\,;1]$ : $f$ y est strictement décroissante, $f([-1\,;1]) = [3\,;7]$, et $4 \in [3\,;7]$ : il y a une unique solution $\beta$ sur cet intervalle.
\item Sur $[1\,;+\infty[$ : $f\big([1\,;+\infty[\big) = [3\,;+\infty[$ et $4 \in [3\,;+\infty[$ : il y a une unique solution $\gamma$.
\end{itemize}
L'équation admet donc en réalité \textbf{trois} solutions réelles ! C'est une excellente illustration du piège signalé plus haut : $f$ n'est pas monotone sur $\mathbb{R}$ tout entier. Intéressons-nous à la solution $\alpha \in ]{-\infty}\,;-1]$, qui correspond au problème concret si l'on cherche la plus petite date.

\textbf{4. Encadrement par balayage.} On calcule :
\[ f(-3) = -27 + 9 + 5 = -13 < 4, \qquad f(-2) = -8 + 6 + 5 = 3 < 4, \qquad f(-1) = 7 > 4. \]
Donc $\alpha \in ]{-2}\,;{-1}[$. On affine au dixième : $f(-2) = 3 < 4$ et
\[ f(-1{,}9) = (-1{,}9)^{3} - 3(-1{,}9) + 5 = -6{,}859 + 5{,}7 + 5 = 3{,}841 < 4. \]
\[ f(-1{,}8) = -5{,}832 + 5{,}4 + 5 = 4{,}568 > 4. \]
Comme $f$ est strictement croissante sur $]{-\infty}\,;-1]$, on conclut :
\[ -1{,}9 < \alpha < -1{,}8. \]
Encadrement d'amplitude $0{,}1$ obtenu. Une dichotomie permettrait de poursuivre jusqu'à $10^{-2}$ près : on trouverait $-1{,}88 < \alpha < -1{,}87$ (en effet $f(-1{,}88) \approx 3{,}995 < 4$ et $f(-1{,}87) \approx 4{,}071 > 4$).
\end{exoresolu}

\begin{exoresolu}[Existence et unicité de la racine cubique]
\textbf{Énoncé.} Montrer que l'équation $x^{3} = 10$ admet une unique solution réelle positive, et l'encadrer entre deux entiers consécutifs. (C'est le problème de l'ingénieur de la SONAREL.)
\tcblower
\textbf{Solution.} Soit $g(x) = x^{3}$ sur $[0\,;+\infty[$. $g$ est une fonction polynôme, donc \textbf{continue}, et $g'(x) = 3x^{2} \geq 0$ ne s'annule qu'en $0$ : $g$ est \textbf{strictement croissante} sur $[0\,;+\infty[$. De plus $g(0) = 0$ et $\lim_{+\infty} g = +\infty$, donc $g\big([0\,;+\infty[\big) = [0\,;+\infty[$. Or $10 \in [0\,;+\infty[$ : d'après le corollaire du TVI, l'équation $x^{3} = 10$ admet une \textbf{unique} solution dans $[0\,;+\infty[$, notée $\sqrt[3]{10}$.

Balayage : $g(2) = 8 < 10$ et $g(3) = 27 > 10$, donc
\[ 2 < \sqrt[3]{10} < 3. \]
Plus finement : $g(2{,}1) = 9{,}261 < 10$ et $g(2{,}2) = 10{,}648 > 10$, d'où $2{,}1 < \sqrt[3]{10} < 2{,}2$. L'ingénieur peut donc dimensionner sa cuve avec $\sqrt[3]{10} \approx 2{,}15$, valeur garantie à $5 \times 10^{-2}$ près.
\end{exoresolu}

\begin{savaistu}[Le TVI, un théorème aux applications surprenantes]
Le théorème des valeurs intermédiaires a des conséquences étonnantes. Par exemple, il permet de prouver qu'à tout instant, il existe sur Terre deux points diamétralement opposés ayant exactement la même température ! Il est aussi à la base des algorithmes de recherche de zéros programmés dans les calculatrices et les logiciels utilisés par les ingénieurs de CAMTEL ou de la SONAREL pour résoudre numériquement des équations sans solution exacte calculable.
\end{savaistu}

\begin{aretenir}[L'essentiel de la section]
\begin{itemize}
\item L'image d'un \textbf{intervalle} par une fonction \textbf{continue} est un intervalle ; l'image d'un \textbf{segment} $[a\,;b]$ est un segment $[m\,;M]$.
\item \textbf{TVI :} $f$ continue sur $[a\,;b]$ et $c$ entre $f(a)$ et $f(b)$ $\Rightarrow$ $f(x) = c$ admet \emph{au moins une} solution dans $[a\,;b]$.
\item \textbf{Corollaire (bijection) :} $f$ continue \emph{et} strictement monotone sur $I$, $c \in f(I)$ $\Rightarrow$ $f(x) = c$ admet une \emph{unique} solution dans $I$.
\item La solution se localise par \textbf{balayage} ou \textbf{dichotomie} (amplitude divisée par $2$ à chaque étape).
\end{itemize}
\end{aretenir}

\begin{exercice}[Pour s'entraîner]
Soit $f$ la fonction définie sur $\mathbb{R}$ par $f(x) = x^{3} + x - 1$.
\begin{enumerate}
\item Étudier les variations de $f$ sur $\mathbb{R}$ et dresser son tableau de variations complet.
\item Montrer que l'équation $f(x) = 0$ admet une unique solution $\alpha$ dans $\mathbb{R}$.
\item Justifier que $\alpha \in {]0\,;1[}$, puis donner un encadrement de $\alpha$ d'amplitude $10^{-2}$ par dichotomie.
\end{enumerate}
\end{exercice}

\begin{corrige}[Exercice 1]
\textbf{1.} $f$ est un polynôme, continue et dérivable sur $\mathbb{R}$. $f'(x) = 3x^{2} + 1 > 0$ pour tout réel $x$ : $f$ est \textbf{strictement croissante} sur $\mathbb{R}$. Limites : $\lim_{-\infty} f = -\infty$ et $\lim_{+\infty} f = +\infty$.

\textbf{2.} $f$ est continue et strictement croissante sur $\mathbb{R}$, et $f(\mathbb{R}) = {]{-\infty}\,;+\infty[} = \mathbb{R}$. Or $0 \in \mathbb{R}$ : d'après le corollaire du TVI, l'équation $f(x) = 0$ admet une \textbf{unique} solution $\alpha$ dans $\mathbb{R}$.

\textbf{3.} $f(0) = -1 < 0$ et $f(1) = 1 > 0$ : comme $f$ est strictement croissante, $\alpha \in {]0\,;1[}$. Dichotomie : $f(0{,}5) = -0{,}375 < 0$ donc $\alpha \in {]0{,}5\,;1[}$ ; $f(0{,}75) \approx 0{,}172 > 0$ donc $\alpha \in {]0{,}5\,;0{,}75[}$ ; $f(0{,}625) \approx -0{,}131 < 0$ donc $\alpha \in {]0{,}625\,;0{,}75[}$ ; $f(0{,}6875) \approx 0{,}0125 > 0$ donc $\alpha \in {]0{,}625\,;0{,}6875[}$ ; $f(0{,}65625) \approx -0{,}0613 < 0$ donc $\alpha \in {]0{,}65625\,;0{,}6875[}$ ; $f(0{,}671875) \approx -0{,}0248 < 0$ donc $\alpha \in {]0{,}671875\,;0{,}6875[}$ ; enfin $f(0{,}68) \approx -0{,}0056 < 0$ et $f(0{,}69) \approx 0{,}0185 > 0$. Conclusion :
\[ 0{,}68 < \alpha < 0{,}69, \]
encadrement d'amplitude $10^{-2}$.
\end{corrige}

\coursec{Fonction réciproque d'une bijection}

Dans la section précédente, nous avons établi un résultat fondamental : si $f$ est \cle{continue} et \cle{strictement monotone} sur un intervalle $I$, alors pour tout nombre $c$ compris entre les bornes de $f(I)$, l'équation $f(x) = c$ admet une \textbf{unique} solution dans $I$. Autrement dit, à chaque valeur $y$ de $J = f(I)$ correspond un \emph{unique} antécédent $x$ dans $I$. Ce constat ouvre une porte immense : on peut définir une \emph{nouvelle fonction} qui, partant de $y$, ramène $x$. C'est cette fonction « retour » que nous allons étudier : la \cle{fonction réciproque}.

\courssub{Injection, surjection, bijection : les rappels indispensables}

Avant de « remonter » une fonction, il faut s'assurer que chaque point d'arrivée possède un antécédent, et un seul. Formalisons cela.

\begin{definition}[Injection, surjection, bijection]
Soit $f$ une fonction définie sur un ensemble $E$ à valeurs dans un ensemble $F$.
\begin{itemize}
\item $f$ est \cle{injective} si deux éléments distincts de $E$ ont toujours des images distinctes :
\[ \forall x_1, x_2 \in E, \quad f(x_1) = f(x_2) \implies x_1 = x_2. \]
\item $f$ est \cle{surjective} (de $E$ sur $F$) si tout élément de $F$ possède au moins un antécédent dans $E$ :
\[ \forall y \in F, \; \exists x \in E, \; f(x) = y. \]
\item $f$ est \cle{bijective} (on dit que $f$ est une \cle{bijection}) si elle est à la fois injective et surjective : tout élément de $F$ possède alors \textbf{exactement un} antécédent dans $E$.
\end{itemize}
\end{definition}

\begin{exemplebox}[Tester l'injectivité en pratique]
La fonction $f : x \mapsto x^2$ définie sur $\mathbb{R}$ tout entier n'est \textbf{pas} injective : $f(-2) = f(2) = 4$, donc $4$ a deux antécédents. En revanche, restreinte à $[0 ; +\infty[$, elle devient injective : si $x_1^2 = x_2^2$ avec $x_1, x_2 \geq 0$, alors $x_1 = x_2$. C'est pourquoi la racine carrée n'existe que parce qu'on a \emph{restreint} le domaine de départ.
\end{exemplebox}

\begin{attention}[Une bijection dépend des ensembles de départ ET d'arrivée]
Dire « $f$ est bijective » n'a de sens que si l'on précise \emph{de quel ensemble vers quel ensemble}. La fonction $x \mapsto x^2$ n'est pas une bijection de $\mathbb{R}$ dans $\mathbb{R}$, mais c'est une bijection de $[0 ; +\infty[$ sur $[0 ; +\infty[$. Au Bac, la première étape consiste souvent à \emph{restreindre} l'ensemble de départ pour obtenir une bijection.
\end{attention}

\courssub{Le théorème de la bijection}

Voici le pont entre la section précédente et celle-ci : la continuité et la stricte monotonie garantissent l'existence d'une bijection.

\begin{propriete}[Théorème de la bijection]
Soit $f$ une fonction \cle{continue} et \cle{strictement monotone} sur un intervalle $I$. Alors $f$ réalise une bijection de $I$ sur l'intervalle image $J = f(I)$.
\end{propriete}

\begin{experience}[Démonstration guidée du théorème de la bijection]
Traitons le cas où $f$ est strictement croissante sur $I$.
\begin{enumerate}
\item \textbf{Injectivité.} Soient $x_1 \neq x_2$ dans $I$. Supposons $x_1 < x_2$. Comme $f$ est \emph{strictement} croissante, $f(x_1) < f(x_2)$, donc $f(x_1) \neq f(x_2)$. Deux éléments distincts ont des images distinctes : $f$ est injective.
\item \textbf{Surjectivité sur $J = f(I)$.} Par définition même de l'ensemble image, tout $y \in J$ s'écrit $y = f(x)$ pour un certain $x \in I$ : $f$ est surjective de $I$ sur $J$. (Le théorème des valeurs intermédiaires garantit d'ailleurs que $J$ est bien un intervalle et que toute valeur intermédiaire est atteinte.)
\item \textbf{Conclusion.} $f$ est injective et surjective de $I$ sur $J$ : c'est une bijection de $I$ sur $J$.
\end{enumerate}
Le cas strictement décroissant se traite de manière analogue.
\end{experience}

\courssub{Définition de la fonction réciproque}

Puisque chaque $y$ de $J$ possède un unique antécédent dans $I$, on peut construire la fonction « retour ».

\begin{definition}[Fonction réciproque]
Soit $f$ une bijection d'un intervalle $I$ sur un intervalle $J$. La \cle{fonction réciproque} de $f$, notée $f^{-1}$, est l'unique fonction de $J$ dans $I$ définie par :
\[ \forall x \in I, \; \forall y \in J, \quad f^{-1}(y) = x \iff y = f(x). \]
Elle vérifie les deux relations fondamentales :
\[ f^{-1} \circ f = \mathrm{id}_I \quad \text{et} \quad f \circ f^{-1} = \mathrm{id}_J, \]
c'est-à-dire : pour tout $x \in I$, $f^{-1}\bigl(f(x)\bigr) = x$, et pour tout $y \in J$, $f\bigl(f^{-1}(y)\bigr) = y$.
\end{definition}

\begin{popfigure}
\begin{center}
\begin{tikzpicture}[scale=1.0, >=Stealth]
% Ensemble I
\draw[thick, popblue] (0,0) ellipse (1.3 and 1.9);
\node[popblue, font=\bfseries] at (0,2.3) {$I$};
\fill[popink] (-0.3,-0.4) circle (2pt) node[left=2pt, popdark] {$x$};
% Ensemble J
\draw[thick, poporange] (6,0) ellipse (1.3 and 1.9);
\node[poporange, font=\bfseries] at (6,2.3) {$J$};
\fill[popink] (6.3,-0.4) circle (2pt) node[right=2pt, popdark] {$y = f(x)$};
% Flèche f
\draw[->, very thick, popblue] (0.9,0.6) .. controls (2.2,1.4) and (3.8,1.4) .. (5.1,0.6)
node[midway, above, popblue] {$f$};
% Flèche f^{-1}
\draw[->, very thick, poporange] (5.1,-1.2) .. controls (3.8,-2.0) and (2.2,-2.0) .. (0.9,-1.2)
node[midway, below, poporange] {$f^{-1}$};
% Points reliés
\draw[->, thick, poppurple, dashed] (-0.3,-0.4) -- (6.3,-0.4);
\node[popmuted, font=\small] at (3,-3.0) {$f^{-1}\bigl(f(x)\bigr) = x$ \quad et \quad $f\bigl(f^{-1}(y)\bigr) = y$};
\end{tikzpicture}
\end{center}
\legende{La bijection $f$ envoie $I$ sur $J$ ; sa réciproque $f^{-1}$ ramène chaque $y$ sur son unique antécédent.}
\end{popfigure}

\begin{attention}[$f^{-1}(x)$ n'est PAS $\dfrac{1}{f(x)}$]
C'est la confusion la plus fréquente (et la plus sanctionnée) ! La notation $f^{-1}$ désigne la \emph{réciproque} (la fonction « retour »), pas l'\emph{inverse}. Par exemple, si $f(x) = x^2$ sur $[0;+\infty[$, alors $f^{-1}(x) = \sqrt{x}$, tandis que $\dfrac{1}{f(x)} = \dfrac{1}{x^2}$ : deux fonctions totalement différentes. L'inverse de $f$ se note $\dfrac{1}{f}$, jamais $f^{-1}$.
\end{attention}

\courssub{Continuité et monotonie de la réciproque}

\begin{propriete}[Continuité et monotonie de la fonction réciproque]
Soit $f$ une fonction continue et strictement monotone sur un intervalle $I$, réalisant une bijection de $I$ sur $J = f(I)$. Alors :
\begin{itemize}
\item sa réciproque $f^{-1}$ est \textbf{continue} sur $J$ (résultat admis) ;
\item $f^{-1}$ est \textbf{strictement monotone sur $J$, de même sens de variation que $f$}.
\end{itemize}
\end{propriete}

L'existence de $f^{-1}$ découle du théorème de la bijection démontré plus haut. La continuité de $f^{-1}$ est admise (sa démonstration dépasse le programme). Le fait que $f^{-1}$ ait le \emph{même sens} de variation que $f$ se comprend bien : si $f$ « monte », remonter la courbe en sens inverse la fait monter aussi. Rigoureusement, si $f$ est strictement croissante et si $y_1 < y_2$ dans $J$, posons $x_1 = f^{-1}(y_1)$ et $x_2 = f^{-1}(y_2)$. Si l'on avait $x_1 \geq x_2$, la croissance de $f$ donnerait $f(x_1) \geq f(x_2)$, c'est-à-dire $y_1 \geq y_2$ : contradiction. Donc $x_1 < x_2$ et $f^{-1}$ est strictement croissante.

\courssub{Dérivabilité de la fonction réciproque}

\begin{propriete}[Dérivée de la fonction réciproque]
Soit $f$ une bijection de $I$ sur $J$, dérivable en un point $a \in I$ avec $f'(a) \neq 0$. Alors $f^{-1}$ est dérivable en $b = f(a)$ et :
\[ \left(f^{-1}\right)'(b) = \frac{1}{f'\bigl(f^{-1}(b)\bigr)} = \frac{1}{f'(a)}. \]
Si $f'(a) = 0$, alors $f^{-1}$ n'est pas dérivable en $b = f(a)$ : sa courbe admet en ce point une tangente \emph{verticale}.
\end{propriete}

\begin{experience}[Démonstration guidée à partir du taux d'accroissement]
Posons $b = f(a)$ et soit $y \in J$, $y \neq b$. Posons $x = f^{-1}(y)$, de sorte que $y = f(x)$. Comme $f$ est bijective, $y \neq b$ équivaut à $x \neq a$. Le taux d'accroissement de $f^{-1}$ entre $b$ et $y$ s'écrit :
\[ \frac{f^{-1}(y) - f^{-1}(b)}{y - b} = \frac{x - a}{f(x) - f(a)} = \frac{1}{\dfrac{f(x) - f(a)}{x - a}}. \]
Quand $y \to b$, la continuité de $f^{-1}$ en $b$ donne $x = f^{-1}(y) \to f^{-1}(b) = a$. Le dénominateur $\dfrac{f(x)-f(a)}{x-a}$ tend alors vers $f'(a)$, qui est non nul par hypothèse. Par passage à la limite du quotient :
\[ \lim_{y \to b} \frac{f^{-1}(y) - f^{-1}(b)}{y - b} = \frac{1}{f'(a)}. \]
Donc $f^{-1}$ est dérivable en $b$ et $(f^{-1})'(b) = \dfrac{1}{f'(a)} = \dfrac{1}{f'\bigl(f^{-1}(b)\bigr)}$.
\end{experience}

\begin{attention}[N'oublie pas de composer par $f^{-1}$]
Erreur classique : écrire $(f^{-1})'(x) = \dfrac{1}{f'(x)}$. C'est faux ! La formule correcte est
\[ (f^{-1})'(x) = \frac{1}{f'\bigl(f^{-1}(x)\bigr)}. \]
La dérivée de $f$ doit être évaluée \emph{au point $f^{-1}(x)$}, pas au point $x$. Vérifie avec $f(x) = x^2$ sur $[0;+\infty[$ : $f'(x) = 2x$, donc $(f^{-1})'(x) = \dfrac{1}{2\sqrt{x}}$, qui est bien la dérivée de $\sqrt{x}$ — alors que $\dfrac{1}{f'(x)} = \dfrac{1}{2x}$ serait faux.
\end{attention}

\begin{aretenir}[Les trois propriétés clés de la réciproque]
Si $f$ est continue et strictement monotone sur un intervalle $I$ (donc bijective de $I$ sur $J=f(I)$) :
\begin{enumerate}
\item \textbf{Continuité :} $f^{-1}$ est continue sur $J$.
\item \textbf{Monotonie :} $f^{-1}$ a le \emph{même sens} de variation que $f$ sur $J$.
\item \textbf{Dérivabilité :} Si $f$ est dérivable en $a$ avec $f'(a)\neq 0$, alors $f^{-1}$ est dérivable en $b=f(a)$ et $(f^{-1})'(b) = \dfrac{1}{f'(a)}$.
\item \textbf{Géométrie :} Les courbes de $f$ et $f^{-1}$ sont symétriques par rapport à la droite $y=x$.
\end{enumerate}
\end{aretenir}

\courssub{Symétrie des courbes par rapport à la droite $y = x$}

Voici la propriété géométrique la plus spectaculaire — et la plus utile pour les tracés.

\begin{propriete}[Symétrie des courbes de $f$ et $f^{-1}$]
Dans un repère orthonormé, les courbes représentatives de $f$ et de $f^{-1}$ sont \cle{symétriques} l'une de l'autre par rapport à la droite d'équation $y = x$ (la première bissectrice).
\end{propriete}

En effet, le point $M(x\,;\,y)$ appartient à la courbe de $f$ si et seulement si $y = f(x)$, ce qui équivaut à $x = f^{-1}(y)$, c'est-à-dire que le point $M'(y\,;\,x)$ appartient à la courbe de $f^{-1}$. Or, dans un repère orthonormé, échanger les coordonnées revient exactement à effectuer la symétrie orthogonale par rapport à la droite $y = x$.

\begin{methode}[Tracer la courbe de $f^{-1}$ à partir de celle de $f$]
\begin{enumerate}
\item Place quelques points remarquables de la courbe de $f$ : intersections avec les axes, extremums, points dont les coordonnées sont entières.
\item Pour chaque point $M(a\,;\,b)$, construis son symétrique $M'(b\,;\,a)$ par rapport à la droite $y = x$ : trace la droite $y = x$, puis reporte $M'$ en échangeant abscisse et ordonnée (la perpendiculaire à $y = x$ passant par $M$ coupe $y = x$ en son milieu).
\item Relie les points symétriques par une courbe lisse, en respectant le sens de variation (le même que celui de $f$).
\item N'oublie pas : une tangente de pente $m$ en $M$ devient une tangente de pente $\dfrac{1}{m}$ en $M'$ ; une tangente horizontale devient verticale.
\end{enumerate}
\end{methode}

\courssub{Exemples de référence}

\textbf{Exemple 1 : la fonction carré sur $[0 ; +\infty[$ et la racine carrée.}

La fonction $f : x \mapsto x^2$ est continue et strictement croissante sur $I = [0 ; +\infty[$ (sa dérivée $2x$ est strictement positive sur $]0 ; +\infty[$). D'après le théorème de la bijection, elle réalise une bijection de $[0 ; +\infty[$ sur $J = f(I) = [0 ; +\infty[$. Sa réciproque est la fonction \cle{racine carrée} :
\[ f^{-1} : [0 ; +\infty[ \to [0 ; +\infty[, \quad x \mapsto \sqrt{x}. \]
Calculons sa dérivée avec la formule : pour $x > 0$, $f'(t) = 2t$, donc
\[ (f^{-1})'(x) = \frac{1}{f'\bigl(f^{-1}(x)\bigr)} = \frac{1}{2\sqrt{x}}. \]
On retrouve la dérivée classique de $\sqrt{x}$. Remarque : $f'(0) = 0$, donc $f^{-1}$ n'est pas dérivable en $0$ — la courbe de $\sqrt{x}$ admet une tangente verticale à l'origine.

\begin{popfigure}
\begin{center}
\begin{tikzpicture}
\begin{axis}[
  width=0.85\linewidth, height=7.5cm,
  axis lines=middle, xlabel={$x$}, ylabel={$y$},
  xmin=-0.6, xmax=4.4, ymin=-0.6, ymax=4.4,
  xtick={1,2,3,4}, ytick={1,2,3,4},
  grid=both, grid style={popline!40},
  legend style={at={(0.03,0.97)}, anchor=north west, font=\small}
]
\addplot[popcurve, popblue, domain=0:2.05, samples=120] {x^2};
\addlegendentry{$y = x^2$ sur $[0;+\infty[$}
\addplot[popcurve, poporange, domain=0:4.3, samples=150] {sqrt(max(0,x))};
\addlegendentry{$y = \sqrt{x}$}
\addplot[dashed, popmuted, domain=-0.4:4.3] {x};
\addlegendentry{$y = x$}
% Couple de points symétriques
\addplot[only marks, mark=*, popink] coordinates {(2,4) (4,2)};
\draw[thick, poppurple, dashed] (axis cs:2,4) -- (axis cs:4,2);
\node[popdark, font=\small, anchor=south east] at (axis cs:2,4) {$M(2\,;\,4)$};
\node[popdark, font=\small, anchor=north west] at (axis cs:4,2) {$M'(4\,;\,2)$};
\end{axis}
\end{tikzpicture}
\end{center}
\legende{Les courbes de $x \mapsto x^2$ (sur $[0;+\infty[$) et de $x \mapsto \sqrt{x}$ sont symétriques par rapport à la droite $y = x$ : les points $M(2\,;\,4)$ et $M'(4\,;\,2)$ se correspondent.}
\end{popfigure}

\textbf{Exemple 2 : la fonction cube et la racine cubique.}

La fonction $g : x \mapsto x^3$ est continue et strictement croissante sur $\mathbb{R}$ tout entier (sa dérivée $3x^2$ est positive et ne s'annule qu'en $0$, sans changement de signe). Elle réalise donc une bijection de $\mathbb{R}$ sur $\mathbb{R}$, et sa réciproque est la fonction \cle{racine cubique} :
\[ g^{-1} : \mathbb{R} \to \mathbb{R}, \quad x \mapsto \sqrt[3]{x}. \]
Contrairement à la racine carrée, la racine cubique est définie sur $\mathbb{R}$ entier : $\sqrt[3]{-8} = -2$ car $(-2)^3 = -8$. Sa dérivée pour $x \neq 0$ :
\[ (g^{-1})'(x) = \frac{1}{3\left(\sqrt[3]{x}\right)^2}. \]

\begin{popfigure}
\begin{center}
\begin{tikzpicture}
\begin{axis}[
  width=0.85\linewidth, height=7.5cm,
  axis lines=middle, xlabel={$x$}, ylabel={$y$},
  xmin=-2.4, xmax=2.4, ymin=-2.4, ymax=2.4,
  xtick={-2,-1,1,2}, ytick={-2,-1,1,2},
  grid=both, grid style={popline!40},
  legend style={at={(0.03,0.97)}, anchor=north west, font=\small}
]
\addplot[popcurve, popblue, domain=-1.32:1.32, samples=150] {x^3};
\addlegendentry{$y = x^3$}
\addplot[popcurve, poporange, domain=0:2.3, samples=120] {x^(1/3)};
\addplot[popcurve, poporange, domain=-2.3:0, samples=120] {-(-x)^(1/3)};
\addlegendentry{$y = \sqrt[3]{x}$}
\addplot[dashed, popmuted, domain=-2.3:2.3] {x};
\addlegendentry{$y = x$}
\addplot[only marks, mark=*, popink] coordinates {(1.26,2) (2,1.26)};
\draw[thick, poppurple, dashed] (axis cs:1.26,2) -- (axis cs:2,1.26);
\end{axis}
\end{tikzpicture}
\end{center}
\legende{La fonction cube est bijective sur $\mathbb{R}$ ; sa réciproque, la racine cubique, lui est symétrique par rapport à $y = x$.}
\end{popfigure}

\courssub{Déterminer l'expression explicite de $f^{-1}$}

\begin{methode}[Calculer l'expression de $f^{-1}$]
Pour déterminer explicitement la réciproque d'une bijection $f : I \to J$ :
\begin{enumerate}
\item Vérifie (ou justifie) que $f$ est continue et strictement monotone sur $I$, et détermine $J = f(I)$.
\item Pose $y = f(x)$ avec $x \in I$ et $y \in J$.
\item Résous cette équation \textbf{d'inconnue $x$} : tu obtiens $x$ en fonction de $y$.
\item Conclus : $f^{-1}(y) = $ l'expression trouvée, puis (par convention) renomme la variable : $f^{-1}(x) = \ldots$, en précisant son ensemble de définition $J$.
\end{enumerate}
\end{methode}

\begin{exoresolu}[Réciproque d'une fonction homographique]
Soit $f$ la fonction définie sur $I = ]1 ; +\infty[$ par $f(x) = \dfrac{2x + 1}{x - 1}$.
\begin{enumerate}
\item Montrer que $f$ réalise une bijection de $I$ sur un intervalle $J$ à préciser.
\item Déterminer l'expression de $f^{-1}$.
\end{enumerate}
\tcblower
\textbf{1.} $f$ est une fonction rationnelle définie, continue et dérivable sur $I = ]1 ; +\infty[$. Pour tout $x \in I$ :
\[ f'(x) = \frac{2(x-1) - (2x+1)}{(x-1)^2} = \frac{-3}{(x-1)^2} < 0. \]
Donc $f$ est strictement décroissante sur $I$. De plus, $\displaystyle\lim_{x \to 1^+} f(x) = +\infty$ (numérateur tend vers $3 > 0$, dénominateur vers $0^+$) et $\displaystyle\lim_{x \to +\infty} f(x) = 2$. Par le théorème de la bijection, $f$ réalise une bijection de $I$ sur $J = ]2 ; +\infty[$.

\textbf{2.} Soit $y \in J$. Résolvons $y = \dfrac{2x+1}{x-1}$ d'inconnue $x$ :
\begin{align*}
y(x-1) &= 2x + 1 \\
yx - y &= 2x + 1 \\
x(y - 2) &= y + 1 \\
x &= \frac{y+1}{y-2} \quad (\text{licite car } y > 2).
\end{align*}
Conclusion : $\displaystyle f^{-1}(x) = \frac{x+1}{x-2}$, définie sur $J = ]2 ; +\infty[$. On vérifie au passage que $f^{-1}$ est bien décroissante sur $J$, comme $f$.
\end{exoresolu}

\begin{exoresolu}[Calculer $(f^{-1})'(b)$ sans connaître $f^{-1}$]
Soit $f$ la fonction définie sur $\mathbb{R}$ par $f(x) = x^3 + 2x - 1$. On admet que $f$ est une bijection de $\mathbb{R}$ sur $\mathbb{R}$. Calculer $\left(f^{-1}\right)'(2)$.
\tcblower
L'expression explicite de $f^{-1}$ est introuvable (il faudrait résoudre une équation du troisième degré), mais la formule de la dérivée de la réciproque s'applique. On cherche d'abord l'antécédent de $2$ : on remarque que $f(1) = 1 + 2 - 1 = 2$, donc $f^{-1}(2) = 1$.
Ensuite, $f'(x) = 3x^2 + 2$, donc $f'(1) = 5 \neq 0$. La formule donne :
\[ \left(f^{-1}\right)'(2) = \frac{1}{f'\bigl(f^{-1}(2)\bigr)} = \frac{1}{f'(1)} = \frac{1}{5}. \]
\end{exoresolu}

\begin{savaistu}[La notation $f^{-1}$ et le piège classique du Bac]
La notation $f^{-1}$ pour la réciproque vient de l'algèbre des applications : composer par $f^{-1}$ « annule » l'effet de $f$, comme multiplier par $a^{-1} = \frac{1}{a}$ annule la multiplication par $a$. Mais pour les fonctions, $f^{-1}$ ne signifie \emph{jamais} $\frac{1}{f}$ ! Au Baccalauréat, une question extrêmement classique consiste à calculer $(f^{-1})'(b)$ \textbf{sans} connaître l'expression de $f^{-1}$ : il suffit de trouver l'antécédent $a$ tel que $f(a) = b$ (souvent une valeur évidente à tester : $0$, $1$, $2$\ldots), puis d'appliquer $(f^{-1})'(b) = \dfrac{1}{f'(a)}$. Les fonctions trigonométriques réciproques (Arc sinus, Arc cosinus, Arc tangente), que tu rencontreras peut-être en enseignement supérieur, sont construites exactement selon ce schéma, par restriction à un intervalle de monotonie.
\end{savaistu}

\begin{exercice}[Pour t'entraîner]
\begin{enumerate}
\item Soit $f(x) = \sqrt{x} + 1$ définie sur $[0 ; +\infty[$. Montrer que $f$ est une bijection de $[0;+\infty[$ sur un intervalle $J$ à préciser, puis déterminer $f^{-1}$.
\item Soit $g(x) = x^3 + x$ sur $\mathbb{R}$. On admet que $g$ est bijective. Calculer $(g^{-1})'(2)$ et $(g^{-1})'(0)$.
\item Soit $h(x) = \dfrac{x}{x+2}$ sur $]-2 ; +\infty[$. Déterminer $h^{-1}$ et tracer sur un même graphique les courbes de $h$ et $h^{-1}$.
\end{enumerate}
\end{exercice}

\begin{corrige}[Exercice 1]
\textbf{1.} $f$ est continue sur $[0;+\infty[$ et $f'(x) = \dfrac{1}{2\sqrt{x}} > 0$ sur $]0;+\infty[$ : $f$ est strictement croissante. Comme $f(0) = 1$ et $\lim_{+\infty} f = +\infty$, $J = [1 ; +\infty[$. On résout $y = \sqrt{x} + 1$ : $\sqrt{x} = y - 1$, d'où $x = (y-1)^2$ (licite car $y \geq 1$). Donc $f^{-1}(x) = (x-1)^2$ sur $[1 ; +\infty[$.

\textbf{2.} $g(1) = 2$ donc $g^{-1}(2) = 1$ ; $g'(x) = 3x^2 + 1$, $g'(1) = 4$, donc $(g^{-1})'(2) = \dfrac{1}{4}$. De même $g(0) = 0$ donc $(g^{-1})'(0) = \dfrac{1}{g'(0)} = 1$.

\textbf{3.} $h'(x) = \dfrac{2}{(x+2)^2} > 0$ : $h$ est strictement croissante, bijective de $]-2;+\infty[$ sur $]-\infty ; 1[$ (limites : $-\infty$ en $-2^+$, $1$ en $+\infty$). On résout $y = \dfrac{x}{x+2}$ : $y(x+2) = x$, d'où $x(y-1) = -2y$ et $x = \dfrac{2y}{1-y}$. Donc $h^{-1}(x) = \dfrac{2x}{1-x}$ sur $]-\infty ; 1[$. Pour le tracé, on place la droite $y = x$ et on construit les symétriques de points remarquables : $(0;0)$ est commun aux deux courbes, $(2\,;\,\tfrac12)$ sur la courbe de $h$ correspond à $(\tfrac12\,;\,2)$ sur celle de $h^{-1}$.
\end{corrige}

\coursec{Racines n-ièmes et puissances rationnelles}

Dans la section précédente, nous avons vu que toute bijection $f : I \to J$ admet une fonction réciproque $f^{-1} : J \to I$, continue, de même sens de variation, et dont la courbe est symétrique de celle de $f$ par rapport à la droite d'équation $y = x$. Nous allons maintenant mettre cette belle machinerie en action sur une famille de fonctions que tu connais depuis longtemps : les fonctions puissances $x \mapsto x^n$. Leur réciproque va nous offrir un nouvel outil de calcul fondamental : la \cle{racine n-ième}, généralisation de la racine carrée, qui permet de résoudre les équations $x^n = a$ et de donner un sens précis à des écritures comme $8^{2/3}$.

\courssub{La fonction $x \mapsto x^n$ sur $[0\,;\,+\infty[$}

Soit $n$ un entier naturel supérieur ou égal à $2$. Considérons la fonction $f_n : x \mapsto x^n$, définie sur $\mathbb{R}$ tout entier, mais que nous restreindrons d'abord à l'intervalle $[0\,;\,+\infty[$.

\begin{propriete}[{Bijectivité de $x \mapsto x^n$ sur $[0\,;\,+\infty[$}]
Pour tout entier $n \geq 2$, la fonction $f_n : x \mapsto x^n$ est \textbf{continue} et \textbf{strictement croissante} sur $[0\,;\,+\infty[$. De plus :
\[ f_n(0) = 0 \qquad \text{et} \qquad \lim_{x \to +\infty} x^n = +\infty. \]
Par conséquent, $f_n$ réalise une \textbf{bijection} de $[0\,;\,+\infty[$ sur $[0\,;\,+\infty[$.
\end{propriete}

\begin{experience}[Démonstration guidée]
\begin{enumerate}
\item \textbf{Continuité.} $f_n$ est une fonction polynôme, donc continue sur $\mathbb{R}$, et en particulier sur $[0\,;\,+\infty[$.
\item \textbf{Stricte croissance.} $f_n$ est dérivable sur $[0\,;\,+\infty[$ et $f_n'(x) = n\,x^{n-1}$. Pour tout $x > 0$, on a $x^{n-1} > 0$, donc $f_n'(x) > 0$ : la fonction est strictement croissante sur $[0\,;\,+\infty[$.
\item \textbf{Limites aux bornes.} $f_n(0) = 0^n = 0$, et comme $n \geq 1$, $\lim_{x \to +\infty} x^n = +\infty$.
\item \textbf{Conclusion.} $f_n$ est continue et strictement monotone sur l'intervalle $[0\,;\,+\infty[$ : d'après le théorème de la bijection (conséquence du théorème des valeurs intermédiaires), elle réalise une bijection de $[0\,;\,+\infty[$ sur l'intervalle image $[0\,;\,+\infty[$.
\end{enumerate}
\end{experience}

L'intuition est simple : quand $x$ parcourt les réels positifs en croissant, $x^n$ parcourt, lui aussi en croissant et sans sauter aucune valeur, tous les réels positifs. Chaque réel positif $a$ est donc atteint \emph{une fois et une seule} : c'est exactement ce qui va nous permettre de définir la racine n-ième.

\begin{popfigure}
\begin{center}
\begin{tikzpicture}
\begin{axis}[
  width=0.85\linewidth, height=7cm,
  axis lines=middle,
  xlabel={$x$}, ylabel={$y$},
  xmin=0, xmax=2.2, ymin=0, ymax=4.5,
  xtick={0,0.5,1,1.5,2}, ytick={1,2,3,4},
  legend style={at={(0.02,0.98)}, anchor=north west, font=\small},
  grid=both, grid style={popline!40},
]
\addplot[popcurve, popblue, domain=0:2, samples=100]{x^2};
\addlegendentry{$y = x^2$}
\addplot[popcurve, poporange, domain=0:1.6, samples=100]{x^3};
\addlegendentry{$y = x^3$}
\addplot[popcurve, poppurple, domain=0:1.42, samples=100]{x^4};
\addlegendentry{$y = x^4$}
% point (a ; racine) : exemple a = 2, n = 2
\draw[popink, dashed] (axis cs:0,2) -- (axis cs:1.4142,2);
\draw[popink, dashed] (axis cs:1.4142,0) -- (axis cs:1.4142,2);
\addplot[only marks, mark=*, popink] coordinates {(1.4142,2)};
\node[popink, anchor=north, font=\small] at (axis cs:1.4142,0) {$\sqrt{2}$};
\node[popink, anchor=west, font=\small] at (axis cs:0.05,2) {$a=2$};
\end{axis}
\end{tikzpicture}
\end{center}
\legende{Courbes de $x \mapsto x^n$ pour $n = 2, 3, 4$ sur $[0\,;\,2]$. Chacune est strictement croissante : la valeur $a = 2$ est atteinte une seule fois par $x^2$, en $x = \sqrt{2}$.}
\end{popfigure}

\courssub{Définition de la racine n-ième}

\begin{definition}[Racine n-ième d'un réel positif]
Soit $n$ un entier supérieur ou égal à $2$ et $a$ un réel \textbf{positif ou nul}. On appelle \cle{racine n-ième} de $a$, et on note $\sqrt[n]{a}$, l'unique réel \textbf{positif ou nul} dont la puissance n-ième vaut $a$ :
\[ \sqrt[n]{a} \text{ est l'unique } x \geq 0 \text{ tel que } x^n = a. \]
Autrement dit, $\sqrt[n]{a} = f_n^{-1}(a)$, où $f_n^{-1}$ est la réciproque de la bijection $x \mapsto x^n$ de $[0\,;\,+\infty[$ sur $[0\,;\,+\infty[$.
\end{definition}

\begin{aretenir}[Cas particuliers à connaître]
\begin{itemize}
\item Pour $n = 2$ : la \textbf{racine carrée}, notée $\sqrt{a}$ (on omet l'indice $2$). Exemple : $\sqrt{9} = 3$.
\item Pour $n = 3$ : la \textbf{racine cubique}, notée $\sqrt[3]{a}$. Exemple : $\sqrt[3]{27} = 3$ car $3^3 = 27$.
\item Pour tout $a \geq 0$ et tout $n \geq 2$ : $\left(\sqrt[n]{a}\right)^n = a$ et $\sqrt[n]{a^n} = a$.
\end{itemize}
\end{aretenir}

La fonction $x \mapsto \sqrt[n]{x}$ est donc définie, continue et strictement croissante sur $[0\,;\,+\infty[$ (propriétés héritées de la réciproque), et sa courbe est la symétrique de celle de $x \mapsto x^n$ par rapport à la droite $y = x$.

\begin{popfigure}
\begin{center}
\begin{tikzpicture}
\begin{axis}[
  width=0.85\linewidth, height=7cm,
  axis lines=middle,
  xlabel={$x$}, ylabel={$y$},
  xmin=0, xmax=4.3, ymin=0, ymax=2.3,
  xtick={0,1,2,3,4}, ytick={0.5,1,1.5,2},
  legend style={at={(0.98,0.35)}, anchor=east, font=\small},
  grid=both, grid style={popline!40},
]
\addplot[popcurve, popblue, domain=0:4, samples=150]{sqrt(max(0,x))};
\addlegendentry{$y = \sqrt{x}$}
\addplot[popcurve, poporange, domain=0:4, samples=150]{x^(1/3)};
\addlegendentry{$y = \sqrt[3]{x}$}
\addplot[only marks, mark=*, popink] coordinates {(1,1)};
\node[popink, anchor=south west, font=\small] at (axis cs:1,1) {$(1\,;\,1)$};
\end{axis}
\end{tikzpicture}
\end{center}
\legende{Comparaison de $x \mapsto \sqrt{x}$ et $x \mapsto \sqrt[3]{x}$ sur $[0\,;\,4]$. Les deux courbes passent par $(0\,;\,0)$ et $(1\,;\,1)$ ; pour $x > 1$, $\sqrt{x} > \sqrt[3]{x}$, et pour $0 < x < 1$, c'est l'inverse.}
\end{popfigure}

\courssub{Résolution de l'équation $x^n = a$}

Voici le point où la parité de $n$ joue un rôle décisif. Sur $\mathbb{R}$ tout entier, le comportement de $x \mapsto x^n$ dépend de la parité de $n$ :

\begin{itemize}
\item si $n$ est \textbf{pair}, $x^n \geq 0$ pour tout $x$ réel, et $f_n$ est strictement décroissante sur $]-\infty\,;\,0]$, strictement croissante sur $[0\,;\,+\infty[$ ;
\item si $n$ est \textbf{impair}, $f_n$ est strictement croissante sur $\mathbb{R}$ tout entier, et $x^n$ prend toutes les valeurs réelles, y compris négatives.
\end{itemize}

\begin{propriete}[Nombre de solutions de $x^n = a$]
Soit $n \geq 2$ un entier et $a$ un réel. L'équation $x^n = a$, d'inconnue $x \in \mathbb{R}$, admet :
\begin{itemize}
\item si $a < 0$ : \textbf{aucune solution} lorsque $n$ est pair ; \textbf{une unique solution} $x = -\sqrt[n]{-a}$ lorsque $n$ est impair ;
\item si $a = 0$ : l'unique solution $x = 0$ ;
\item si $a > 0$ : \textbf{deux solutions} $x = \sqrt[n]{a}$ et $x = -\sqrt[n]{a}$ lorsque $n$ est pair ; \textbf{une unique solution} $x = \sqrt[n]{a}$ lorsque $n$ est impair.
\end{itemize}
\end{propriete}

\begin{experience}[Démonstration guidée]
\textbf{Cas $n$ pair.} Pour tout réel $x$, $x^n = (x^2)^{n/2} \geq 0$, donc $x^n = a < 0$ est impossible. Si $a > 0$ : sur $[0\,;\,+\infty[$, $f_n$ est une bijection sur $[0\,;\,+\infty[$, donc l'équation a une unique solution positive $x_0 = \sqrt[n]{a}$. Sur $]-\infty\,;\,0]$, $f_n$ est continue et strictement décroissante, avec $f_n(0) = 0$ et $\lim_{x \to -\infty} x^n = +\infty$ : c'est une bijection de $]-\infty\,;\,0]$ sur $[0\,;\,+\infty[$, donc une unique solution négative existe ; comme $(-x_0)^n = (-1)^n x_0^n = x_0^n = a$ (car $n$ pair), cette solution est $-x_0$.

\textbf{Cas $n$ impair.} $f_n$ est strictement croissante sur $\mathbb{R}$ (sa dérivée $n x^{n-1}$ est positive, ne s'annulant qu'en $0$), continue, avec $\lim_{-\infty} f_n = -\infty$ et $\lim_{+\infty} f_n = +\infty$ : c'est une bijection de $\mathbb{R}$ sur $\mathbb{R}$. L'équation $x^n = a$ a donc une unique solution pour tout $a$ : c'est $\sqrt[n]{a}$ si $a \geq 0$, et $-\sqrt[n]{-a}$ si $a < 0$ (vérification : $(-\sqrt[n]{-a})^n = (-1)^n (\sqrt[n]{-a})^n = -(-a) = a$ car $n$ impair).
\end{experience}

\begin{exemplebox}[Résolutions d'équations $x^n = a$]
\begin{itemize}
\item $x^3 = -8$ : $n = 3$ impair, $a = -8 < 0$, donc une unique solution : $x = -\sqrt[3]{8} = -2$. Vérification : $(-2)^3 = -8$.
\item $x^4 = 16$ : $n = 4$ pair, $a = 16 > 0$, donc deux solutions : $x = \sqrt[4]{16} = 2$ et $x = -2$. Vérification : $2^4 = (-2)^4 = 16$.
\item $x^5 = 32$ : $n = 5$ impair, une unique solution : $x = \sqrt[5]{32} = 2$ car $2^5 = 32$.
\item $x^2 = -4$ : $n$ pair, $a < 0$ : aucune solution réelle.
\end{itemize}
\end{exemplebox}

\begin{attention}[Erreurs fréquentes]
\begin{itemize}
\item \textbf{Oublier la solution négative.} L'équation $x^4 = 16$ a \emph{deux} solutions : $2$ \textbf{et} $-2$. La racine n-ième $\sqrt[n]{a}$ est par définition \emph{positive}, mais l'équation $x^n = a$ avec $n$ pair admet aussi la solution opposée.
\item \textbf{La racine n-ième n'est pas additive :} en général
\[ \sqrt[n]{a + b} \neq \sqrt[n]{a} + \sqrt[n]{b}. \]
Contre-exemple : $\sqrt{9 + 16} = \sqrt{25} = 5$, alors que $\sqrt{9} + \sqrt{16} = 3 + 4 = 7$. De même, $\sqrt[3]{8+8} = \sqrt[3]{16} \approx 2{,}52 \neq 2 + 2 = 4$.
\item \textbf{Écrire $\sqrt[n]{a}$ pour $a < 0$ avec $n$ pair} n'a aucun sens dans $\mathbb{R}$ : $\sqrt{-4}$ n'existe pas (en Terminale, dans $\mathbb{R}$).
\end{itemize}
\end{attention}

\courssub{Puissances à exposant rationnel}

La racine n-ième permet de donner un sens aux exposants fractionnaires. L'idée directrice : on veut que la règle $\left(a^r\right)^s = a^{rs}$ reste vraie. Or, si l'on veut définir $a^{1/n}$, on doit avoir $\left(a^{1/n}\right)^n = a^{1} = a$ : c'est exactement la définition de $\sqrt[n]{a}$.

\begin{definition}[Puissance à exposant rationnel]
Soit $a$ un réel \textbf{strictement positif}, $p$ un entier relatif et $q$ un entier naturel non nul. On pose :
\[ a^{1/q} = \sqrt[q]{a} \qquad \text{et} \qquad a^{p/q} = \left(\sqrt[q]{a}\right)^p = \sqrt[q]{a^p}. \]
\end{definition}

\begin{propriete}[Règles de calcul sur les puissances rationnelles (admises)]
Pour tous réels $a > 0$, $b > 0$ et tous rationnels $r$, $s$ :
\begin{align*}
a^r \times a^s &= a^{r+s} & \dfrac{a^r}{a^s} &= a^{r-s} & \left(a^r\right)^s &= a^{rs} \\
(ab)^r &= a^r b^r & \left(\dfrac{a}{b}\right)^r &= \dfrac{a^r}{b^r} & a^{-r} &= \dfrac{1}{a^r}
\end{align*}
En particulier, pour $a, b > 0$ et $n \geq 2$ :
\[ \sqrt[n]{ab} = \sqrt[n]{a}\,\sqrt[n]{b}, \qquad \sqrt[n]{\dfrac{a}{b}} = \dfrac{\sqrt[n]{a}}{\sqrt[n]{b}}, \qquad \sqrt[n]{a^p} = a^{p/n}, \qquad \left(\sqrt[n]{a}\right)^n = a. \]
\end{propriete}

Ces règles prolongent exactement celles que tu connais pour les exposants entiers : c'est toute la force de la définition. Retiens que l'écriture fractionnaire et l'écriture avec radicaux sont deux visages du même objet, et qu'on passe de l'une à l'autre selon les besoins du calcul.

\begin{exoresolu}[Simplification d'une expression à puissances rationnelles]
Simplifier l'expression $E = \dfrac{8^{2/3} \times 27^{1/3}}{4^{1/2}}$.
\tcblower
\textbf{Solution détaillée.} On exprime chaque base comme puissance d'un nombre premier :
\[ 8 = 2^3, \qquad 27 = 3^3, \qquad 4 = 2^2. \]
On applique ensuite la règle $(a^r)^s = a^{rs}$ :
\begin{align*}
8^{2/3} &= \left(2^3\right)^{2/3} = 2^{3 \times \frac{2}{3}} = 2^2 = 4, \\
27^{1/3} &= \left(3^3\right)^{1/3} = 3^{1} = 3, \\
4^{1/2} &= \left(2^2\right)^{1/2} = 2^{1} = 2.
\end{align*}
Donc :
\[ E = \dfrac{4 \times 3}{2} = \dfrac{12}{2} = 6. \]
\textbf{Vérification par les radicaux :} $8^{2/3} = \left(\sqrt[3]{8}\right)^2 = 2^2 = 4$ ; $27^{1/3} = \sqrt[3]{27} = 3$ ; $4^{1/2} = \sqrt{4} = 2$. Même résultat : $E = 6$.
\end{exoresolu}

\begin{exoresolu}[Équations et simplifications]
\begin{enumerate}
\item Résoudre dans $\mathbb{R}$ l'équation $x^5 = 32$.
\item Simplifier $F = \sqrt[3]{54} \times \sqrt[3]{4} \div 6^{1/3}$.
\end{enumerate}
\tcblower
\textbf{Solution.}
\begin{enumerate}
\item $n = 5$ est impair et $32 > 0$ : unique solution $x = \sqrt[5]{32} = \sqrt[5]{2^5} = 2$.
\item On regroupe sous un même radical grâce à $\sqrt[n]{a}\,\sqrt[n]{b} = \sqrt[n]{ab}$ :
\[ F = \dfrac{\sqrt[3]{54 \times 4}}{6^{1/3}} = \dfrac{\sqrt[3]{216}}{\sqrt[3]{6}} = \sqrt[3]{\dfrac{216}{6}} = \sqrt[3]{36}. \]
On peut aussi écrire $F = 36^{1/3} = 6^{2/3}$. (Vérification numérique : $\sqrt[3]{36} \approx 3{,}30$ ; or $\sqrt[3]{54} \approx 3{,}78$, $\sqrt[3]{4} \approx 1{,}59$, $6^{1/3} \approx 1{,}82$, et $\frac{3{,}78 \times 1{,}59}{1{,}82} \approx 3{,}30$ : cohérent.)
\end{enumerate}
\end{exoresolu}

\courssub{Dérivée de la fonction racine n-ième}

La fonction $g : x \mapsto \sqrt[n]{x}$ est la réciproque de $f : x \mapsto x^n$ sur $[0\,;\,+\infty[$. Appliquons la formule de dérivation des réciproques vue dans la section précédente : si $f$ est dérivable et si $f'(x) \neq 0$, alors
\[ \left(f^{-1}\right)'(y) = \dfrac{1}{f'\left(f^{-1}(y)\right)}. \]
Ici $f'(x) = n x^{n-1}$, qui ne s'annule qu'en $x = 0$. Pour $x > 0$, on obtient donc :

\begin{propriete}[{Dérivée de $x \mapsto \sqrt[n]{x}$}]
La fonction $g : x \mapsto \sqrt[n]{x}$ est dérivable sur $]0\,;\,+\infty[$ (mais \textbf{pas en $0$}) et :
\[ g'(x) = \dfrac{1}{n\,\left(\sqrt[n]{x}\right)^{n-1}} = \dfrac{1}{n}\,x^{\frac{1}{n} - 1}. \]
\end{propriete}

\begin{experience}[Démonstration guidée par la formule de la réciproque]
Posons $y = \sqrt[n]{x}$, c'est-à-dire $x = y^n = f(y)$. Pour $x > 0$, on a $y > 0$ donc $f'(y) = n y^{n-1} \neq 0$. La formule donne :
\[ g'(x) = \dfrac{1}{f'(g(x))} = \dfrac{1}{n\,\left(\sqrt[n]{x}\right)^{n-1}}. \]
Retrouvons la seconde écriture : $\left(\sqrt[n]{x}\right)^{n-1} = \left(x^{1/n}\right)^{n-1} = x^{(n-1)/n}$, donc
\[ g'(x) = \dfrac{1}{n}\,x^{-(n-1)/n} = \dfrac{1}{n}\,x^{\frac{1}{n}-1}. \]
Remarque : en $x = 0$, le taux d'accroissement $\dfrac{\sqrt[n]{h} - 0}{h} = \dfrac{1}{h^{(n-1)/n}}$ tend vers $+\infty$ quand $h \to 0^+$ : la courbe admet en l'origine une \textbf{tangente verticale} (symétrique de la tangente horizontale de $x \mapsto x^n$ en $0$), et $g$ n'y est pas dérivable.
\end{experience}

\begin{exemplebox}[Application : dérivée de la racine cubique]
Pour $n = 3$ : $\left(\sqrt[3]{x}\right)' = \dfrac{1}{3\left(\sqrt[3]{x}\right)^2} = \dfrac{1}{3}\,x^{-2/3}$ pour tout $x > 0$.\\
Pour $n = 2$, on retrouve la formule connue : $(\sqrt{x})' = \dfrac{1}{2\sqrt{x}}$.
\end{exemplebox}

\begin{savaistu}[Une belle cohérence de chapitre]
Remarque l'élégance du résultat : la formule $\left(x^{1/n}\right)' = \frac{1}{n}\,x^{\frac{1}{n}-1}$ est exactement la même que celle des puissances entières, $(x^n)' = n x^{n-1}$, mais appliquée à l'exposant rationnel $\frac{1}{n}$ ! Ce n'est pas un hasard : c'est la formule de dérivation des fonctions réciproques qui garantit cette cohérence. Au Baccalauréat, savoir retrouver rapidement $(\sqrt[n]{x})'$ en écrivant « $\sqrt[n]{x}$ est la réciproque de $x^n$, donc sa dérivée est $\frac{1}{n(\sqrt[n]{x})^{n-1}}$ » est une démarche très valorisée par les correcteurs, car elle montre que tu as compris l'architecture du chapitre plutôt qu'appris des formules par cœur. Historiquement, c'est Isaac Newton qui, au XVII\textsuperscript{e} siècle, a généralisé la notation $a^{p/q}$ aux exposants fractionnaires dans ses travaux sur les binômes — une notation que nous utilisons encore telle quelle aujourd'hui.
\end{savaistu}

\begin{methode}[Simplifier une expression avec des puissances rationnelles]
\begin{enumerate}
\item \textbf{Décomposer} chaque base en produit de facteurs premiers (par exemple $72 = 2^3 \times 3^2$).
\item \textbf{Convertir} les radicaux en exposants fractionnaires : $\sqrt[q]{a^p} = a^{p/q}$.
\item \textbf{Appliquer} les règles de calcul : $(a^r)^s = a^{rs}$, $a^r a^s = a^{r+s}$, $\dfrac{a^r}{a^s} = a^{r-s}$.
\item \textbf{Réduire} les exposants (sommes et différences de fractions, au même dénominateur).
\item \textbf{Revenir} si demandé à une écriture avec radicaux, ou donner la valeur exacte quand les exposants sont entiers.
\end{enumerate}
\end{methode}

\begin{exercice}[Pour t'entraîner]
\begin{enumerate}
\item Résoudre dans $\mathbb{R}$ : a) $x^3 = 27$ ; b) $x^4 = 81$ ; c) $x^6 = -1$ ; d) $x^3 = -\dfrac{1}{8}$.
\item Simplifier : $A = \dfrac{16^{3/4} \times 9^{1/2}}{27^{2/3}}$ et $B = \sqrt[3]{16} \times \sqrt[3]{4} \times 2^{-1/3}$.
\item Soit $g : x \mapsto \sqrt[4]{x}$. Déterminer $g'(x)$ pour $x > 0$, puis l'équation de la tangente à la courbe de $g$ au point d'abscisse $1$.
\end{enumerate}
\end{exercice}

\begin{corrige}[Exercice n°1]
\begin{enumerate}
\item a) $n = 3$ impair : unique solution $x = \sqrt[3]{27} = 3$. \quad b) $n = 4$ pair, $81 > 0$ : deux solutions $x = \sqrt[4]{81} = 3$ et $x = -3$ (car $81 = 3^4$). \quad c) $n = 6$ pair et $-1 < 0$ : aucune solution réelle. \quad d) $n = 3$ impair, $a < 0$ : unique solution $x = -\sqrt[3]{\frac{1}{8}} = -\dfrac{1}{2}$.
\item $A = \dfrac{(2^4)^{3/4} \times (3^2)^{1/2}}{(3^3)^{2/3}} = \dfrac{2^3 \times 3}{3^2} = \dfrac{8 \times 3}{9} = \dfrac{8}{3}$.\\
$B = \sqrt[3]{16 \times 4} \times 2^{-1/3} = \sqrt[3]{64} \times 2^{-1/3} = 4 \times 2^{-1/3} = 2^{2 - 1/3} = 2^{5/3} = \sqrt[3]{32}$.
\item $g'(x) = \dfrac{1}{4\left(\sqrt[4]{x}\right)^3} = \dfrac{1}{4}\,x^{-3/4}$. En $x = 1$ : $g(1) = 1$ et $g'(1) = \dfrac{1}{4}$. Tangente : $y = \dfrac{1}{4}(x - 1) + 1$, soit $y = \dfrac{x}{4} + \dfrac{3}{4}$.
\end{enumerate}
\end{corrige}

\begin{aretenir}[L'essentiel de la section]
\begin{itemize}
\item $x \mapsto x^n$ est une bijection de $[0\,;\,+\infty[$ sur $[0\,;\,+\infty[$ ; sa réciproque est la \cle{racine n-ième} $x \mapsto \sqrt[n]{x}$.
\item $\sqrt[n]{a}$ est toujours \textbf{positif} ; l'équation $x^n = a$ ($a > 0$) a \textbf{deux solutions} si $n$ est pair, \textbf{une seule} si $n$ est impair ; si $a < 0$, une solution si $n$ impair, aucune si $n$ pair.
\item $a^{p/q} = \left(\sqrt[q]{a}\right)^p$ pour $a > 0$, et toutes les règles de calcul des puissances entières se prolongent.
\item $\left(\sqrt[n]{x}\right)' = \dfrac{1}{n\left(\sqrt[n]{x}\right)^{n-1}}$ sur $]0\,;\,+\infty[$, obtenue par la dérivation des réciproques ; tangente verticale en $0$.
\item Pièges : $\sqrt[n]{a+b} \neq \sqrt[n]{a} + \sqrt[n]{b}$ ; ne pas oublier la solution négative quand $n$ est pair.
\end{itemize}
\end{aretenir}

\coursec{Inégalité des accroissements finis}

Dans la section précédente, nous avons appris à manipuler les racines $n$-ièmes et les puissances rationnelles : tu sais désormais donner un sens à $a^{1/n}$ et résoudre $x^n = a$. Mais une question pratique demeure : comment obtenir une \emph{valeur approchée} de $\sqrt{101}$ sans calculatrice, et surtout comment \emph{contrôler l'erreur} commise ? C'est exactement le genre de problème que résout l'\cle{inégalité des accroissements finis}, outil central de l'analyse de Terminale et vedette incontestée des problèmes de Baccalauréat C, D, E et TI.

L'idée intuitive est simple. Imagine un chauffeur de taxi-brousse qui roule entre Douala et Yaoundé. Si sa vitesse instantanée reste toujours comprise entre $\SI{60}{km/h}$ et $\SI{90}{km/h}$, alors la distance parcourue en $3$ heures est forcément comprise entre $180$ et $270$ kilomètres. La vitesse, c'est la dérivée ; la distance parcourue, c'est l'accroissement $f(b) - f(a)$. Encadrer la dérivée, c'est donc encadrer l'accroissement. Toute cette section consiste à rendre cette intuition rigoureuse et à l'exploiter.

\courssub{Rappel : le théorème des accroissements finis}

\begin{propriete}[Théorème des accroissements finis (TAF) — admis]
Soit $f$ une fonction \cle{continue} sur $[a\,;\,b]$ et \cle{dérivable} sur $]a\,;\,b[$. Alors il existe au moins un réel $c \in {}]a\,;\,b[$ tel que :
\[ f(b) - f(a) = f'(c)\,(b - a). \]
\end{propriete}

\begin{aretenir}[Interprétation géométrique]
Le quotient $\dfrac{f(b)-f(a)}{b-a}$ est la \cle{pente de la corde} $[AB]$, où $A(a\,;\,f(a))$ et $B(b\,;\,f(b))$. Le TAF affirme qu'il existe au moins un point $C(c\,;\,f(c))$ de la courbe où la \textbf{tangente est parallèle à la corde} $[AB]$.
\end{aretenir}

Le TAF donne une \emph{égalité}, mais il ne dit pas où se trouve le point $c$ : c'est un théorème d'existence. Pour obtenir des résultats quantitatifs (des encadrements chiffrés), on lui préfère son avatar inégalitaire, que nous allons démontrer.

\courssub{L'inégalité des accroissements finis}

\begin{propriete}[Inégalité des accroissements finis (IAF)]
Soit $f$ une fonction dérivable sur un intervalle $[a\,;\,b]$ (avec $a < b$). On suppose qu'il existe deux réels $m$ et $M$ tels que, pour tout $x \in [a\,;\,b]$ :
\[ m \leq f'(x) \leq M. \]
Alors :
\[ m\,(b-a) \;\leq\; f(b) - f(a) \;\leq\; M\,(b-a). \]
\end{propriete}

\begin{experience}[Démonstration guidée de l'IAF]
L'idée consiste à appliquer le lien entre signe de la dérivée et sens de variation à deux fonctions auxiliaires bien choisies.
\begin{enumerate}
\item \textbf{Première fonction auxiliaire.} Posons $g(x) = f(x) - m\,x$. Comme $f'(x) \geq m$ sur $[a\,;\,b]$, on a $g'(x) = f'(x) - m \geq 0$ : la fonction $g$ est donc \textbf{croissante} sur $[a\,;\,b]$.
\item \textbf{Comparaison des valeurs.} Puisque $a \leq b$ et que $g$ est croissante, $g(a) \leq g(b)$, c'est-à-dire $f(a) - m\,a \leq f(b) - m\,b$, d'où $f(b) - f(a) \geq m(b-a)$.
\item \textbf{Deuxième fonction auxiliaire.} Posons $h(x) = M\,x - f(x)$. Comme $f'(x) \leq M$, on a $h'(x) = M - f'(x) \geq 0$ : $h$ est croissante sur $[a\,;\,b]$, donc $h(a) \leq h(b)$, ce qui donne $f(b) - f(a) \leq M(b-a)$.
\item \textbf{Conclusion.} En regroupant : $m(b-a) \leq f(b) - f(a) \leq M(b-a)$. $\blacksquare$
\end{enumerate}
\end{experience}

\begin{popfigure}
\begin{center}
\begin{tikzpicture}
\begin{axis}[
  width=0.92\linewidth, height=7.5cm,
  axis lines=middle, xlabel={$x$}, ylabel={$y$},
  xmin=0, xmax=6.4, ymin=0, ymax=7.4,
  xtick={1,5}, xticklabels={$a$,$b$},
  ytick=\empty,
  legend style={at={(0.02,0.97)},anchor=north west,font=\small}
]
\addplot[popcurve,domain=0.4:6,samples=200]{0.35*(x-1)^2 + 1};
\addlegendentry{$y=f(x)$}
% Corde AB
\addplot[poporange,thick,domain=1:5]{1 + 1.4*(x-1)};
\addlegendentry{corde $[AB]$ : pente $\frac{f(b)-f(a)}{b-a}$}
% Tangentes de pentes extrêmes
\addplot[popgreen,thick,dashed,domain=0.4:2.6]{1 + 0.35*(x-1)};
\addlegendentry{tangente de pente minimale $m$}
\addplot[popink,thick,dashed,domain=3.4:6]{6.6 + 2.45*(x-5)};
\addlegendentry{tangente de pente maximale $M$}
% Points
\addplot[only marks,mark=*,popdark] coordinates {(1,1) (5,6.6)};
\node[popdark,anchor=east] at (axis cs:1,1) {$A$};
\node[popdark,anchor=south west] at (axis cs:5,6.6) {$B$};
\end{axis}
\end{tikzpicture}
\end{center}
\legende{Interprétation géométrique de l'IAF : la pente de la corde $[AB]$ est comprise entre la plus petite pente $m$ et la plus grande pente $M$ des tangentes à la courbe sur $[a\,;\,b]$.}
\end{popfigure}

\begin{aretenir}[Lecture géométrique]
Diviser l'encadrement $m(b-a) \leq f(b)-f(a) \leq M(b-a)$ par $b-a > 0$ donne :
\[ m \;\leq\; \frac{f(b)-f(a)}{b-a} \;\leq\; M. \]
Autrement dit : \textbf{la pente de toute corde est comprise entre les pentes extrêmes des tangentes}. C'est exactement ce que traduit la figure ci-dessus.
\end{aretenir}

\courssub{Cas particulier fondamental : fonctions lipschitziennes}

\begin{propriete}[Cas $|f'| \leq k$]
Soit $f$ une fonction dérivable sur un intervalle $I$ et $k$ un réel positif. Si pour tout $x \in I$, $|f'(x)| \leq k$, alors pour tous réels $a$ et $b$ de $I$ :
\[ \bigl|f(b) - f(a)\bigr| \;\leq\; k\,|b - a|. \]
On dit que $f$ est \cle{$k$-lipschitzienne} sur $I$.
\end{propriete}

\textbf{Démonstration.} L'hypothèse $|f'(x)| \leq k$ s'écrit $-k \leq f'(x) \leq k$. Supposons $a < b$ et appliquons l'IAF avec $m = -k$ et $M = k$ :
\[ -k(b-a) \leq f(b) - f(a) \leq k(b-a), \]
ce qui équivaut exactement à $|f(b) - f(a)| \leq k(b-a) = k|b-a|$. Si $b < a$, on échange les rôles de $a$ et $b$ ; si $a = b$, l'inégalité est triviale. $\blacksquare$

\begin{attention}[Conditions d'application à ne jamais oublier]
Deux erreurs reviennent sans cesse dans les copies :
\begin{itemize}
\item \textbf{Appliquer l'IAF sans vérifier la dérivabilité sur tout l'intervalle.} Par exemple, $x \mapsto \sqrt{x}$ n'est pas dérivable en $0$ : on ne peut pas appliquer l'IAF sur $[0\,;\,1]$, mais on peut l'appliquer sur $[1\,;\,4]$ sans problème.
\item \textbf{Encadrer $f'$ sur un intervalle plus grand que $[a\,;\,b]$.} Si tu encadres $f'$ sur $[0\,;\,10]$ alors que tu travailles sur $[2\,;\,3]$, ton encadrement reste vrai mais il est \emph{grossier} : tu perds en précision. Encadre toujours $f'$ sur l'intervalle exact $[a\,;\,b]$, en t'appuyant si besoin sur le tableau de variations de $f'$.
\end{itemize}
\end{attention}

\courssub{Méthode et application : encadrement d'une valeur approchée}

\begin{methode}[Encadrer $f(b) - f(a)$ grâce à l'IAF]
\begin{enumerate}
\item Vérifier que $f$ est \textbf{dérivable} sur $[a\,;\,b]$ (et le rédiger !).
\item Calculer $f'(x)$ et \textbf{encadrer $f'$ sur $[a\,;\,b]$} : $m \leq f'(x) \leq M$. Pour cela, on étudie le sens de variation de $f'$ (souvent via le signe de $f''$) ou on utilise des propriétés de monotonie (décroissance de $x \mapsto \dfrac{1}{2\sqrt{x}}$, etc.).
\item \textbf{Multiplier} l'encadrement par $(b-a) > 0$ : $m(b-a) \leq f(b) - f(a) \leq M(b-a)$.
\item Conclure en isolant la quantité cherchée.
\end{enumerate}
\end{methode}

\begin{exemplebox}[Encadrement de $\sqrt{101}$]
Montrons que $0{,}0495 \leq \sqrt{101} - 10 \leq 0{,}05$.

\textbf{Étape 1.} La fonction $f : x \mapsto \sqrt{x}$ est dérivable sur $[100\,;\,101]$ (et même sur $]0\,;\,+\infty[$), avec $f'(x) = \dfrac{1}{2\sqrt{x}}$.

\textbf{Étape 2.} La fonction $x \mapsto \sqrt{x}$ est croissante, donc $f'$ est \textbf{décroissante} sur $[100\,;\,101]$. Ainsi, pour tout $x \in [100\,;\,101]$ :
\[ \frac{1}{2\sqrt{101}} \;\leq\; f'(x) \;\leq\; \frac{1}{2\sqrt{100}} = \frac{1}{20} = 0{,}05. \]
Pour la borne inférieure, on utilise $\sqrt{101} \leq \sqrt{102{,}01} = 10{,}1$, d'où, par décroissance de la fonction inverse sur $]0\,;\,+\infty[$ :
\[ \frac{1}{2\sqrt{101}} \geq \frac{1}{2 \times 10{,}1} = \frac{1}{20{,}2} \geq 0{,}0495 \quad \text{car } 1 \geq 0{,}0495 \times 20{,}2 = 0{,}9999. \]

\textbf{Étape 3.} L'IAF avec $b - a = 1$ donne :
\[ 0{,}0495 \;\leq\; \sqrt{101} - \sqrt{100} \;\leq\; 0{,}05. \]

\textbf{Étape 4.} Comme $\sqrt{100} = 10$, on obtient $10{,}0495 \leq \sqrt{101} \leq 10{,}05$, soit un encadrement d'amplitude $0{,}0005$ — remarquablement précis pour un calcul fait à la main ! (La valeur exacte est $\sqrt{101} \approx 10{,}04987$ : elle est bien dans l'encadrement.)
\end{exemplebox}

\begin{exoresolu}[Encadrer $\sqrt{101} - 10$ à $10^{-3}$ près]
\textbf{Énoncé.} Soit $f(x) = \sqrt{x}$. En appliquant l'inégalité des accroissements finis à $f$ sur $[100\,;\,101]$, montrer que :
\[ 0{,}0497 \leq \sqrt{101} - 10 \leq 0{,}05. \]
(On admettra que $\sqrt{101} \leq 10{,}05$.)
\tcblower
\textbf{Solution détaillée.}

\emph{Étape 1 : dérivabilité.} $f$ est dérivable sur $]0\,;\,+\infty[$, donc en particulier sur $[100\,;\,101]$, et $f'(x) = \dfrac{1}{2\sqrt{x}}$.

\emph{Étape 2 : encadrement de $f'$.} La fonction $f'$ est décroissante sur $[100\,;\,101]$ car $x \mapsto \sqrt{x}$ y est croissante. Donc pour tout $x \in [100\,;\,101]$ :
\[ \frac{1}{2\sqrt{101}} \leq f'(x) \leq \frac{1}{2\sqrt{100}}. \]
D'une part $\dfrac{1}{2\sqrt{100}} = \dfrac{1}{20} = 0{,}05$. D'autre part, de $\sqrt{101} \leq 10{,}05$ on tire $2\sqrt{101} \leq 20{,}1$, puis par décroissance de la fonction inverse sur $]0\,;\,+\infty[$ :
\[ \frac{1}{2\sqrt{101}} \geq \frac{1}{20{,}1} = \frac{10}{201} \geq 0{,}0497 \quad \text{car } 10 \geq 0{,}0497 \times 201 = 9{,}9897. \]

\emph{Étape 3 : application de l'IAF.} Avec $m = 0{,}0497$, $M = 0{,}05$ et $b - a = 1$ :
\[ 0{,}0497 \times 1 \;\leq\; f(101) - f(100) \;\leq\; 0{,}05 \times 1, \]
c'est-à-dire $0{,}0497 \leq \sqrt{101} - 10 \leq 0{,}05$. $\blacksquare$
\end{exoresolu}

\courssub{Application aux suites récurrentes $u_{n+1} = f(u_n)$ (complément Bac)}

Voici l'application la plus stratégique de l'IAF au Baccalauréat. On considère une suite définie par $u_{n+1} = f(u_n)$, et un réel $\ell$ tel que $f(\ell) = \ell$ (on dit que $\ell$ est un \cle{point fixe} de $f$). Supposons que $|f'(x)| \leq k$ sur un intervalle $I$ contenant $\ell$ et tous les termes de la suite, avec $k < 1$. Alors, pour tout entier $n$ :
\[ |u_{n+1} - \ell| = \bigl|f(u_n) - f(\ell)\bigr| \;\leq\; k\,|u_n - \ell|. \]
Par récurrence immédiate :
\[ |u_n - \ell| \;\leq\; k^n\,|u_0 - \ell|. \]
Comme $0 \leq k < 1$, $k^n \to 0$, donc $|u_n - \ell| \to 0$ : \textbf{la suite converge vers $\ell$}. Mieux : l'inégalité donne la \emph{vitesse} de convergence — on dit qu'elle est \cle{géométrique} de raison $k$ — et permet de \textbf{prévoir le nombre d'itérations} nécessaires pour atteindre une précision donnée.

\begin{popfigure}
\begin{center}
\begin{tikzpicture}[scale=1.15]
% Axes
\draw[->,popdark] (0,0) -- (6.2,0) node[below left]{$x$};
\draw[->,popdark] (0,0) -- (0,5.2) node[below left]{$y$};
% Droite y = x
\draw[popgreen,thick] (0,0) -- (5,5) node[pos=0.92,sloped,above,font=\small]{$y=x$};
% Courbe y = f(x), contractante, point fixe en 3
\draw[popblue,very thick,domain=0.6:5.6,samples=100] plot (\x,{3 + 0.45*(\x-3) + 0.35*sin(deg(1.5*(\x-3)))}) node[pos=0.95,above,font=\small]{$y=f(x)$};
% Point fixe
\fill[popdark] (3,3) circle (1.8pt) node[below right,font=\small]{$\ell$};
\draw[dashed,popmuted] (3,0) -- (3,3) -- (0,3);
\node[below,font=\small,popmuted] at (3,0) {$\ell$};
% Toile d'araignée : u0 = 1.1
\pgfmathsetmacro{\u}{1.1}
\pgfmathsetmacro{\fu}{3 + 0.45*(\u-3) + 0.35*sin(deg(1.5*(\u-3)))}
\draw[poporange,thick,->] (\u,0) -- (\u,\fu);
\draw[dashed,popmuted] (\u,0) node[below,font=\small]{$u_0$} -- (\u,0.02);
\pgfmathsetmacro{\v}{\fu}
\pgfmathsetmacro{\fv}{3 + 0.45*(\v-3) + 0.35*sin(deg(1.5*(\v-3)))}
\draw[poporange,thick,->] (\u,\fu) -- (\v,\fu);
\draw[poporange,thick,->] (\v,\fu) -- (\v,\fv);
\pgfmathsetmacro{\w}{\fv}
\pgfmathsetmacro{\fw}{3 + 0.45*(\w-3) + 0.35*sin(deg(1.5*(\w-3)))}
\draw[poporange,thick,->] (\v,\fv) -- (\w,\fv);
\draw[poporange,thick,->] (\w,\fv) -- (\w,\fw);
\pgfmathsetmacro{\z}{\fw}
\pgfmathsetmacro{\fz}{3 + 0.45*(\z-3) + 0.35*sin(deg(1.5*(\z-3)))}
\draw[poporange,thick,->] (\w,\fw) -- (\z,\fw);
\draw[poporange,thick,->] (\z,\fw) -- (\z,\fz);
% Labels u1, u2
\node[below,font=\small,poporange] at (\v,0) {$u_1$};
\draw[dashed,popmuted] (\v,0) -- (\v,0.02);
\node[below,font=\small,poporange] at (\w,0) {$u_2$};
\draw[dashed,popmuted] (\w,0) -- (\w,0.02);
\end{tikzpicture}
\end{center}
\legende{Toile d'araignée (ou spirale d'escargot) d'une suite récurrente $u_{n+1} = f(u_n)$ : quand $|f'| \leq k < 1$ au voisinage du point fixe $\ell$, les marches de l'escalier (ou de la spirale) rétrécissent et les termes $u_n$ se rapprochent inexorablement de $\ell$.}
\end{popfigure}

\begin{methode}[Prouver la convergence d'une suite récurrente par l'IAF]
\begin{enumerate}
\item Déterminer le point fixe : résoudre $f(\ell) = \ell$.
\item Montrer qu'il existe un intervalle $I$ \textbf{stable} par $f$ (c'est-à-dire $f(I) \subset I$) contenant $\ell$ et $u_0$ ; en déduire par récurrence que tous les $u_n$ sont dans $I$.
\item Montrer que $|f'(x)| \leq k$ sur $I$ avec $k < 1$.
\item Appliquer l'IAF entre $u_n$ et $\ell$ : $|u_{n+1} - \ell| \leq k\,|u_n - \ell|$, puis conclure $|u_n - \ell| \leq k^n |u_0 - \ell| \xrightarrow[n\to+\infty]{} 0$.
\end{enumerate}
\end{methode}

\begin{exoresolu}[Convergence d'une suite récurrente]
\textbf{Énoncé.} Soit $f(x) = \dfrac{1}{2}\left(x + \dfrac{2}{x}\right)$ et la suite définie par $u_0 = 2$ et $u_{n+1} = f(u_n)$. On admet que $f(\sqrt{2}) = \sqrt{2}$ et que $u_n \in [\sqrt{2}\,;\,2]$ pour tout $n$.
\begin{enumerate}
\item Montrer que pour tout $x \in [\sqrt{2}\,;\,2]$, $|f'(x)| \leq \dfrac{1}{4}$.
\item En déduire que $|u_n - \sqrt{2}| \leq \dfrac{2 - \sqrt{2}}{4^n}$ et conclure.
\end{enumerate}
\tcblower
\textbf{Solution détaillée.}

\emph{1. Majoration de $|f'|$.} Pour $x > 0$ : $f'(x) = \dfrac{1}{2}\left(1 - \dfrac{2}{x^2}\right)$. Sur $[\sqrt{2}\,;\,2]$, la fonction $x \mapsto -\dfrac{2}{x^2}$ est croissante (car $x \mapsto \dfrac{1}{x^2}$ est décroissante), donc $f'$ est croissante, avec :
\[ f'(\sqrt{2}) = \frac{1}{2}\left(1 - \frac{2}{2}\right) = 0 \quad \text{et} \quad f'(2) = \frac{1}{2}\left(1 - \frac{2}{4}\right) = \frac{1}{4}. \]
Ainsi $0 \leq f'(x) \leq \dfrac{1}{4}$ sur $[\sqrt{2}\,;\,2]$, donc $|f'(x)| \leq \dfrac{1}{4}$.

\emph{2. Convergence.} $f$ est dérivable sur $[\sqrt{2}\,;\,2]$ et $|f'| \leq \dfrac14$ : l'IAF appliquée entre $u_n$ et $\sqrt{2}$ (qui appartiennent tous deux à $[\sqrt{2}\,;\,2]$) donne :
\[ |u_{n+1} - \sqrt{2}| = \bigl|f(u_n) - f(\sqrt{2})\bigr| \leq \frac{1}{4}\,|u_n - \sqrt{2}|. \]
Par récurrence : $|u_n - \sqrt{2}| \leq \left(\dfrac{1}{4}\right)^n |u_0 - \sqrt{2}| = \dfrac{2 - \sqrt{2}}{4^n}$. Comme $\dfrac{1}{4^n} \to 0$, le théorème des gendarmes donne $|u_n - \sqrt{2}| \to 0$ : la suite \textbf{converge vers $\sqrt{2}$}.

\emph{Estimation d'erreur.} Pour obtenir $\sqrt{2}$ à $10^{-6}$ près, il suffit que $\dfrac{0{,}6}{4^n} \leq 10^{-6}$, soit $4^n \geq 6 \times 10^{5}$ : $n = 10$ convient largement ($4^{10} > 10^{6}$). C'est la \emph{méthode de Héron}, utilisée depuis l'Antiquité pour extraire des racines carrées !
\end{exoresolu}

\begin{savaistu}[L'IAF, star des problèmes de Bac]
L'inégalité des accroissements finis est l'un des outils les plus fréquemment exigés dans les problèmes de Baccalauréat C et E au Cameroun : elle apparaît presque chaque année, le plus souvent dans l'étude d'une suite récurrente $u_{n+1} = f(u_n)$ associée à une fonction étudiée dans les questions précédentes. Le schéma « point fixe + majoration de $|f'|$ par $k < 1$ + récurrence » est un grand classique. La méthode de Héron vue ci-dessus était déjà connue des mathématiciens babyloniens il y a près de 4000 ans pour calculer des racines carrées : c'est elle qui, cachée dans les microprocesseurs, permet encore aujourd'hui aux calculatrices et aux téléphones d'évaluer $\sqrt{x}$ en une fraction de seconde.
\end{savaistu}

\begin{exercice}[À toi de jouer]
\textbf{Exercice 1.} En appliquant l'IAF à $f(x) = \sqrt{x}$ sur $[49\,;\,50]$, montrer que $7{,}070 \leq \sqrt{50} \leq 7{,}072$ (on admettra $\sqrt{50} \leq 7{,}1$).

\textbf{Exercice 2.} Soit $f(x) = \dfrac{1}{3}\cos(x) + 1$ et $(u_n)$ définie par $u_0 = 0$ et $u_{n+1} = f(u_n)$. On admet que l'équation $f(x) = x$ admet une unique solution $\ell$ dans $\mathbb{R}$.
\begin{enumerate}
\item Montrer que pour tout réel $x$, $|f'(x)| \leq \dfrac{1}{3}$.
\item En déduire que $|u_{n+1} - \ell| \leq \dfrac{1}{3}|u_n - \ell|$, puis que $(u_n)$ converge vers $\ell$.
\item Déterminer un entier $n$ tel que $|u_n - \ell| \leq 10^{-3}$.
\end{enumerate}
\end{exercice}

\begin{corrige}[Exercice 1]
$f$ est dérivable sur $[49\,;\,50]$ et $f'(x) = \dfrac{1}{2\sqrt{x}}$ est décroissante sur cet intervalle (car $x \mapsto \sqrt{x}$ y est croissante). Donc pour tout $x \in [49\,;\,50]$ :
\[ \frac{1}{2\sqrt{50}} \leq f'(x) \leq \frac{1}{2\sqrt{49}} = \frac{1}{14}. \]
Pour la borne supérieure : $\dfrac{1}{14} \leq 0{,}0715$ car $1 \leq 0{,}0715 \times 14 = 1{,}001$.

Pour la borne inférieure : de $\sqrt{50} \leq 7{,}1$ (admis) on tire $2\sqrt{50} \leq 14{,}2$, puis, par décroissance de la fonction inverse sur $]0\,;\,+\infty[$ :
\[ \frac{1}{2\sqrt{50}} \geq \frac{1}{14{,}2} \geq 0{,}0704 \quad \text{car } 1 \geq 0{,}0704 \times 14{,}2 = 0{,}99968. \]
L'IAF avec $b - a = 1$ donne alors :
\[ 0{,}0704 \;\leq\; \sqrt{50} - \sqrt{49} \;\leq\; 0{,}0715, \]
c'est-à-dire, puisque $\sqrt{49} = 7$ :
\[ 7{,}0704 \leq \sqrt{50} \leq 7{,}0715, \]
ce qui implique a fortiori l'encadrement demandé : $7{,}070 \leq \sqrt{50} \leq 7{,}072$. (La valeur exacte est $\sqrt{50} \approx 7{,}0711$.)
\end{corrige}

\begin{corrige}[Exercice 2]
\begin{enumerate}
\item $f'(x) = -\dfrac{1}{3}\sin(x)$, donc $|f'(x)| = \dfrac{1}{3}|\sin x| \leq \dfrac{1}{3}$ pour tout réel $x$.
\item $f$ est dérivable sur $\mathbb{R}$ avec $|f'| \leq \dfrac13$ : par l'IAF appliquée entre $u_n$ et $\ell$ (avec $f(\ell) = \ell$) :
\[ |u_{n+1} - \ell| = |f(u_n) - f(\ell)| \leq \frac{1}{3}|u_n - \ell|. \]
Par récurrence, $|u_n - \ell| \leq \dfrac{|u_0 - \ell|}{3^n} = \dfrac{|\ell|}{3^n}$. Comme $\dfrac{1}{3^n} \to 0$, le théorème des gendarmes donne $|u_n - \ell| \to 0$ : $(u_n)$ converge vers $\ell$.
\item On peut majorer $|\ell|$ : comme $f(x) \in \left[\dfrac23\,;\,\dfrac43\right]$ pour tout $x$, le point fixe vérifie $\ell \in \left[\dfrac23\,;\,\dfrac43\right]$, donc $|\ell| \leq \dfrac43 \leq 2$. Il suffit alors que $\dfrac{2}{3^n} \leq 10^{-3}$, soit $3^n \geq 2000$. Or $3^6 = 729 < 2000$ et $3^7 = 2187 \geq 2000$ : $n = 7$ convient.
\end{enumerate}
\end{corrige}

\begin{aretenir}[L'essentiel de la section]
\begin{itemize}
\item \textbf{IAF} : si $m \leq f'(x) \leq M$ sur $[a\,;\,b]$, alors $m(b-a) \leq f(b) - f(a) \leq M(b-a)$.
\item \textbf{Cas lipschitzien} : si $|f'(x)| \leq k$, alors $|f(b) - f(a)| \leq k|b-a|$.
\item \textbf{Suites récurrentes} : si $|f'| \leq k < 1$ sur un intervalle stable contenant le point fixe $\ell$, alors $|u_n - \ell| \leq k^n |u_0 - \ell| \to 0$.
\item Toujours \textbf{vérifier la dérivabilité} et \textbf{encadrer $f'$ sur l'intervalle exact} avant d'appliquer l'IAF.
\end{itemize}
\end{aretenir}

\coursec{Limites usuelles et branches infinies}

Dans la section précédente, l'inégalité des accroissements finis nous a permis de contrôler finement les variations d'une fonction entre deux points. Mais que se passe-t-il lorsque la variable s'approche d'une valeur interdite, ou lorsqu'elle s'éloigne vers l'infini ? La courbe peut alors « filer » vers le haut, se coller contre une droite horizontale, ou suivre de très près une droite inclinée : ce sont les \cle{branches infinies}. Pour les étudier, nous avons besoin d'une boîte à outils de limites usuelles, notamment trigonométriques et irrationnelles. C'est l'objet de cette section, capitale pour l'étude complète de fonctions au Baccalauréat.

\courssub{Limites usuelles des fonctions trigonométriques}

\textbf{Une intuition géométrique.} Sur le cercle trigonométrique de rayon $1$, prends un angle $x$ (en radians) petit et positif. La longueur de l'arc correspondant vaut $x$, tandis que la hauteur du point image sur le cercle vaut $\sin x$. Quand $x$ devient très petit, l'arc et la corde se « confondent » : le rapport $\dfrac{\sin x}{x}$ doit donc tendre vers $1$. C'est cette intuition que formalise le théorème suivant.

\begin{propriete}[Limites trigonométriques usuelles]
\[
\lim_{x \to 0} \frac{\sin x}{x} = 1 \qquad ; \qquad
\lim_{x \to 0} \frac{\tan x}{x} = 1 \qquad ; \qquad
\lim_{x \to 0} \frac{1 - \cos x}{x^2} = \frac{1}{2} \qquad ; \qquad
\lim_{x \to 0} \frac{\cos x - 1}{x} = 0.
\]
\end{propriete}

\begin{experience}[Démonstration guidée : $\lim_{x\to 0}\dfrac{\sin x}{x}=1$]
On raisonne pour $x \in \left]0\,;\,\dfrac{\pi}{2}\right[$, le cas $x<0$ s'obtenant par parité (la fonction $x\mapsto\dfrac{\sin x}{x}$ est paire).
\begin{enumerate}
\item Sur le cercle trigonométrique, compare les trois aires : celle du triangle $OAM$ (avec $M$ le point image de $x$), celle du secteur angulaire d'angle $x$, et celle du triangle $OAT$ où $T$ est l'intersection de la demi-droite $[OM)$ avec la tangente au cercle en $A(1\,;0)$.
\item On obtient : $\dfrac{1}{2}\sin x \leq \dfrac{x}{2} \leq \dfrac{1}{2}\tan x$, c'est-à-dire $\sin x \leq x \leq \tan x$.
\item En divisant par $\sin x > 0$ : $1 \leq \dfrac{x}{\sin x} \leq \dfrac{1}{\cos x}$, puis en passant aux inverses (tous les termes sont strictement positifs) :
\[
\cos x \leq \frac{\sin x}{x} \leq 1.
\]
\item Comme $\lim_{x\to 0}\cos x = 1$, le théorème des gendarmes donne $\lim_{x\to 0}\dfrac{\sin x}{x}=1$.
\end{enumerate}
Les trois autres limites s'en déduisent algébriquement :
\[
\frac{\tan x}{x} = \frac{\sin x}{x}\times\frac{1}{\cos x} \xrightarrow[x\to0]{} 1\times 1 = 1 ;
\]
\[
\frac{1-\cos x}{x^2} = \frac{(1-\cos x)(1+\cos x)}{x^2(1+\cos x)} = \frac{\sin^2 x}{x^2}\times\frac{1}{1+\cos x} = \left(\frac{\sin x}{x}\right)^2 \times \frac{1}{1+\cos x} \xrightarrow[x\to0]{} 1\times\frac{1}{2} = \frac{1}{2} ;
\]
\[
\frac{\cos x - 1}{x} = -x \times \frac{1-\cos x}{x^2} \xrightarrow[x\to0]{} 0 \times \frac{1}{2} = 0.
\]
\end{experience}

\begin{methode}[Se ramener à $0$ par changement de variable]
Pour calculer une limite trigonométrique quand $x \to a$, on pose $X = x - a$ (alors $X \to 0$) et on utilise les formules d'addition pour faire apparaître $\sin X$ ou $1 - \cos X$.
\end{methode}

\begin{exemplebox}[Limites trigonométriques calculées]
\textbf{a)} Calculons $\lim_{x\to 0}\dfrac{\sin 3x}{x}$. On écrit $\dfrac{\sin 3x}{x} = 3\times\dfrac{\sin 3x}{3x}$. En posant $X = 3x$, on a $X \to 0$, donc la limite vaut $3 \times 1 = 3$.

\textbf{b)} Calculons $\lim_{x\to 0}\dfrac{1-\cos x}{x^2}$ : c'est la limite usuelle, elle vaut $\dfrac{1}{2}$.

\textbf{c)} Calculons $\lim_{x\to \pi/2}\dfrac{\cos x}{x - \pi/2}$. Posons $X = x - \dfrac{\pi}{2}$. Alors $\cos x = \cos\left(X + \dfrac{\pi}{2}\right) = -\sin X$, donc le quotient devient $\dfrac{-\sin X}{X} \to -1$.
\end{exemplebox}

\begin{attention}[Radians obligatoires]
Ces limites ne sont valables que si $x$ est exprimé en \cle{radians}. De plus, $\lim_{x\to 0}\dfrac{\sin x}{x}=1$ ne se généralise pas naïvement : $\dfrac{\sin 3x}{x}$ tend vers $3$, pas vers $1$. Pense toujours à faire apparaître exactement le même argument au numérateur et au dénominateur.
\end{attention}

\courssub{Limites des fonctions irrationnelles}

\textbf{Le problème.} Une fonction comme $f(x) = \sqrt{x^2 + x} - x$ donne, en $+\infty$, une forme indéterminée « $+\infty - \infty$ ». Deux techniques permettent de lever l'indétermination.

\begin{methode}[Quantité conjuguée pour les formes $\sqrt{u}-\sqrt{v}$]
\begin{enumerate}
\item Multiplie numérateur et dénominateur par la \cle{quantité conjuguée} $\sqrt{u}+\sqrt{v}$.
\item Utilise l'identité $(\sqrt{u}-\sqrt{v})(\sqrt{u}+\sqrt{v}) = u - v$.
\item La nouvelle expression n'est en général plus indéterminée : conclus.
\end{enumerate}
\end{methode}

\begin{exemplebox}[Calcul de $\lim_{x\to+\infty}\left(\sqrt{x^2+x}-x\right)$]
\begin{align*}
\sqrt{x^2+x}-x &= \frac{\left(\sqrt{x^2+x}-x\right)\left(\sqrt{x^2+x}+x\right)}{\sqrt{x^2+x}+x} = \frac{x^2+x-x^2}{\sqrt{x^2+x}+x} = \frac{x}{\sqrt{x^2+x}+x}.
\end{align*}
Pour $x>0$, $\sqrt{x^2+x} = x\sqrt{1+\frac{1}{x}}$, donc
\[
\frac{x}{\sqrt{x^2+x}+x} = \frac{x}{x\sqrt{1+\frac{1}{x}}+x} = \frac{1}{\sqrt{1+\frac{1}{x}}+1} \xrightarrow[x\to+\infty]{} \frac{1}{2}.
\]
\end{exemplebox}

\begin{methode}[Factorisation par le terme prépondérant en $\pm\infty$]
Sous un radical, factorise le terme dominant, puis sors-le de la racine en n'oubliant pas la valeur absolue : pour $x>0$, $\sqrt{x^2 u(x)} = x\sqrt{u(x)}$ ; pour $x<0$, $\sqrt{x^2 u(x)} = -x\sqrt{u(x)}$. Cette étape est indispensable avant tout calcul de $\dfrac{f(x)}{x}$.
\end{methode}

\begin{attention}[Le piège de la valeur absolue]
Pour $x \to -\infty$, $\sqrt{x^2} = |x| = -x$ et non $x$. Ainsi $\dfrac{\sqrt{x^2+1}}{x} = \dfrac{|x|\sqrt{1+1/x^2}}{x}$ tend vers $1$ en $+\infty$ mais vers $-1$ en $-\infty$. Oublier cela fausse toute l'étude des branches infinies.
\end{attention}

\courssub{Branches infinies : asymptotes verticales et horizontales}

Lorsqu'une courbe « s'échappe » à l'infini, elle peut le faire en se rapprochant indéfiniment d'une droite. Cette droite est appelée \cle{asymptote}.

\begin{definition}[Asymptote verticale]
Si $\lim_{x\to a} f(x) = +\infty$ ou $-\infty$ (limite à gauche ou à droite), alors la droite d'équation $x = a$ est \cle{asymptote verticale} à la courbe $\mathcal{C}_f$.
\end{definition}

\begin{definition}[Asymptote horizontale]
Si $\lim_{x\to +\infty} f(x) = b$ (ou en $-\infty$) avec $b$ réel, alors la droite d'équation $y = b$ est \cle{asymptote horizontale} à $\mathcal{C}_f$ en $+\infty$ (resp. $-\infty$).
\end{definition}

\begin{popfigure}
\begin{center}
\begin{tikzpicture}
\begin{axis}[
  width=0.8\linewidth, height=6.5cm,
  axis lines=middle, xlabel={$x$}, ylabel={$y$},
  xmin=-1, xmax=7, ymin=-6, ymax=8,
  xtick={2}, ytick={1},
  xticklabels={$2$}, yticklabels={$1$},
  clip=true]
\addplot[popcurve, domain=-1:1.85, samples=100]{(x)/(x-2)};
\addplot[popcurve, domain=2.15:7, samples=100]{(x)/(x-2)};
\addplot[poporange, dashed, thick] coordinates {(2,-6) (2,8)};
\addplot[popgreen, dashed, thick, domain=-1:7]{1};
\node[poporange, anchor=west] at (axis cs:2.1,-4.5) {$x=2$};
\node[popgreen, anchor=south west] at (axis cs:5.2,1.05) {$y=1$};
\node[popblue, anchor=west] at (axis cs:3.1,4.6) {$y=\dfrac{x}{x-2}$};
\end{axis}
\end{tikzpicture}
\end{center}
\legende{L'hyperbole $y=\dfrac{x}{x-2}$ : asymptote verticale $x=2$ (la limite est infinie) et asymptote horizontale $y=1$ (la limite en $\pm\infty$ vaut $1$).}
\end{popfigure}

\courssub{Asymptote oblique}

\textbf{Intuition.} Quand $f(x)$ se comporte en $+\infty$ « presque comme » une fonction affine $ax+b$, la courbe se plaque contre la droite $y=ax+b$ : l'écart vertical $f(x)-(ax+b)$ tend vers $0$.

\begin{definition}[Asymptote oblique]
Si $\lim_{x\to+\infty}\big[f(x)-(ax+b)\big]=0$ (ou en $-\infty$), avec $a \neq 0$, alors la droite d'équation $y = ax+b$ est \cle{asymptote oblique} à $\mathcal{C}_f$.
\end{definition}

\begin{methode}[Détermination d'une asymptote oblique]
\begin{enumerate}
\item Calcule $a = \lim_{x\to\pm\infty}\dfrac{f(x)}{x}$. Si $a$ est un réel \textbf{non nul}, continue ; si $a=0$ et $f$ a une limite finie, c'est une asymptote horizontale ; si $a = \pm\infty$, c'est une branche parabolique de direction $(Oy)$.
\item Calcule $b = \lim_{x\to\pm\infty}\big[f(x)-ax\big]$. Si $b$ est un réel fini, alors $y=ax+b$ est asymptote oblique.
\item \textbf{Position relative} : étudie le signe de $f(x)-(ax+b)$. S'il est positif au voisinage de l'infini, $\mathcal{C}_f$ est \emph{au-dessus} de l'asymptote ; s'il est négatif, elle est \emph{en dessous}.
\end{enumerate}
\end{methode}

\begin{attention}[Le piège classique du Bac]
Il ne suffit \textbf{pas} que $\dfrac{f(x)}{x}$ ait une limite finie $a$ pour conclure à une asymptote oblique ! Il faut encore vérifier que $f(x)-ax$ a une limite \emph{finie}. Contre-exemple : $f(x) = x + \sqrt{x}$. On a $\dfrac{f(x)}{x} = 1 + \dfrac{1}{\sqrt{x}} \to 1$, mais $f(x) - x = \sqrt{x} \to +\infty$ : il n'y a pas d'asymptote oblique, seulement une \cle{direction asymptotique} de pente $1$ (branche parabolique de direction la droite $y=x$).
\end{attention}

\begin{exoresolu}[Asymptotes de $f(x)=\dfrac{2x^2-3x+1}{x+1}$]
\textbf{Énoncé.} Soit $f(x)=\dfrac{2x^2-3x+1}{x+1}$. Déterminer toutes les asymptotes de $\mathcal{C}_f$ et préciser la position relative de la courbe par rapport à son asymptote oblique.
\tcblower
\textbf{Solution.} $f$ est définie sur $\mathbb{R}\setminus\{-1\}$.

\textbf{Asymptote verticale.} En $-1$, le numérateur vaut $2+3+1=6\neq 0$ et le dénominateur s'annule : $\lim_{x\to -1^-}f(x)=-\infty$ et $\lim_{x\to -1^+}f(x)=+\infty$. La droite $x=-1$ est asymptote verticale.

\textbf{Asymptote oblique.} En $\pm\infty$, $f(x)$ se comporte comme $\dfrac{2x^2}{x}=2x$, donc
\[
a=\lim_{x\to\pm\infty}\frac{f(x)}{x}=\lim_{x\to\pm\infty}\frac{2x^2-3x+1}{x^2+x}=2.
\]
Ensuite :
\[
f(x)-2x=\frac{2x^2-3x+1-2x(x+1)}{x+1}=\frac{-5x+1}{x+1}\xrightarrow[x\to\pm\infty]{}-5.
\]
Donc $b=-5$ : la droite $y=2x-5$ est asymptote oblique en $+\infty$ et en $-\infty$.

\textbf{Position relative.} $f(x)-(2x-5)=\dfrac{-5x+1}{x+1}+5=\dfrac{-5x+1+5x+5}{x+1}=\dfrac{6}{x+1}$.
\begin{itemize}
\item Si $x>-1$ (en particulier en $+\infty$), $\dfrac{6}{x+1}>0$ : la courbe est \emph{au-dessus} de son asymptote.
\item Si $x<-1$ (en particulier en $-\infty$), $\dfrac{6}{x+1}<0$ : la courbe est \emph{en dessous} de son asymptote.
\end{itemize}
\end{exoresolu}

\begin{popfigure}
\begin{center}
\begin{tikzpicture}
\begin{axis}[
  width=0.85\linewidth, height=7.5cm,
  axis lines=middle, xlabel={$x$}, ylabel={$y$},
  xmin=-4, xmax=8, ymin=-8, ymax=16,
  xtick={1}, ytick={2},
  xticklabels={$1$}, yticklabels={$2$},
  clip=true]
\addplot[popcurve, domain=-4:0.86, samples=120]{(x^2+x)/(x-1)};
\addplot[popcurve, domain=1.14:8, samples=120]{(x^2+x)/(x-1)};
\addplot[poporange, dashed, thick, domain=-4:8]{x+2};
\addplot[popgreen, dashed, thick] coordinates {(1,-8) (1,16)};
\node[poporange, anchor=north west] at (axis cs:5.5,7.6) {$y=x+2$};
\node[popgreen, anchor=west] at (axis cs:1.1,-6) {$x=1$};
\node[popblue, anchor=south east] at (axis cs:-1.5,0.4) {$\mathcal{C}_f$};
\end{axis}
\end{tikzpicture}
\end{center}
\legende{La courbe de $f(x)=\dfrac{x^2+x}{x-1}$ : asymptote verticale $x=1$ et asymptote oblique $y=x+2$. Comme $f(x)-(x+2)=\dfrac{2}{x-1}$, la courbe est au-dessus de l'asymptote pour $x>1$ et en dessous pour $x<1$.}
\end{popfigure}

\courssub{Directions asymptotiques et branches paraboliques}

Que faire quand le plan d'étude en deux étapes échoue ? On distingue deux situations.

\begin{definition}[Branche parabolique]
Soit $f$ définie au voisinage de $+\infty$ (ou $-\infty$).
\begin{itemize}
\item Si $\lim \dfrac{f(x)}{x} = a$ (réel) mais $\lim\big[f(x)-ax\big] = \pm\infty$, on dit que $\mathcal{C}_f$ admet une \cle{branche parabolique de direction} la droite $y = ax$.
\item Si $\lim \dfrac{f(x)}{x} = \pm\infty$, on dit que $\mathcal{C}_f$ admet une \cle{branche parabolique de direction} l'axe $(Oy)$.
\end{itemize}
\end{definition}

\begin{exemplebox}[Deux branches paraboliques de référence]
\textbf{a)} $f(x)=\sqrt{x}$ en $+\infty$ : $\dfrac{f(x)}{x}=\dfrac{1}{\sqrt{x}}\to 0$ et $f(x)\to+\infty$ : branche parabolique de direction l'axe $(Ox)$ (la pente tend vers $0$ mais la courbe monte sans fin).

\textbf{b)} $f(x)=x^2$ en $+\infty$ : $\dfrac{f(x)}{x}=x\to+\infty$ : branche parabolique de direction $(Oy)$ (la courbe monte de plus en plus verticalement).
\end{exemplebox}

\begin{popfigure}
\begin{center}
\begin{tikzpicture}
\begin{axis}[
  width=0.85\linewidth, height=6.5cm,
  axis lines=middle, xlabel={$x$}, ylabel={$y$},
  xmin=-0.5, xmax=9, ymin=-1, ymax=8,
  clip=true]
\addplot[popcurve, domain=0:9, samples=150]{sqrt(max(0,x))};
\addplot[popcurve, domain=0:2.6, samples=150]{x^2};
\addplot[poporange, dashed, thick, domain=0:8]{0.5*x+1.5};
\node[popblue, anchor=south] at (axis cs:6.5,2.55) {$y=\sqrt{x}$};
\node[popblue, anchor=west] at (axis cs:2.15,5.6) {$y=x^2$};
\node[poporange, anchor=north west, rotate=26] at (axis cs:5.2,4.6) {asymptote oblique};
\end{axis}
\end{tikzpicture}
\end{center}
\legende{Trois comportements à l'infini : $y=\sqrt{x}$ s'écrase lentement (branche parabolique de direction $(Ox)$), $y=x^2$ se redresse (direction $(Oy)$), tandis qu'une courbe à asymptote oblique se plaque contre sa droite.}
\end{popfigure}

\courssub{Tableau récapitulatif des six cas de branches infinies}

\begin{aretenir}[Les six cas à connaître par cœur]
\begin{center}
\begin{tabularx}{\linewidth}{|l|X|X|}
\hline
\rowcolor{popblueL} \textbf{Hypothèse sur la limite} & \textbf{Conclusion} & \textbf{Nature} \\
\hline
$\lim_{x\to a}f(x)=\pm\infty$ & Droite $x=a$ & Asymptote verticale \\
\hline
$\lim_{x\to\pm\infty}f(x)=b$ & Droite $y=b$ & Asymptote horizontale \\
\hline
$\lim\big[f(x)-(ax+b)\big]=0$, $a\neq 0$ & Droite $y=ax+b$ & Asymptote oblique \\
\hline
$\lim\dfrac{f(x)}{x}=0$ et $\lim f(x)=\pm\infty$ & Direction $y=0$ & Branche parabolique de direction $(Ox)$ \\
\hline
$\lim\dfrac{f(x)}{x}=a\neq 0$ et $\lim\big[f(x)-ax\big]=\pm\infty$ & Direction $y=ax$ & Branche parabolique de direction $y=ax$ \\
\hline
$\lim\dfrac{f(x)}{x}=\pm\infty$ & Direction $(Oy)$ & Branche parabolique de direction $(Oy)$ \\
\hline
\end{tabularx}
\end{center}
\end{aretenir}

\begin{exercice}[Pour t'entraîner]
Soit $g(x)=\dfrac{x^2+3x-1}{x+2}$ et $h(x)=\sqrt{x^2-4}+x$.
\begin{enumerate}
\item Montrer que la droite $y=x+1$ est asymptote oblique à $\mathcal{C}_g$ et préciser la position relative.
\item Étudier les branches infinies de $h$ en $+\infty$ et en $-\infty$ (on montrera en particulier que $y=2x$ est asymptote en $+\infty$ et $y=0$ asymptote en $-\infty$).
\end{enumerate}
\end{exercice}

\begin{corrige}[Exercice : branches infinies]
\textbf{1.} $\dfrac{g(x)}{x}=\dfrac{x^2+3x-1}{x^2+2x}\to 1$, donc $a=1$. Puis $g(x)-x=\dfrac{x^2+3x-1-x(x+2)}{x+2}=\dfrac{x-1}{x+2}\to 1$, donc $b=1$ : $y=x+1$ est asymptote oblique en $\pm\infty$. Position : $g(x)-(x+1)=\dfrac{x-1}{x+2}-1=\dfrac{-3}{x+2}$. Pour $x>-2$, la courbe est en dessous ; pour $x<-2$, elle est au-dessus.

\textbf{2.} En $+\infty$ : $\dfrac{h(x)}{x}=\sqrt{1-\frac{4}{x^2}}+1\to 2$, et $h(x)-2x=\sqrt{x^2-4}-x=\dfrac{-4}{\sqrt{x^2-4}+x}\to 0$ (quantité conjuguée) : $y=2x$ est asymptote oblique, et comme $h(x)-2x<0$, la courbe est en dessous. En $-\infty$ : par quantité conjuguée, $h(x)=\dfrac{(\sqrt{x^2-4}+x)(\sqrt{x^2-4}-x)}{\sqrt{x^2-4}-x}=\dfrac{-4}{\sqrt{x^2-4}-x}\to 0$ car le dénominateur tend vers $+\infty$ (on a $\sqrt{x^2-4}\to+\infty$ et $-x\to+\infty$) : $y=0$ est asymptote horizontale en $-\infty$.
\end{corrige}

\begin{savaistu}[Les asymptotes dans la vie courante]
Les asymptotes ne sont pas qu'un jeu de mathématiciens ! En économie, le coût unitaire de production d'une entreprise de Douala diminue quand la quantité produite augmente, mais sans jamais descendre sous un certain seuil : il admet une asymptote horizontale. En physique, la vitesse d'un parachutiste tend vers une vitesse limite sans jamais la dépasser : là encore, une asymptote horizontale. Et la courbe de charge d'un réseau mobile en heure de pointe peut présenter des branches quasi verticales : les ingénieurs de MTN ou Orange doivent anticiper ces « explosions » de trafic.
\end{savaistu}

Tu disposes maintenant de tous les outils : limites usuelles, quantité conjuguée, recherche d'asymptotes et étude des positions relatives. Il ne reste plus qu'à assembler ces pièces dans la dernière section, consacrée à l'\emph{étude complète de fonctions} et au tracé de leurs courbes.

\coursec{Étude complète de fonctions et tracé de courbes}

Dans la section précédente, tu as appris à traquer les \cle{branches infinies} : asymptotes verticales, horizontales, obliques, directions asymptotiques. Ces outils, ajoutés au calcul de limites, à la dérivation, au théorème des valeurs intermédiaires et à l'inégalité des accroissements finis, forment désormais une boîte à outils complète. Il est temps de tout assembler : c'est l'objet de cette dernière section, qui est aussi \emph{la} compétence reine du problème de Baccalauréat. Une \cle{étude complète de fonction}, c'est l'enquête méthodique qui mène d'une formule $f(x)$ à une courbe fidèle, cohérente et exploitable. Suivons le plan, étape par étape.

\courssub{Le plan type d'une étude complète}

\begin{methode}[Les 8 étapes du plan d'étude — fiche à mémoriser pour le Bac]
\begin{enumerate}
\item \textbf{Ensemble de définition} $D_f$ : repérer les valeurs interdites (dénominateur nul, quantité sous une racine carrée négative, argument de logarithme non positif, etc.) et écrire $D_f$ comme réunion d'intervalles.
\item \textbf{Restriction du domaine d'étude} : tester la parité ($f(-x) = f(x)$ : fonction paire ; $f(-x) = -f(x)$ : impaire) et la périodicité ($f(x+T) = f(x)$). On étudie alors sur un domaine réduit.
\item \textbf{Limites aux bornes} de $D_f$ : calculer et justifier chaque limite (bornes finies ouvertes, $+\infty$, $-\infty$).
\item \textbf{Continuité et dérivabilité} : justifier sur chaque intervalle de $D_f$ (fonction rationnelle, composée, etc.).
\item \textbf{Dérivée et signe de $f'$} : calculer $f'(x)$, factoriser, dresser le tableau de signes.
\item \textbf{Tableau de variations complet} : y reporter le signe de $f'$, les limites calculées à l'étape 3 et les valeurs remarquables (extremums).
\item \textbf{Branches infinies et asymptotes} : déduire des limites les asymptotes verticales, horizontales, obliques ; étudier la position relative courbe/asymptote par le signe de $f(x) - (ax+b)$.
\item \textbf{Points remarquables et tracé} : intersections avec les axes, tangentes horizontales, demi-tangentes (points anguleux), puis tracé soigné cohérent avec le tableau.
\end{enumerate}
\end{methode}

\begin{methode}[Exploiter parité et périodicité pour restreindre le domaine]
\begin{enumerate}
\item Vérifier d'abord que $D_f$ est \textbf{symétrique par rapport à 0} (sinon, ni parité ni imparité possibles).
\item Si $f$ est \textbf{paire} : étudier sur $D_f \cap [0\,; +\infty[$, puis compléter la courbe par \textbf{symétrie d'axe $(Oy)$}.
\item Si $f$ est \textbf{impaire} : étudier sur $D_f \cap [0\,; +\infty[$, puis compléter par \textbf{symétrie de centre $O$}.
\item Si $f$ est \textbf{$T$-périodique} : étudier sur un intervalle d'amplitude $T$ (par exemple $[0\,;T]$ ou $[-T/2\,;T/2]$), puis compléter par \textbf{translations de vecteur $kT\vec{\imath}$}, $k \in \mathbb{Z}$.
\item Parité $+$ périodicité : on peut souvent réduire l'étude à $[0\,; T/2]$.
\end{enumerate}
\end{methode}

\begin{attention}[Les trois fautes qui coûtent cher]
\begin{itemize}
\item \textbf{Oublier une valeur interdite} : si $f(x) = \dfrac{x^2}{x-1}$, écrire $D_f = \mathbb{R}$ est une faute éliminatoire dès la première ligne : $D_f = \mathbb{R} \setminus \{1\}$.
\item \textbf{Tracer la courbe qui traverse une asymptote verticale} : la courbe ne touche jamais une asymptote verticale ; elle « file » vers $+\infty$ ou $-\infty$ de part et d'autre.
\item \textbf{Incohérence tracé / tableau} : si le tableau annonce une fonction croissante de $-\infty$ vers $3$, la courbe ne peut pas redescendre ni dépasser $3$. Le correcteur vérifie toujours cette cohérence.
\end{itemize}
\end{attention}

\begin{attention}[Règle d'or de rédaction au Bac]
Toute limite qui apparaît dans le tableau de variations doit avoir été \textbf{calculée et justifiée dans une question précédente}. Un tableau contenant une limite non démontrée est sanctionné, même si la valeur est exacte. Rédige donc les calculs de limites \emph{avant} de dresser le tableau.
\end{attention}

\courssub{Étude complète d'une fonction rationnelle avec asymptote oblique}

\begin{exemplebox}[Étude guidée de $f(x) = \dfrac{x^2 - 2x + 2}{x - 1}$]
\textbf{Étape 1 — Ensemble de définition.} Le dénominateur s'annule en $x = 1$ : $D_f = \mathbb{R} \setminus \{1\} = ]-\infty\,; 1[ \cup ]1\,; +\infty[$.

\textbf{Étape 2 — Parité.} $D_f$ n'est pas symétrique par rapport à $0$ : pas de réduction possible.

\textbf{Étape 3 — Limites.} En $1$ : le numérateur vaut $1 - 2 + 2 = 1 > 0$. Donc
\[ \lim_{\substack{x \to 1 \\ x < 1}} f(x) = -\infty \quad \text{et} \quad \lim_{\substack{x \to 1 \\ x > 1}} f(x) = +\infty. \]
En l'infini, $f(x) \sim \dfrac{x^2}{x} = x$, donc $\lim\limits_{x \to +\infty} f(x) = +\infty$ et $\lim\limits_{x \to -\infty} f(x) = -\infty$.

\textbf{Étape 4 — Dérivabilité.} $f$ est une fonction rationnelle : elle est continue et dérivable sur chaque intervalle de $D_f$.

\textbf{Étape 5 — Dérivée.} Avec $u = x^2 - 2x + 2$ et $v = x - 1$ :
\[ f'(x) = \frac{(2x-2)(x-1) - (x^2 - 2x + 2)}{(x-1)^2} = \frac{x^2 - 2x}{(x-1)^2} = \frac{x(x-2)}{(x-1)^2}. \]
Le dénominateur est strictement positif sur $D_f$ : le signe de $f'$ est celui de $x(x-2)$. Donc $f' > 0$ sur $]-\infty\,;0[$ et $]2\,;+\infty[$, $f' < 0$ sur $]0\,;1[$ et $]1\,;2[$. Extremums : $f(0) = -2$ (maximum local) et $f(2) = 2$ (minimum local).

\textbf{Étape 6 — Tableau de variations.}
\begin{center}
\begin{tabular}{|c|ccccccccc|}
\hline
$x$ & $-\infty$ & & $0$ & & $1$ & & $2$ & & $+\infty$ \\
\hline
$f'(x)$ & & $+$ & $0$ & $-$ & $\|$ & $-$ & $0$ & $+$ & \\
\hline
$f(x)$ & $-\infty$ & $\nearrow$ & $-2$ & $\searrow$ & $-\infty \ \| \ +\infty$ & $\searrow$ & $2$ & $\nearrow$ & $+\infty$ \\
\hline
\end{tabular}
\end{center}

\textbf{Étape 7 — Branches infinies.} La droite $x = 1$ est \textbf{asymptote verticale}. Pour l'infini, effectuons la division euclidienne : $x^2 - 2x + 2 = (x-1)(x-1) + 1$, d'où
\[ f(x) = x - 1 + \frac{1}{x-1}. \]
Comme $\lim\limits_{x \to \pm\infty} \bigl(f(x) - (x-1)\bigr) = \lim\limits_{x \to \pm\infty} \dfrac{1}{x-1} = 0$, la droite d'équation $y = x - 1$ est \textbf{asymptote oblique}. Position relative : $f(x) - (x-1) = \dfrac{1}{x-1}$, positif si $x > 1$ (courbe \emph{au-dessus}), négatif si $x < 1$ (courbe \emph{en dessous}).

\textbf{Étape 8 — Tracé.} La courbe coupe $(Oy)$ en $(0\,; -2)$, admet des tangentes horizontales en $x = 0$ et $x = 2$, et ne coupe pas $(Ox)$ (le numérateur $x^2 - 2x + 2$ a pour discriminant $-4 < 0$).
\end{exemplebox}

\begin{popfigure}
\begin{tikzpicture}
\begin{axis}[width=0.95\linewidth, height=8.5cm, axis lines=middle, xlabel=$x$, ylabel=$y$, xmin=-4.5, xmax=6.5, ymin=-7, ymax=9, xtick={-4,-2,1,2,4,6}, ytick={-6,-4,-2,2,4,6,8}, legend style={at={(0.02,0.98)}, anchor=north west, font=\small}]
\addplot[popcurve, domain=-4.5:0.88, samples=200]{(x^2-2*x+2)/(x-1)};
\addplot[popcurve, domain=1.12:6.3, samples=200]{(x^2-2*x+2)/(x-1)};
\addplot[dashed, poporange, thick] coordinates {(1,-7) (1,9)};
\addplot[dashed, popgreen, thick, domain=-4.5:6.5]{x-1};
\addplot[only marks, mark=*, poppink, mark size=2pt] coordinates {(0,-2) (2,2)};
\legend{$\mathcal{C}_f$, , $x=1$, $y=x-1$, extremums}
\end{axis}
\end{tikzpicture}
\legende{Courbe de $f(x) = \dfrac{x^2-2x+2}{x-1}$ : asymptote verticale $x=1$, asymptote oblique $y = x-1$, maximum local $-2$ en $0$, minimum local $2$ en $2$.}
\end{popfigure}

\courssub{Étude d'une fonction irrationnelle : demi-tangentes verticales}

\begin{exoresolu}[Étude de $f(x) = \sqrt{x^2 - 4}$]
Étudier complètement $f$ et tracer sa courbe.
\tcblower
\textbf{Ensemble de définition.} Il faut $x^2 - 4 \geq 0$, soit $|x| \geq 2$ : $D_f = ]-\infty\,; -2] \cup [2\,; +\infty[$.

\textbf{Parité.} $f(-x) = \sqrt{(-x)^2 - 4} = f(x)$ : $f$ est \textbf{paire}. On étudie sur $[2\,; +\infty[$ et on complète par symétrie d'axe $(Oy)$.

\textbf{Limites.} $\lim\limits_{x \to +\infty} \sqrt{x^2 - 4} = +\infty$. En $2$ : $f(2) = 0$ (fonction continue à droite).

\textbf{Dérivabilité.} $f$ est dérivable là où $x^2 - 4 > 0$, c'est-à-dire sur $]2\,; +\infty[$. \textbf{Attention} : en $x = 2$, la quantité sous la racine s'annule, $f$ n'est \emph{pas} dérivable en $2$ au sens usuel. Étudions le taux d'accroissement :
\[ \frac{f(x) - f(2)}{x - 2} = \frac{\sqrt{x^2-4}}{x-2} = \frac{\sqrt{(x-2)(x+2)}}{x-2} = \sqrt{\frac{x+2}{x-2}} \xrightarrow[x \to 2^+]{} +\infty. \]
La courbe admet donc en $(2\,; 0)$ une \cle{demi-tangente verticale} dirigée vers le haut. Par parité, demi-tangente verticale aussi en $(-2\,; 0)$.

\textbf{Dérivée et variations.} Pour $x > 2$ : $f'(x) = \dfrac{2x}{2\sqrt{x^2-4}} = \dfrac{x}{\sqrt{x^2-4}} > 0$. Donc $f$ est strictement croissante sur $[2\,; +\infty[$, de $0$ vers $+\infty$.

\textbf{Branches infinies.} Pour $x > 0$, $\sqrt{x^2 - 4} = x\sqrt{1 - \dfrac{4}{x^2}}$, donc
\[ f(x) - x = x\left(\sqrt{1 - \tfrac{4}{x^2}} - 1\right) = \frac{-4x}{x^2\left(\sqrt{1 - \frac{4}{x^2}} + 1\right)} \xrightarrow[x \to +\infty]{} 0. \]
La droite $y = x$ est asymptote en $+\infty$, et par parité $y = -x$ est asymptote en $-\infty$ : la courbe a pour branche asymptotique la droite $y = |x|$. Comme $f(x) - x < 0$ pour $x > 2$, la courbe est \emph{en dessous} de son asymptote.

\textbf{Tracé.} Voir la figure ci-dessous : la courbe est la partie supérieure de l'hyperbole équilatère $x^2 - y^2 = 4$.
\end{exoresolu}

\begin{popfigure}
\begin{tikzpicture}
\begin{axis}[width=0.95\linewidth, height=7.5cm, axis lines=middle, xlabel=$x$, ylabel=$y$, xmin=-7, xmax=7, ymin=-1, ymax=6.5, xtick={-6,-4,-2,2,4,6}, ytick={2,4,6}, legend style={at={(0.02,0.98)}, anchor=north west, font=\small}]
\addplot[popcurve, domain=2:6.8, samples=200]{sqrt(max(0,x^2-4))};
\addplot[popcurve, domain=-6.8:-2, samples=200]{sqrt(max(0,x^2-4))};
\addplot[dashed, poporange, thick, domain=-6.8:6.8]{abs(x)};
\addplot[only marks, mark=*, poppink, mark size=2pt] coordinates {(2,0) (-2,0)};
\legend{$\mathcal{C}_f$, , $y = |x|$, points à demi-tangente verticale}
\end{axis}
\end{tikzpicture}
\legende{Courbe de $f(x) = \sqrt{x^2 - 4}$ : demi-tangentes verticales en $(\pm 2\,; 0)$, asymptotes $y = x$ en $+\infty$ et $y = -x$ en $-\infty$.}
\end{popfigure}

\begin{attention}[Demi-tangente verticale : le piège des fonctions racines]
Pour $g(x) = \sqrt{u(x)}$, la dérivée $g'(x) = \dfrac{u'(x)}{2\sqrt{u(x)}}$ n'existe que si $u(x) > 0$. Aux bornes où $u$ s'annule, ne conclus jamais trop vite : calcule la limite du \textbf{taux d'accroissement}. Si elle est infinie, il y a une demi-tangente verticale (la courbe « démarre » à la verticale), et non une tangente horizontale.
\end{attention}

\courssub{Étude d'une fonction trigonométrique}

\begin{exoresolu}[{Étude de $f(x) = x + \sin x$ sur $[-2\pi\,; 2\pi]$}]
Étudier $f$ et tracer sa courbe.
\tcblower
\textbf{Domaine et symétrie.} $f$ est définie sur $\mathbb{R}$. De plus $f(-x) = -x - \sin x = -f(x)$ : $f$ est \textbf{impaire}. On étudie sur $[0\,; 2\pi]$ et on complète par symétrie de centre $O$.

\textbf{Dérivée.} $f'(x) = 1 + \cos x \geq 0$ pour tout $x$, et $f'(x) = 0$ uniquement lorsque $\cos x = -1$, c'est-à-dire en $x = \pi$ (sur $[0\,;2\pi]$). La dérivée ne s'annule qu'en un point isolé \textbf{sans changer de signe} : $f$ est strictement croissante sur $\mathbb{R}$, avec une \textbf{tangente horizontale} au point $(\pi\,; \pi)$ — c'est un point d'inflexion à tangente horizontale.

\textbf{Valeurs et limites.} $f(0) = 0$, $f(\pi) = \pi \approx \num{3,14}$, $f(2\pi) = 2\pi \approx \num{6,28}$. Sur $\mathbb{R}$ : $\lim\limits_{x \to +\infty} f(x) = +\infty$ car $f(x) \geq x - 1$.

\textbf{Branches infinies.} $\dfrac{f(x)}{x} = 1 + \dfrac{\sin x}{x} \to 1$ et $f(x) - x = \sin x$ n'a pas de limite : la droite $y = x$ est une \textbf{direction asymptotique} (la courbe ondule autour de $y = x$ sans s'en rapprocher), mais pas une asymptote au sens strict.

\textbf{Tracé.} La courbe suit globalement la droite $y = x$, avec des tangentes horizontales en tous les points $(\pi + 2k\pi\,; \pi + 2k\pi)$, $k \in \mathbb{Z}$.
\end{exoresolu}

\begin{popfigure}
\begin{tikzpicture}
\begin{axis}[width=0.95\linewidth, height=8cm, axis lines=middle, xlabel=$x$, ylabel=$y$, xmin=-7, xmax=7, ymin=-7.5, ymax=7.5, xtick={-6.2832,-3.1416,3.1416,6.2832}, xticklabels={$-2\pi$,$-\pi$,$\pi$,$2\pi$}, ytick={-6.2832,-3.1416,3.1416,6.2832}, yticklabels={$-2\pi$,$-\pi$,$\pi$,$2\pi$}, legend style={at={(0.02,0.98)}, anchor=north west, font=\small}]
\addplot[popcurve, domain=-6.2832:6.2832, samples=300]{x + sin(deg(x))};
\addplot[dashed, poporange, thick, domain=-6.8:6.8]{x};
\addplot[only marks, mark=*, poppink, mark size=2pt] coordinates {(-3.1416,-3.1416) (3.1416,3.1416)};
\legend{$\mathcal{C}_f$, $y = x$, tangentes horizontales}
\end{axis}
\end{tikzpicture}
\legende{Courbe de $f(x) = x + \sin x$ sur $[-2\pi\,; 2\pi]$ : croissante, elle ondule autour de la direction asymptotique $y = x$ avec des tangentes horizontales en $\pm\pi$.}
\end{popfigure}

\courssub{Exploitation d'une étude : équations et fonction réciproque}

Une étude complète n'est pas une fin en soi : elle sert à \emph{répondre à des questions}. Les deux exploitations classiques au Bac sont les suivantes.

\textbf{Résolution d'équations.} Le tableau de variations, combiné au \cle{théorème des valeurs intermédiaires}, permet de dénombrer et d'encadrer les solutions de $f(x) = c$.

\begin{exemplebox}[Application du TVI avec la fonction rationnelle étudiée]
Reprenons $f(x) = \dfrac{x^2 - 2x + 2}{x - 1}$. Montrons que l'équation $f(x) = 3$ admet une unique solution $\alpha$ dans $]1\,; +\infty[$. Sur $]1\,; 2]$, $f$ est continue et strictement décroissante de $+\infty$ vers $2$ ; or $3 \in ]2\,; +\infty[$, donc le TVI garantit une unique solution dans $]1\,; 2]$. Sur $[2\,; +\infty[$, $f$ est strictement croissante de $2$ vers $+\infty$ : le TVI donne une seconde solution. Bilan : exactement \textbf{deux} solutions sur $]1\,;+\infty[$. Comme $f(\num{1,5}) = \dfrac{2,25 - 3 + 2}{0,5} = \num{2,5} < 3$ et $f(\num{1,3}) \approx \num{3,63} > 3$, on encadre la première : $\alpha \in ]\num{1,3}\,; \num{1,5}[$.
\end{exemplebox}

\textbf{Courbe de la réciproque.} Si $f$ est continue et strictement monotone sur un intervalle $I$, elle réalise une bijection de $I$ sur $J = f(I)$, et la courbe de $f^{-1}$ s'obtient par \textbf{symétrie d'axe $y = x$}. Par exemple, $f(x) = x + \sin x$ est strictement croissante sur $\mathbb{R}$ avec $f(\mathbb{R}) = \mathbb{R}$ : elle admet une réciproque $f^{-1} : \mathbb{R} \to \mathbb{R}$, dérivable partout sauf aux points images des tangentes horizontales, où $(f^{-1})'$ « explose » (demi-tangentes verticales en $\pi + 2k\pi$).

\begin{methode}[Tracer la courbe de $f^{-1}$ à partir de celle de $f$]
\begin{enumerate}
\item Justifier que $f$ est continue et strictement monotone sur $I$ (tableau de variations).
\item Déterminer $J = f(I)$ grâce aux limites du tableau.
\item Tracer la droite $y = x$ en pointillés.
\item Reporter les points remarquables : $(a\,; b) \in \mathcal{C}_f \iff (b\,; a) \in \mathcal{C}_{f^{-1}}$.
\item Les asymptotes se transposent : une asymptote horizontale $y = c$ pour $f$ devient une asymptote verticale $x = c$ pour $f^{-1}$.
\end{enumerate}
\end{methode}

\begin{savaistu}[La répétition générale du problème de Bac]
Aux Baccalauréats C et E, le « problème » de l'épreuve de mathématiques se termine presque toujours par une étude complète de fonction : on y mobilise successivement le calcul de limites, le TVI pour localiser une solution $\alpha$, l'inégalité des accroissements finis pour encadrer une valeur, la recherche d'asymptotes, puis le tracé final de la courbe — parfois accompagné de celui de la réciproque. Cette section est donc ta répétition générale : chaque étude que tu rédiges à l'entraînement doit suivre les 8 étapes de la fiche méthode, sans exception. Le jour de l'examen, tu ne réfléchiras plus au plan : tu l'appliqueras, et tu gagneras un temps précieux pour soigner la rédaction et le graphique.
\end{savaistu}

\begin{exercice}[Pour t'entraîner]
Soit $g(x) = \dfrac{x}{\sqrt{x^2 + 1}}$.
\begin{enumerate}
\item Déterminer $D_g$, montrer que $g$ est impaire et restreindre le domaine d'étude.
\item Calculer les limites aux bornes et en déduire les asymptotes.
\item Montrer que $g'(x) = \dfrac{1}{(x^2+1)\sqrt{x^2+1}}$ et dresser le tableau de variations.
\item Montrer que $g$ réalise une bijection de $\mathbb{R}$ sur $]-1\,; 1[$ et préciser les asymptotes de la courbe de $g^{-1}$.
\item Tracer $\mathcal{C}_g$ et $\mathcal{C}_{g^{-1}}$ dans un même repère.
\end{enumerate}
\end{exercice}

\begin{exercice}[Fonction trigonométrique]
Soit $h(x) = \dfrac{\sin x}{2 + \cos x}$.
\begin{enumerate}
\item Montrer que $h$ est définie sur $\mathbb{R}$, impaire et $2\pi$-périodique. Justifier qu'on peut étudier $h$ sur $[0\,; \pi]$.
\item Calculer $h'(x)$ et montrer que $h'(x) = \dfrac{2\cos x + 1}{(2 + \cos x)^2}$.
\item Dresser le tableau de variations sur $[0\,; \pi]$ (on utilisera $\cos x = -\tfrac12 \iff x = \tfrac{2\pi}{3}$).
\item Tracer la courbe sur $[-2\pi\,; 2\pi]$ en complétant par symétrie puis par translations.
\end{enumerate}
\end{exercice}

\newpage

\coursec{Activité d'intégration}

\coursec{Leçon M03 : Étude de fonctions numériques — Activité d'intégration : Le forage de Bafoussam}

\courssub{Objectifs de l'activité}

À l'issue de cette activité d'intégration, tu seras capable de :

\begin{itemize}
\item mener l'\cle{étude complète} d'une fonction rationnelle (ensemble de définition, dérivée, variations, limites, \cle{branches infinies}) et interpréter concrètement une \cle{asymptote oblique} ;
\item démontrer l'\cle{existence} et l'\cle{unicité} d'une solution de l'équation $C(x)=c$ grâce au \cle{théorème des valeurs intermédiaires} (TVI), puis l'encadrer par \cle{dichotomie} ;
\item justifier qu'une fonction réalise une \cle{bijection} et exploiter sa \cle{fonction réciproque}, notamment la formule de la \cle{dérivée de la réciproque} sans expliciter $C^{-1}$ ;
\item utiliser l'\cle{inégalité des accroissements finis} (IAF) pour majorer un accroissement et valider ou invalider une affirmation pratique ;
\item tracer une courbe et déduire l'allure de la courbe de la fonction réciproque par \cle{symétrie} par rapport à la droite d'équation $y=x$.
\end{itemize}

\courssub{Situation}

La commune de Bafoussam, chef-lieu de la région de l'Ouest, connaît une croissance démographique rapide. Pour améliorer l'accès à l'eau potable dans les quartiers périphériques (Tougang, Banengo, Djeleng), la mairie lance un programme de \textbf{forages} (puits profonds équipés de pompes). Le bureau d'études retenu a établi, à partir de forages comparables réalisés dans la région, un modèle de coût.

On désigne par $x$ la \textbf{profondeur} du forage, exprimée en \textbf{dizaines de mètres} (ainsi $x=3$ correspond à \SI{30}{m}), et par $C(x)$ le \textbf{coût total} du forage, exprimé en \textbf{centaines de milliers de FCFA} (ainsi $C(x)=15$ correspond à \num{1500000}~FCFA). Le modèle retenu est :
\[ C(x)=\dfrac{x^{2}+2x+4}{x+1},\qquad x\geq 1. \]

La mairie dispose, pour le premier forage pilote du quartier Tougang, d'une enveloppe de \textbf{\num{1500000}~FCFA}. Un sous-traitant local affirme par ailleurs que «~faire passer un forage de \SI{30}{m} à \SI{40}{m} coûte au plus $1{,}2$ fois le coût d'une dizaine de mètres supplémentaires, c'est-à-dire au plus \num{120000}~FCFA de surcoût~». La mairie te missionne, toi et ton groupe, pour éclairer sa décision à l'aide des outils mathématiques de Terminale.

\courssub{Consignes}

Travail en groupes de quatre. Chaque groupe rédigera un rapport et préparera une affiche comportant le tracé des courbes.

\textbf{Partie A : étude du modèle de coût}
\begin{enumerate}
\item Justifier que $C$ est définie et dérivable sur $[1\,;\,+\infty[$. Montrer que pour tout $x\geq 1$ :
\[ C(x)=x+1+\dfrac{3}{x+1}. \]
\item Calculer $C'(x)$, étudier son signe sur $[1\,;\,+\infty[$ et dresser le tableau de variations complet de $C$ (avec les limites aux bornes).
\item Déterminer les limites de $C$ en $+\infty$, puis montrer que la droite $(\Delta)$ d'équation $y=x+1$ est \cle{asymptote oblique} à la courbe $(\mathcal{C})$ de $C$. Préciser la position relative de $(\mathcal{C})$ et $(\Delta)$. Interpréter cette asymptote en termes de «~coût de structure~» pour un forage très profond.
\end{enumerate}

\textbf{Partie B : le budget de la mairie}
\begin{enumerate}
\setcounter{enumi}{3}
\item Démontrer, à l'aide du théorème des valeurs intermédiaires, que l'équation $C(x)=15$ admet une \textbf{unique} solution $\alpha$ dans $[1\,;\,+\infty[$.
\item Par dichotomie (ou balayage), donner un encadrement de $\alpha$ à $10^{-1}$ près. En déduire la profondeur maximale, en mètres, que la mairie peut financer avec son enveloppe.
\end{enumerate}

\textbf{Partie C : la fonction réciproque, outil du planificateur}
\begin{enumerate}
\setcounter{enumi}{5}
\item Justifier que $C$ réalise une bijection de $[1\,;\,+\infty[$ sur un intervalle $J$ à préciser. Que représente concrètement $C^{-1}(y)$ pour la mairie ?
\item Sans expliciter $C^{-1}$, calculer $\left(C^{-1}\right)'(15)$. Interpréter ce nombre : de combien varie approximativement la profondeur finançable lorsque le budget augmente de \num{100000}~FCFA au voisinage de \num{1500000}~FCFA ?
\end{enumerate}

\textbf{Partie D : l'affirmation du sous-traitant}
\begin{enumerate}
\setcounter{enumi}{7}
\item Montrer que $C'$ est croissante sur $[3\,;\,4]$ et en déduire que pour tout $x\in[3\,;\,4]$, $C'(x)\leq \dfrac{22}{25}$.
\item À l'aide de l'inégalité des accroissements finis, majorer $C(4)-C(3)$. L'affirmation du sous-traitant est-elle fondée ? Comparer avec la valeur exacte de $C(4)-C(3)$.
\end{enumerate}

\textbf{Partie E : restitution graphique}
\begin{enumerate}
\setcounter{enumi}{9}
\item Tracer soigneusement $(\mathcal{C})$, l'asymptote $(\Delta)$ et la droite d'équation $y=x$ dans un repère orthonormé (unité : \SI{0,5}{cm} par unité). Placer le point d'abscisse $\alpha$.
\item En déduire, par symétrie, l'allure de la courbe de $C^{-1}$. Présenter l'affiche au conseil de classe.
\end{enumerate}

\courssub{Ressources à mobiliser}

\begin{aretenir}[Boîte à outils de l'activité]
\begin{itemize}
\item \textbf{TVI (bijection)} : si $f$ est \cle{continue} et \cle{strictement monotone} sur un intervalle $I$, alors $f$ réalise une bijection de $I$ sur $f(I)$ ; l'équation $f(x)=c$ admet une unique solution dans $I$ dès que $c\in f(I)$.
\item \textbf{Dérivée de la réciproque} : si $f$ est une bijection dérivable et $f'(x_0)\neq 0$, alors $\left(f^{-1}\right)'(y_0)=\dfrac{1}{f'(x_0)}$ avec $y_0=f(x_0)$.
\item \textbf{IAF} : si $|f'(x)|\leq M$ sur $[a\,;\,b]$, alors $|f(b)-f(a)|\leq M\,|b-a|$.
\item \textbf{Asymptote oblique} : si $f(x)-(ax+b)\to 0$ en $+\infty$, la droite $y=ax+b$ est asymptote à la courbe.
\item \textbf{Symétrie} : les courbes de $f$ et $f^{-1}$ sont symétriques par rapport à la droite $y=x$.
\end{itemize}
\end{aretenir}

\courssub{Démarche et corrigé commenté}

\begin{exoresolu}[Forage de Bafoussam : corrigé complet]
On considère $C(x)=\dfrac{x^{2}+2x+4}{x+1}$ pour $x\geq 1$. Questions : étude de $C$, résolution de $C(x)=15$, étude de $C^{-1}$, majoration de $C(4)-C(3)$, tracés.
\tcblower
\textbf{Partie A.} \textbf{1.} Le dénominateur $x+1$ ne s'annule pas sur $[1\,;\,+\infty[$ (car $x+1\geq 2$), donc $C$, fonction rationnelle, y est définie, continue et dérivable. En effectuant la division euclidienne du numérateur par le dénominateur :
\[ x^{2}+2x+4 = (x+1)(x+1) + 3, \]
d'où, pour tout $x\geq 1$,
\[ C(x)=x+1+\dfrac{3}{x+1}. \]

\textbf{2.} On dérive cette forme simplifiée :
\[ C'(x)=1-\dfrac{3}{(x+1)^{2}}=\dfrac{(x+1)^{2}-3}{(x+1)^{2}}. \]
Pour $x\geq 1$, on a $x+1\geq 2$, donc $(x+1)^{2}\geq 4 > 3$. Le numérateur est strictement positif, le dénominateur aussi, donc $C'(x)>0$ sur $[1\,;\,+\infty[$. La fonction $C$ est \textbf{strictement croissante} sur $[1\,;\,+\infty[$.
On calcule $C(1)=\dfrac{1+2+4}{2}=\dfrac{7}{2}=3{,}5$ et $\lim_{x\to+\infty}C(x)=+\infty$ (car $x+1\to+\infty$ et $\frac{3}{x+1}\to 0$).

\begin{center}
\begin{tabular}{|c|cccccc|}\hline
$x$ & $1$ & & & & & $+\infty$ \\ \hline
$C'(x)$ & & $+$ & & & & \\ \hline
$C(x)$ & $\dfrac{7}{2}$ & \nearrow & & & & $+\infty$ \\ \hline
\end{tabular}
\end{center}

\textbf{3.} On a $C(x)-(x+1)=\dfrac{3}{x+1}$. Comme $\lim_{x\to+\infty}\dfrac{3}{x+1}=0$, la droite $(\Delta)$ d'équation $y=x+1$ est \cle{asymptote oblique} à la courbe $(\mathcal{C})$ en $+\infty$.
Puisque $\dfrac{3}{x+1}>0$ pour tout $x\geq 1$, on a $C(x) > x+1$ : la courbe $(\mathcal{C})$ est \textbf{au-dessus} de son asymptote $(\Delta)$.
\emph{Interprétation économique :} pour un forage très profond ($x\to+\infty$), le terme $\frac{3}{x+1}$ (représentant les frais fixes amortis sur la profondeur) devient négligeable. Le coût marginal tend vers la pente de l'asymptote, soit $1$ centaine de milliers de FCFA par dizaine de mètres supplémentaire : c'est le «~coût de structure~» du forage.

\textbf{Partie B.} \textbf{4.} D'après la question 2, $C$ est continue et strictement croissante sur l'intervalle $[1\,;\,+\infty[$. Elle réalise donc une bijection de $[1\,;\,+\infty[$ sur son ensemble image $J=\left[\dfrac{7}{2}\,;\,+\infty\right[$. Or $15\in J$ car $15 > 3{,}5$. D'après le \textbf{théorème des valeurs intermédiaires} (version bijection), l'équation $C(x)=15$ admet une \textbf{unique} solution $\alpha$ dans $[1\,;\,+\infty[$.

\textbf{5.} On encadre $\alpha$ par dichotomie (valeurs approchées au millième) :
\begin{align*}
C(13) &= 14 + \frac{3}{14} \approx 14{,}214 < 15, \\
C(14) &= 15 + \frac{3}{15} = 15{,}2 > 15, \quad \text{donc } \alpha\in[13\,;\,14]. \\
C(13{,}7) &= 14{,}7 + \frac{3}{14{,}7} \approx 14{,}904 < 15, \\
C(13{,}8) &= 14{,}8 + \frac{3}{14{,}8} \approx 15{,}003 > 15, \quad \text{donc } \alpha\in[13{,}7\,;\,13{,}8].
\end{align*}
On a $13{,}7 < \alpha < 13{,}8$. La profondeur maximale finançable est d'environ \SI{137}{m} (on retient \SI{137}{m} par sécurité pour ne pas dépasser le budget).

\textbf{Partie C.} \textbf{6.} D'après la question 4, $C$ est une bijection de $[1\,;\,+\infty[$ sur $J=\left[\dfrac{7}{2}\,;\,+\infty\right[$. Sa fonction réciproque $C^{-1}$ est définie sur $J$ et associe à tout budget $y$ (en centaines de milliers de FCFA) la \textbf{profondeur exacte} (en dizaines de mètres) du forage réalisable. C'est l'outil de planification de la mairie : «~quel budget pour quelle profondeur ?~» devient «~quelle profondeur pour quel budget ?~».

\textbf{7.} $C$ est dérivable sur $[1\,;\,+\infty[$ et $C'(x)>0$ partout, en particulier $C'(\alpha)\neq 0$. La fonction réciproque $C^{-1}$ est donc dérivable en $15$ et la formule de la dérivée de la réciproque donne :
\[ \left(C^{-1}\right)'(15)=\dfrac{1}{C'(\alpha)}. \]
On détermine $\alpha$ exactement. L'équation $C(x)=15$ s'écrit $x+1+\frac{3}{x+1}=15$. En posant $t=x+1\ (\geq 2)$, on obtient $t+\frac{3}{t}=15$, soit $t^{2}-15t+3=0$.
Le discriminant est $\Delta = 225 - 12 = 213$. Les racines sont $t=\frac{15\pm\sqrt{213}}{2}$. Comme $t\geq 2$, on garde $t=\frac{15+\sqrt{213}}{2}\approx 14{,}797$ (l'autre racine $\approx 0{,}203$ est rejetée).
Alors $C'(\alpha)=1-\dfrac{3}{t^{2}}$. Numériquement, $t^{2}\approx 218{,}95$, donc $\frac{3}{t^{2}}\approx 0{,}0137$ et $C'(\alpha)\approx 0{,}9863$.
Finalement $\left(C^{-1}\right)'(15)\approx \dfrac{1}{0{,}9863}\approx 1{,}014$.
\emph{Interprétation :} au voisinage de \num{1500000}~FCFA, une augmentation du budget de \num{100000}~FCFA (soit $1$ unité sur l'axe des ordonnées) permet de gagner environ $1{,}014$ dizaines de mètres, soit environ \SI{10,1}{m} de profondeur supplémentaire.

\textbf{Partie D.} \textbf{8.} On dérive $C'$ : $C''(x)=\dfrac{6}{(x+1)^{3}}$. Sur $[3\,;\,4]$, $x+1>0$ donc $C''(x)>0$ : $C'$ est \textbf{strictement croissante} sur $[3\,;\,4]$. Son maximum est atteint en $x=4$ :
\[ \forall x\in[3\,;\,4],\quad C'(x)\leq C'(4)=1-\dfrac{3}{25}=\dfrac{22}{25}=0{,}88. \]

\textbf{9.} On applique l'\textbf{inégalité des accroissements finis} à $C$ sur $[3\,;\,4]$. Comme $C'>0$, $C(4)-C(3) = |C(4)-C(3)| \leq \max_{[3,4]}|C'| \times (4-3) \leq \dfrac{22}{25}\times 1 = 0{,}88$.
Le surcoût est donc majoré par $0{,}88$ centaines de milliers de FCFA, soit \num{88000}~FCFA. Cette valeur est bien inférieure à \num{120000}~FCFA ($1{,}2$ unités) : l'affirmation du sous-traitant est \textbf{fondée} (et même prudente).
Vérification exacte : $C(3)=\dfrac{9+6+4}{4}=\dfrac{19}{4}=4{,}75$ et $C(4)=\dfrac{16+8+4}{5}=\dfrac{28}{5}=5{,}6$, donc $C(4)-C(3)=0{,}85$, soit \num{85000}~FCFA. On confirme $0{,}85 \leq 0{,}88$.

\textbf{Partie E.} \textbf{10.} On place dans un repère orthonormé (unité \SI{0,5}{cm}) :
\begin{itemize}
\item la droite $(\Delta): y=x+1$ (pointillés orange) ;
\item la droite $y=x$ (pointillés gris) ;
\item des points de $(\mathcal{C})$ : $(1\,;\,3{,}5)$, $(3\,;\,4{,}75)$, $(4\,;\,5{,}6)$, $(\alpha\,;\,15)$ avec $\alpha\approx 13{,}8$ ;
\item la courbe $(\mathcal{C})$, strictement croissante, concave (car $C''>0$), au-dessus de $(\Delta)$ et la frôlant.
\end{itemize}

\textbf{11.} La courbe de $C^{-1}$ s'obtient par symétrie par rapport à la droite $y=x$ : elle part de $(3{,}5\,;\,1)$, passe par $(4{,}75\,;\,3)$, $(5{,}6\,;\,4)$, $(15\,;\,\alpha)$, et admet pour asymptote la symétrique de $(\Delta)$, c'est-à-dire la droite d'équation $y=x-1$ (pointillés verts).
\end{exoresolu}

\begin{popfigure}
\begin{center}
\begin{tikzpicture}[scale=0.42]
% Axes
\draw[->,popmuted] (-0.5,0) -- (17,0) node[below left] {\small $x$ (dizaines de m)};
\draw[->,popmuted] (0,-0.5) -- (0,17) node[below left] {\small $y$ (centaines de milliers FCFA)};
% Droite y=x
\draw[popline,dashed] (-0.5,-0.5) -- (16.5,16.5) node[above left,popmuted] {\small $y=x$};
% Asymptote y=x+1
\draw[poporange,dashed,thick,domain=-0.5:15.5,samples=50] plot (\x,{\x+1}) node[pos=0.85,sloped,above] {\small $y=x+1$};
% Asymptote y=x-1 (symétrique)
\draw[popgreen,dashed,thick,domain=0.5:16.5,samples=50] plot (\x,{\x-1}) node[pos=0.15,sloped,below] {\small $y=x-1$};
% Courbe C
\draw[popblue,very thick,domain=1:15.5,samples=200] plot (\x,{\x+1+3/(\x+1)});
% Courbe C^{-1} (paramétrée par t = y original)
\draw[popgreen,very thick,domain=3.5:16,samples=200] plot ({\x+1+3/(\x+1)},\x);
% Points clés
\fill[popink] (1,3.5) circle (2.5pt) node[above left] {\small $(1\,;\,3{,}5)$};
\fill[popink] (3.5,1) circle (2.5pt) node[below right] {\small $(3{,}5\,;\,1)$};
\fill[popink] (3,4.75) circle (2pt);
\fill[popink] (4.75,3) circle (2pt);
\fill[popink] (4,5.6) circle (2pt);
\fill[popink] (5.6,4) circle (2pt);
\fill[popink] (13.797,15) circle (3pt) node[above right] {\small $(\alpha\,;\,15)$};
\fill[popink] (15,13.797) circle (3pt) node[below right] {\small $(15\,;\,\alpha)$};
\end{tikzpicture}
\end{center}
\legende{En bleu, la courbe $(\mathcal{C})$ de $C$ et son asymptote oblique $(\Delta): y=x+1$ (orange) ; en vert, la courbe de $C^{-1}$ et son asymptote $y=x-1$, symétriques par rapport à la droite $y=x$ (gris).}
\end{popfigure}

\courssub{Bilan de l'activité}

Cette activité a permis de mettre en évidence que :

\begin{itemize}
\item le \textbf{TVI} est l'outil qui transforme une question concrète («~ce budget suffit-il ?~») en un théorème d'existence \emph{et} d'unicité, la stricte monotonie étant la clé de l'unicité ;
\item la \textbf{fonction réciproque} n'est pas un objet abstrait : c'est l'instrument du décideur, qui lit la profondeur en fonction du budget, et sa dérivée se calcule \emph{sans jamais l'expliciter} grâce à la formule $\left(f^{-1}\right)'(y)=\dfrac{1}{f'(x)}$ ;
\item l'\textbf{IAF} permet de trancher un débat commercial (l'affirmation du sous-traitant) par une majoration rigoureuse, sans calcul exact ;
\item l'\textbf{asymptote oblique} porte une information économique : le coût marginal de structure d'un forage profond ;
\item enfin, la \textbf{symétrie des courbes} de $f$ et $f^{-1}$ par rapport à $y=x$ offre une lecture graphique immédiate des deux points de vue : «~coût en fonction de la profondeur~» et «~profondeur en fonction du coût~».
\end{itemize}

\begin{savaistu}[L'eau, un enjeu national]
Selon les programmes camerounais d'hydraulique rurale et urbaine (MINEE, PNUD), des milliers de forages sont réalisés chaque année, notamment dans les régions de l'Ouest, du Nord et de l'Extrême-Nord. Le dimensionnement d'un forage (profondeur, pompe, budget) mobilise exactement les outils que tu viens d'utiliser : modélisation par une fonction, résolution d'équations et majoration des coûts. Les mathématiques de Terminale sont déjà un métier !
\end{savaistu}

\newpage

\coursec{Méthodes et exercices résolus}

\courssub{Méthode 1 : Démontrer l'existence et l'unicité d'une solution de $f(x)=c$ par le TVI}

\begin{methode}[Existence et unicité par le théorème des valeurs intermédiaires]
Pour résoudre une équation du type $f(x)=c$ sur un intervalle $I$ :
\begin{enumerate}
\item \textbf{Vérifier la continuité} de $f$ sur $I$ (fonction usuelle, somme, produit, quotient ou composée de fonctions continues).
\item \textbf{Étudier la monotonie stricte} : calculer $f'(x)$ et montrer que $f'$ garde un signe strict sur $I$ (sauf éventuellement en des points isolés où elle s'annule).
\item \textbf{Calculer les limites (ou valeurs) aux bornes} de $I$ et vérifier que $c$ est compris entre elles.
\item \textbf{Conclure} : $f$ réalise une bijection de $I$ sur $f(I)$, donc l'équation $f(x)=c$ admet une \cle{unique solution} $\alpha$ dans $I$.
\item \textbf{Encadrer $\alpha$} par balayage ou dichotomie : trouver $a<b$ tels que $f(a)$ et $f(b)$ encadrent $c$, puis affiner jusqu'à l'amplitude demandée.
\end{enumerate}
\end{methode}

\begin{exoresolu}[Équation $x^3+x-3=0$]
Montrer que l'équation $x^3+x-3=0$ admet une unique solution $\alpha$ dans $\mathbb{R}$, puis donner un encadrement de $\alpha$ d'amplitude $10^{-1}$.
\tcblower
Posons $f(x)=x^3+x-3$, définie sur $\mathbb{R}$.
\begin{enumerate}
\item $f$ est une fonction polynôme, donc \textbf{continue sur $\mathbb{R}$}.
\item $f'(x)=3x^2+1$. Pour tout $x\in\mathbb{R}$, $3x^2\geq 0$ donc $f'(x)\geq 1>0$. Ainsi $f$ est \textbf{strictement croissante} sur $\mathbb{R}$.
\item Limites aux bornes :
\[ \lim_{x\to-\infty}f(x)=\lim_{x\to-\infty}x^3=-\infty \qquad ; \qquad \lim_{x\to+\infty}f(x)=\lim_{x\to+\infty}x^3=+\infty. \]
\item $f$ est continue et strictement croissante sur $\mathbb{R}$ : elle réalise une bijection de $\mathbb{R}$ sur $\mathbb{R}$. Comme $0\in\mathbb{R}$, le théorème des valeurs intermédiaires garantit que l'équation $f(x)=0$ admet une \textbf{unique solution} $\alpha\in\mathbb{R}$.
\item Balayage : $f(1)=1+1-3=-1<0$ et $f(1{,}5)=3{,}375+1{,}5-3=1{,}875>0$, donc $\alpha\in]1\,;\,1{,}5[$. On affine au dixième : $f(1{,}2)=1{,}728+1{,}2-3=-0{,}072<0$ et $f(1{,}3)=2{,}197+1{,}3-3=0{,}497>0$.
\end{enumerate}
\textbf{Conclusion :} l'équation admet une unique solution $\alpha$ et $1{,}2<\alpha<1{,}3$ (encadrement d'amplitude $0{,}1=10^{-1}$).
\end{exoresolu}

\courssub{Méthode 2 : Étudier une fonction réciproque et tracer sa courbe}

\begin{methode}[Fonction réciproque]
\begin{enumerate}
\item \textbf{Justifier la bijectivité} : montrer que $f$ est continue et strictement monotone sur $I$ ; alors $f$ est une bijection de $I$ sur $J=f(I)$.
\item \textbf{Préciser $f^{-1}$} : ensemble de définition $J$, sens de variation (le \cle{même} que celui de $f$), continuité sur $J$.
\item \textbf{Dérivée de la réciproque} : si $f$ est dérivable en $a$ et $f'(a)\neq 0$, alors $f^{-1}$ est dérivable en $b=f(a)$ et
\[ (f^{-1})'(b)=\frac{1}{f'\bigl(f^{-1}(b)\bigr)}=\frac{1}{f'(a)}. \]
\item \textbf{Tracé} : la courbe de $f^{-1}$ est la \cle{symétrique} de celle de $f$ par rapport à la droite d'équation $y=x$ (première bissectrice). Un point $M(a\,;\,b)$ de $\mathcal{C}_f$ donne le point $M'(b\,;\,a)$ de $\mathcal{C}_{f^{-1}}$.
\end{enumerate}
\end{methode}

\begin{exoresolu}[{Réciproque de $f(x)=x^2+1$ sur $[0\,;\,+\infty[$}]
Soit $f(x)=x^2+1$ définie sur $I=[0\,;\,+\infty[$. (1) Montrer que $f$ admet une fonction réciproque $f^{-1}$ dont on précisera l'ensemble de définition. (2) Expliciter $f^{-1}(x)$. (3) Calculer $(f^{-1})'(5)$ de deux façons.
\tcblower
\begin{enumerate}
\item $f$ est polynôme donc continue sur $I$. $f'(x)=2x>0$ pour $x>0$, donc $f$ est strictement croissante sur $I$. Par le théorème de la bijection, $f$ réalise une bijection de $I$ sur $J=f(I)=[f(0)\,;\,\lim_{+\infty}f[=[1\,;\,+\infty[$. Donc $f^{-1}$ existe, définie sur $[1\,;\,+\infty[$, à valeurs dans $[0\,;\,+\infty[$.
\item Pour $x\geq 1$, résolvons $y=f^{-1}(x)$, c'est-à-dire $f(y)=x$ avec $y\geq 0$ :
\[ y^2+1=x \iff y^2=x-1 \iff y=\sqrt{x-1}\quad(\text{car } y\geq 0). \]
Donc $f^{-1}(x)=\sqrt{x-1}$.
\item \textbf{Première méthode} (dérivation directe) : $(f^{-1})'(x)=\dfrac{1}{2\sqrt{x-1}}$, donc $(f^{-1})'(5)=\dfrac{1}{2\sqrt{4}}=\dfrac{1}{4}$.
\textbf{Seconde méthode} (formule de la réciproque) : on cherche $a$ tel que $f(a)=5$ : $a^2+1=5$ et $a\geq 0$ donnent $a=2$. Comme $f'(2)=4\neq 0$, la formule s'applique :
\[ (f^{-1})'(5)=\frac{1}{f'\bigl(f^{-1}(5)\bigr)}=\frac{1}{f'(2)}=\frac{1}{2\times 2}=\frac{1}{4}. \]
Les deux résultats coïncident.
\end{enumerate}
\textbf{Conclusion :} $f^{-1}(x)=\sqrt{x-1}$ sur $[1\,;\,+\infty[$ et $(f^{-1})'(5)=\dfrac{1}{4}$.
\end{exoresolu}

\courssub{Méthode 3 : Résoudre $x^n=a$ et simplifier les puissances rationnelles}

\begin{methode}[Équation $x^n=a$ et racines n-ièmes]
\begin{enumerate}
\item \textbf{Discuter selon la parité de $n$ et le signe de $a$} :
\begin{itemize}
\item Si $a>0$ : une solution positive $\sqrt[n]{a}=a^{1/n}$ ; si $n$ est pair, il y a aussi la solution négative $-\sqrt[n]{a}$.
\item Si $a=0$ : unique solution $x=0$.
\item Si $a<0$ : solution $-\sqrt[n]{|a|}$ (négative) si $n$ est impair ; aucune solution réelle si $n$ est pair.
\end{itemize}
\item \textbf{Utiliser les règles de calcul} valables pour $a,b>0$ :
\[ \sqrt[n]{ab}=\sqrt[n]{a}\,\sqrt[n]{b},\qquad \sqrt[n]{a^n}=a,\qquad a^{p/q}=\sqrt[q]{a^p}=\bigl(\sqrt[q]{a}\bigr)^p,\qquad a^{-r}=\frac{1}{a^r}. \]
\item \textbf{Simplifier} en décomposant le radicande pour faire apparaître des puissances n-ièmes parfaites.
\end{enumerate}
\end{methode}

\begin{exoresolu}[Résolutions et simplifications]
(1) Résoudre dans $\mathbb{R}$ : $x^3=-27$ puis $x^4=16$. (2) Simplifier $A=\sqrt[3]{54}$ et $B=8^{2/3}\times 4^{-1/2}$.
\tcblower
\begin{enumerate}
\item \textbf{Équation $x^3=-27$ :} $n=3$ est impair, donc unique solution : $x=\sqrt[3]{-27}=-\sqrt[3]{27}=-3$. Vérification : $(-3)^3=-27$. $\mathcal{S}=\{-3\}$.
\textbf{Équation $x^4=16$ :} $n=4$ est pair et $a=16>0$, donc deux solutions : $x=\sqrt[4]{16}=2$ ou $x=-2$ (car $16=2^4$). $\mathcal{S}=\{-2\,;\,2\}$.
\item \textbf{Simplification de $A$ :} $54=27\times 2=3^3\times 2$, donc
\[ A=\sqrt[3]{3^3\times 2}=\sqrt[3]{3^3}\times\sqrt[3]{2}=3\sqrt[3]{2}. \]
\textbf{Simplification de $B$ :} $8^{2/3}=\bigl(\sqrt[3]{8}\bigr)^2=2^2=4$ et $4^{-1/2}=\dfrac{1}{4^{1/2}}=\dfrac{1}{\sqrt{4}}=\dfrac{1}{2}$. Donc
\[ B=4\times\frac{1}{2}=2. \]
\end{enumerate}
\textbf{Conclusion :} $\mathcal{S}_1=\{-3\}$, $\mathcal{S}_2=\{-2\,;\,2\}$, $A=3\sqrt[3]{2}$ et $B=2$.
\end{exoresolu}

\courssub{Méthode 4 : Encadrer $f(b)-f(a)$ par l'inégalité des accroissements finis}

\begin{methode}[Inégalité des accroissements finis (IAF)]
\begin{enumerate}
\item \textbf{Vérifier les hypothèses} : $f$ dérivable sur un intervalle contenant $[a\,;\,b]$.
\item \textbf{Encadrer la dérivée} : trouver deux réels $m$ et $M$ tels que pour tout $x\in[a\,;\,b]$, $m\leq f'(x)\leq M$.
\item \textbf{Appliquer la formule} (pour $a<b$) :
\[ m\,(b-a)\leq f(b)-f(a)\leq M\,(b-a). \]
Cas particulier utile : si $|f'(x)|\leq k$ sur $[a\,;\,b]$, alors $|f(b)-f(a)|\leq k\,|b-a|$ (\cle{majoration de l'erreur}).
\item \textbf{Conclure} par l'encadrement demandé de $f(b)$ à partir de $f(a)$.
\end{enumerate}
\end{methode}

\begin{exoresolu}[Encadrement de $\sqrt{102}$]
Soit $f(x)=\sqrt{x}$. (1) Montrer que pour tout $x\in[100\,;\,102]$, $\dfrac{1}{2\sqrt{102}}\leq f'(x)\leq \dfrac{1}{20}$. (2) En déduire, sans utiliser la valeur approchée de $\sqrt{102}$ donnée par une calculatrice, un encadrement de $\sqrt{102}$ d'amplitude inférieure à $10^{-2}$.
\tcblower
\begin{enumerate}
\item $f$ est dérivable sur $]0\,;\,+\infty[$ et $f'(x)=\dfrac{1}{2\sqrt{x}}$. La fonction $x\mapsto \sqrt{x}$ est croissante, donc $f'$ est \textbf{décroissante} sur $[100\,;\,102]$. Ainsi, pour tout $x\in[100\,;\,102]$ :
\[ f'(102)\leq f'(x)\leq f'(100) \iff \frac{1}{2\sqrt{102}}\leq f'(x)\leq \frac{1}{2\sqrt{100}}=\frac{1}{20}. \]
\item Appliquons l'IAF à $f$ sur $[100\,;\,102]$ avec $m=\dfrac{1}{2\sqrt{102}}$ et $M=\dfrac{1}{20}$ :
\[ \frac{1}{2\sqrt{102}}\,(102-100)\leq \sqrt{102}-\sqrt{100}\leq \frac{1}{20}\,(102-100) \]
\[ \frac{1}{\sqrt{102}}\leq \sqrt{102}-10\leq \frac{1}{10}. \]
La borne de gauche fait intervenir $\sqrt{102}$, que l'on cherche justement à encadrer : on la remplace par une minoration grossière mais rigoureuse. Comme $11^2=121>102$, on a $\sqrt{102}<11$, donc $\dfrac{1}{\sqrt{102}}>\dfrac{1}{11}$. Alors :
\[ \frac{1}{11}\leq \sqrt{102}-10\leq \frac{1}{10} \iff 10+\frac{1}{11}\leq \sqrt{102}\leq 10+\frac{1}{10}. \]
Or $\dfrac{1}{11}\approx 0{,}0909$ et $\dfrac{1}{10}=0{,}1$, d'où $10{,}090\leq \sqrt{102}\leq 10{,}100$.
\end{enumerate}
\textbf{Conclusion :} $10{,}09\leq \sqrt{102}\leq 10{,}1$, encadrement d'amplitude $0{,}01$ obtenu rigoureusement par l'inégalité des accroissements finis.
\end{exoresolu}

\courssub{Méthode 5 : Déterminer les branches infinies et les asymptotes}

\begin{methode}[Branches infinies : le réflexe en 4 questions]
On étudie le comportement de $f$ au voisinage d'une borne $x_0$ (finie ou infinie) où $|f|\to\infty$ ou $|x|\to\infty$.
\begin{enumerate}
\item \textbf{Asymptote verticale} : si $\lim_{x\to x_0}f(x)=\pm\infty$ ($x_0$ fini), la droite $x=x_0$ est asymptote verticale.
\item \textbf{Asymptote horizontale} : si $\lim_{x\to\pm\infty}f(x)=\ell$ (fini), la droite $y=\ell$ est asymptote horizontale.
\item \textbf{Sinon, direction asymptotique} : calculer $a=\lim_{x\to\pm\infty}\dfrac{f(x)}{x}$.
\begin{itemize}
\item Si $a=\pm\infty$ : \cle{branche parabolique} de direction $(Oy)$.
\item Si $a=0$ : branche parabolique de direction $(Ox)$.
\item Si $a$ fini non nul : calculer $b=\lim\bigl(f(x)-ax\bigr)$. Si $b$ fini : \textbf{asymptote oblique} $y=ax+b$ ; si $b=\pm\infty$ : branche parabolique de direction $y=ax$.
\end{itemize}
\item \textbf{Position relative} : étudier le signe de $f(x)-(ax+b)$ pour placer la courbe par rapport à l'asymptote oblique.
\end{enumerate}
\end{methode}

\begin{exoresolu}[Branches infinies de $f(x)=\dfrac{x^2-3x+3}{x-1}$]
Soit $f(x)=\dfrac{x^2-3x+3}{x-1}$. Déterminer toutes les branches infinies de $\mathcal{C}_f$ et préciser la position de la courbe par rapport à son asymptote oblique.
\tcblower
$\mathcal{D}_f=\mathbb{R}\setminus\{1\}$.
\begin{enumerate}
\item \textbf{En $x=1$ :} le numérateur en $1$ vaut $1-3+3=1\neq 0$. $\lim_{x\to 1^+}(x-1)=0^+$ donc $\lim_{x\to 1^+}f(x)=+\infty$ ; de même $\lim_{x\to 1^-}f(x)=-\infty$. La droite d'équation $x=1$ est \textbf{asymptote verticale}.
\item \textbf{En $\pm\infty$ :} $\lim_{x\to\pm\infty}f(x)=\lim_{x\to\pm\infty} \dfrac{x^2}{x}=\lim_{x\to\pm\infty} x=\pm\infty$ : pas d'asymptote horizontale. On calcule alors :
\[ \lim_{x\to\pm\infty}\frac{f(x)}{x}=\lim_{x\to\pm\infty}\frac{x^2}{x^2}=1=a\ (\text{fini, non nul}). \]
\item \textbf{Recherche de $b$ :} effectuons la division euclidienne du numérateur par $x-1$ : $x^2-3x+3=(x-1)(x-2)+1$, donc
\[ f(x)=x-2+\frac{1}{x-1}. \]
Alors $\lim_{x\to\pm\infty}\bigl(f(x)-(x-2)\bigr)=\lim_{x\to\pm\infty}\dfrac{1}{x-1}=0$ (limite finie). La droite d'équation $y=x-2$ est \textbf{asymptote oblique} en $-\infty$ et en $+\infty$.
\item \textbf{Position relative :} $f(x)-(x-2)=\dfrac{1}{x-1}$.
\begin{itemize}
\item Si $x>1$ : $f(x)-(x-2)>0$ : la courbe est \textbf{au-dessus} de l'asymptote.
\item Si $x<1$ : $f(x)-(x-2)<0$ : la courbe est \textbf{en dessous} de l'asymptote.
\end{itemize}
\end{enumerate}
\textbf{Conclusion :} $\mathcal{C}_f$ admet l'asymptote verticale d'équation $x=1$ et l'asymptote oblique d'équation $y=x-2$ ; la courbe est au-dessus de l'asymptote oblique pour $x>1$ et en dessous pour $x<1$ (le changement de position se fait de part et d'autre de l'asymptote verticale, la courbe n'étant pas définie en $1$).
\end{exoresolu}

\courssub{Méthode 6 : Mener une étude complète de fonction et tracer la courbe}

\begin{methode}[Plan d'étude complet (à rédiger dans cet ordre)]
\begin{enumerate}
\item \textbf{Ensemble de définition} $\mathcal{D}_f$ (dénominateur non nul, radicande positif ou nul, etc.) ; éventuelles symétries (parité, périodicité) pour \cle{réduire le domaine d'étude}.
\item \textbf{Limites aux bornes} de $\mathcal{D}_f$ et \textbf{branches infinies} (asymptotes).
\item \textbf{Dérivabilité et calcul de $f'$}, signe de $f'$ (tableau de signes si nécessaire).
\item \textbf{Tableau de variations} complet : limites, valeurs aux extrema, flèches.
\item \textbf{Éléments graphiques} : tangentes horizontales, points remarquables (intersections avec les axes), éventuel point d'inflexion.
\item \textbf{Tracé} : asymptotes d'abord, points remarquables, puis la courbe en respectant le tableau de variations.
\end{enumerate}
\end{methode}

\begin{exoresolu}[Étude complète de $f(x)=\dfrac{2x}{1+x^2}$]
Étudier complètement la fonction $f(x)=\dfrac{2x}{1+x^2}$ et tracer sa courbe représentative.
\tcblower
\textbf{1. Ensemble de définition.} Pour tout $x\in\mathbb{R}$, $1+x^2\geq 1>0$ : le dénominateur ne s'annule jamais. Donc $\mathcal{D}_f=\mathbb{R}$. De plus $f(-x)=\dfrac{-2x}{1+x^2}=-f(x)$ : $f$ est \textbf{impaire}, on peut restreindre l'étude à $[0\,;\,+\infty[$ et compléter par symétrie centrale de centre $O$.

\textbf{2. Limites.} $\lim_{x\to\pm\infty}f(x)=\lim_{x\to\pm\infty}\dfrac{2x}{x^2}=\lim_{x\to\pm\infty}\dfrac{2}{x}=0$. La droite d'équation $y=0$ (axe des abscisses) est \textbf{asymptote horizontale} en $-\infty$ et en $+\infty$. Il n'y a pas de branche infinie verticale.

\textbf{3. Dérivée.} $f$ est dérivable sur $\mathbb{R}$ (quotient de polynômes dont le dénominateur ne s'annule pas) :
\[ f'(x)=\frac{2(1+x^2)-2x\cdot 2x}{(1+x^2)^2}=\frac{2-2x^2}{(1+x^2)^2}=\frac{2(1-x)(1+x)}{(1+x^2)^2}. \]
Le dénominateur est strictement positif ; le signe de $f'$ est celui de $1-x^2$ :
$f'(x)>0$ sur $]-1\,;\,1[$, $f'(x)<0$ sur $]-\infty\,;\,-1[$ et sur $]1\,;\,+\infty[$, et $f'(-1)=f'(1)=0$.

\textbf{4. Tableau de variations.} Valeurs remarquables : $f(-1)=-1$, $f(0)=0$, $f(1)=1$.
\begin{center}
\begin{tabular}{|c|ccccccc|}
\hline
$x$ & $-\infty$ & & $-1$ & & $1$ & & $+\infty$ \\
\hline
$f'(x)$ & & $-$ & $0$ & $+$ & $0$ & $-$ & \\
\hline
$f(x)$ & $0$ & $\searrow$ & $-1$ & $\nearrow$ & $1$ & $\searrow$ & $0$ \\
\hline
\end{tabular}
\end{center}

\textbf{5. Points remarquables.} Intersection avec les axes : $f(x)=0\iff x=0$ : la courbe passe par l'origine. Tangentes horizontales en $A(-1\,;\,-1)$ (minimum) et $B(1\,;\,1)$ (maximum).

\textbf{6. Tracé.}
\begin{popfigure}
\begin{center}
\begin{tikzpicture}[scale=1.6]
\draw[popline,->] (-3.4,0) -- (3.4,0) node[below left] {$x$};
\draw[popline,->] (0,-1.8) -- (0,1.8) node[below left] {$y$};
\draw[popblue,very thick,domain=-3.3:3.3,samples=200,smooth] plot (\x,{2*\x/(1+\x*\x)});
\fill[popink] (1,1) circle (0.045) node[above right] {\small $B(1\,;\,1)$};
\fill[popink] (-1,-1) circle (0.045) node[below left] {\small $A(-1\,;\,-1)$};
\fill[popink] (0,0) circle (0.045) node[below right] {\small $O$};
\draw[poporange,dashed] (-1.4,-1) -- (1.4,-1);
\draw[poporange,dashed] (-1.4,1) -- (1.4,1);
\foreach \x in {-2,-1,1,2} \draw (\x,0.05) -- (\x,-0.05) node[below] {\small $\x$};
\end{tikzpicture}
\end{center}
\legende{Courbe de $f(x)=\dfrac{2x}{1+x^2}$ : impaire, asymptote $y=0$, extrema en $\pm 1$.}
\end{popfigure}
\textbf{Conclusion :} $f$ est impaire, croissante sur $[-1\,;\,1]$, décroissante sur $]-\infty\,;\,-1]$ et sur $[1\,;\,+\infty[$, bornée entre $-1$ et $1$, avec l'axe des abscisses pour asymptote horizontale.
\end{exoresolu}

\courssub{Méthode 7 : Lever les indéterminations avec les limites usuelles (trigonométrie et irrationnelles)}

\begin{methode}[Limites usuelles à mobiliser]
\begin{enumerate}
\item \textbf{Formes à reconnaître immédiatement} :
\[ \lim_{x\to 0}\frac{\sin x}{x}=1,\qquad \lim_{x\to 0}\frac{\tan x}{x}=1,\qquad \lim_{x\to 0}\frac{1-\cos x}{x^2}=\frac{1}{2},\qquad \lim_{x\to 0}\frac{\sin(ax)}{x}=a. \]
\item \textbf{Changement de variable} : si $x\to x_0$, poser $t=x-x_0$ pour se ramener en $0$.
\item \textbf{Fonctions irrationnelles} : pour une forme $\infty-\infty$ avec des racines carrées, multiplier par la \cle{quantité conjuguée} :
\[ \sqrt{A}-\sqrt{B}=\frac{A-B}{\sqrt{A}+\sqrt{B}}. \]
\item \textbf{Encadrement} : si la limite usuelle ne s'applique pas directement, utiliser $|\sin x|\leq 1$, $|\cos x|\leq 1$ avec le théorème des gendarmes.
\end{enumerate}
\end{methode}

\begin{exoresolu}[Trois limites classiques]
Calculer : (1) $\lim_{x\to 0}\dfrac{\sin(3x)}{2x}$ ; (2) $\lim_{x\to+\infty}\bigl(\sqrt{x^2+3x}-x\bigr)$ ; (3) $\lim_{x\to+\infty}\dfrac{\sin x}{x}$.
\tcblower
\begin{enumerate}
\item \textbf{Limite 1.} Forme indéterminée $\dfrac{0}{0}$. On fait apparaître la limite usuelle :
\[ \frac{\sin(3x)}{2x}=\frac{3}{2}\times\frac{\sin(3x)}{3x}. \]
Posons $t=3x$ : quand $x\to 0$, $t\to 0$, et $\lim_{t\to 0}\dfrac{\sin t}{t}=1$. Donc
\[ \lim_{x\to 0}\frac{\sin(3x)}{2x}=\frac{3}{2}\times 1=\frac{3}{2}. \]
\item \textbf{Limite 2.} Forme $\infty-\infty$. Multiplions par la quantité conjuguée :
\[ \sqrt{x^2+3x}-x=\frac{(x^2+3x)-x^2}{\sqrt{x^2+3x}+x}=\frac{3x}{\sqrt{x^2+3x}+x}. \]
Pour $x>0$, $\sqrt{x^2+3x}=x\sqrt{1+\frac{3}{x}}$, donc
\[ \frac{3x}{x\sqrt{1+\frac{3}{x}}+x}=\frac{3}{\sqrt{1+\frac{3}{x}}+1}\xrightarrow[x\to+\infty]{}\frac{3}{\sqrt{1}+1}=\frac{3}{2}. \]
(Interprétation géométrique : la courbe d'équation $y=\sqrt{x^2+3x}$ admet en $+\infty$ l'asymptote oblique d'équation $y=x+\frac{3}{2}$.)
\item \textbf{Limite 3.} Le numérateur n'a pas de limite, mais il est \textbf{borné} : pour $x>0$,
\[ \left|\frac{\sin x}{x}\right|\leq \frac{1}{x}\xrightarrow[x\to+\infty]{}0. \]
Par le théorème des gendarmes, $\lim_{x\to+\infty}\dfrac{\sin x}{x}=0$.
\end{enumerate}
\textbf{Conclusion :} les trois limites valent respectivement $\dfrac{3}{2}$, $\dfrac{3}{2}$ et $0$.
\end{exoresolu}

\begin{attention}[Les 5 erreurs qui coûtent des points au Bac]
\begin{enumerate}
\item \textbf{Appliquer le TVI sans vérifier les hypothèses.} Il faut rédiger explicitement : « $f$ est \emph{continue} sur $I$ » ET « $f$ est \emph{strictement monotone} sur $I$ ». Sans la stricte monotonie, on obtient l'existence mais pas l'\textbf{unicité}. Oublier l'une des deux conditions fait perdre les points de justification.
\item \textbf{Confondre $(f^{-1})'(x)$ et $\dfrac{1}{f'(x)}$.} La formule est $(f^{-1})'(b)=\dfrac{1}{f'\bigl(f^{-1}(b)\bigr)}$ : on évalue $f'$ en $f^{-1}(b)$, pas en $b$. Et la formule n'est valable que si $f'\bigl(f^{-1}(b)\bigr)\neq 0$.
\item \textbf{Oublier la solution négative de $x^n=a$ avec $n$ pair.} $x^4=16$ a \emph{deux} solutions : $2$ et $-2$. À l'inverse, écrire $\sqrt[4]{16}=-2$ est faux : la racine n-ième ($n$ pair) d'un réel positif est toujours \textbf{positive}.
\item \textbf{Conclure trop vite sur les branches infinies.} $\lim_{x\to\pm\infty} f(x)=\pm\infty$ ne signifie pas « asymptote verticale » ; il faut calculer $\lim \dfrac{f(x)}{x}$. Et une asymptote oblique exige que $\lim\bigl(f(x)-ax\bigr)$ soit \emph{finie} : si elle est infinie, c'est une branche parabolique, pas une asymptote.
\item \textbf{Erreurs de calcul sur les dérivées de quotients et les quantités conjuguées.} Dans $\left(\dfrac{u}{v}\right)'=\dfrac{u'v-uv'}{v^2}$, l'ordre $u'v-uv'$ est crucial (signe inversé sinon). Pour $\sqrt{A}-\sqrt{B}$, le conjugué donne $\dfrac{A-B}{\sqrt{A}+\sqrt{B}}$ : le numérateur est $A-B$, et non $A+B$. Vérifie chaque étape avant de conclure.
\end{enumerate}
\end{attention}

\newpage

\coursec{Exercices}

\coursec{Étude de fonctions numériques d'une variable réelle}

\courssub{Image d'un intervalle par une fonction continue et Théorème des valeurs intermédiaires}

\begin{aretenir}[Théorème des valeurs intermédiaires (TVI)]
Soit $f$ une fonction continue sur un intervalle fermé $[a\,;b]$ (avec $a<b$).
\begin{itemize}
\item \textbf{Existence :} Pour tout réel $k$ compris entre $f(a)$ et $f(b)$, il existe au moins un $c\in[a\,;b]$ tel que $f(c)=k$.
\item \textbf{Unicité :} Si de plus $f$ est \textbf{strictement monotone} sur $[a\,;b]$, alors ce $c$ est \textbf{unique}.
\end{itemize}
\end{aretenir}

\begin{propriete}[Image d'un intervalle par une fonction continue]
Si $f$ est continue sur un intervalle $I$ (quel qu'il soit : ouvert, fermé, borné ou non), alors l'image $f(I)$ est aussi un intervalle.
\end{propriete}

\begin{methode}[Montrer l'existence et l'unicité d'une solution de $f(x)=c$]
\begin{enumerate}
\item Vérifier que $f$ est continue sur un intervalle $I=[a\,;b]$ (ou $]a\,;b]$, $[a\,;b[$, $]a\,;b[$ en utilisant les limites aux bornes).
\item Étudier les variations de $f$ sur $I$ (signe de $f'$) pour établir sa monotonie stricte.
\item Calculer $f(a)$ et $f(b)$ (ou les limites aux bornes si l'intervalle est ouvert) et vérifier que $c$ est bien \emph{strictement} compris entre ces deux valeurs.
\item Conclure par le TVI : l'équation $f(x)=c$ admet une unique solution dans $I$.
\end{enumerate}
\end{methode}

\begin{exemplebox}[Application du TVI]
Soit $f(x)=x^3+2x-3$ sur $\mathbb{R}$.
\begin{itemize}
\item $f$ est continue sur $\mathbb{R}$ (fonction polynôme).
\item $f'(x)=3x^2+2 > 0$ sur $\mathbb{R}$, donc $f$ est strictement croissante.
\item $\lim_{x\to-\infty}f(x)=-\infty$ et $\lim_{x\to+\infty}f(x)=+\infty$.
\item Pour $c=0$, $0\in]-\infty\,;+\infty[$. L'équation $f(x)=0$ admet une unique solution $\alpha\in\mathbb{R}$.
\item Comme $f(0)=-3<0$ et $f(2)=9>0$, on a $\alpha\in\,]0\,;2[$.
\end{itemize}
\end{exemplebox}

\begin{attention}[Pièges à éviter]
\begin{itemize}
\item Ne pas oublier de vérifier la \textbf{continuité} sur l'intervalle \textbf{fermé} $[a\,;b]$ (si l'intervalle est ouvert, il faut prolonger par continuité ou utiliser les limites).
\item L'unicité nécessite une \textbf{monotonie stricte} (dérivée de signe constant et non nulle sur un intervalle, ou étude de signe rigoureuse).
\item Le TVI ne donne pas la valeur de la solution, seulement son existence. Pour l'approcher, on utilise la dichotomie ou des suites.
\end{itemize}
\end{attention}

\courssub{Fonction réciproque d'une bijection}

\begin{definition}[Bijection et fonction réciproque]
Une fonction $f : I \to J$ est une \textbf{bijection} si elle est à la fois injective (valeurs distinctes ont des images distinctes) et surjective (tout élément de $J$ a un antécédent dans $I$).
La \textbf{fonction réciproque} $f^{-1} : J \to I$ est définie par : $f^{-1}(y)=x \iff f(x)=y$.
\end{definition}

\begin{propriete}[Propriétés de la fonction réciproque]
Soit $f$ une bijection d'un intervalle $I$ sur un intervalle $J$.
\begin{enumerate}
\item \textbf{Continuité :} Si $f$ est continue sur $I$, alors $f^{-1}$ est continue sur $J$.
\item \textbf{Dérivabilité :} Si $f$ est dérivable sur $I$ et si $f'(x)\neq 0$ pour tout $x\in I$, alors $f^{-1}$ est dérivable sur $J$ et pour tout $y\in J$ :
\[ (f^{-1})'(y) = \frac{1}{f'(f^{-1}(y))}. \]
\item \textbf{Symétrie graphique :} Dans un repère orthonormé, les courbes $\mathcal{C}_f$ et $\mathcal{C}_{f^{-1}}$ sont symétriques l'une de l'autre par rapport à la droite d'équation $y=x$ (première bissectrice).
\end{enumerate}
\end{propriete}

\begin{methode}[Déterminer l'expression de $f^{-1}$]
\begin{enumerate}
\item Vérifier que $f$ est une bijection de $I$ sur $J$ (continuité + monotonie stricte + bornes).
\item Poser $y = f(x)$ avec $x\in I$, $y\in J$.
\item Exprimer $x$ en fonction de $y$ (résoudre l'équation par rapport à $x$).
\item Échanger les lettres $x$ et $y$ (usage) : on obtient $f^{-1}(x) = \ldots$ pour $x\in J$.
\end{enumerate}
\end{methode}

\begin{exemplebox}[{Réciproque de $g(x)=x^2-4x+5$ sur $[2\,;+\infty[$}]
\begin{itemize}
\item $g(x)=(x-2)^2+1$. Sur $[2\,;+\infty[$, $g$ est strictement croissante, $g(2)=1$, $\lim_{+\infty}g=+\infty$. Bijection $[2\,;+\infty[ \to [1\,;+\infty[$.
\item $y = (x-2)^2+1 \implies (x-2)^2 = y-1 \implies x-2 = \sqrt{y-1}$ (car $x\ge 2$).
\item $x = 2+\sqrt{y-1}$. Donc $g^{-1}(x) = 2+\sqrt{x-1}$ pour $x\in[1\,;+\infty[$.
\item Dérivée : $(g^{-1})'(x) = \frac{1}{2\sqrt{x-1}}$. En $x=2$ : $(g^{-1})'(2)=\frac{1}{2}$.
\item Vérification par la formule : $g^{-1}(2)=3$, $g'(3)=2(3)-4=2$, $(g^{-1})'(2)=\frac{1}{g'(3)}=\frac{1}{2}$.
\end{itemize}
\end{exemplebox}

\begin{attention}[Condition $f'(x)\neq 0$]
La formule $(f^{-1})'(y) = 1/f'(f^{-1}(y))$ \textbf{nécessite} que $f'(f^{-1}(y))\neq 0$. Si $f'(x_0)=0$, alors $f^{-1}$ n'est \textbf{pas dérivable} en $y_0=f(x_0)$ (la tangente à $\mathcal{C}_{f^{-1}}$ est verticale).
\end{attention}

\courssub{Racines n-ièmes et puissances rationnelles}

\begin{definition}[Racine n-ième]
Soit $n\in\mathbb{N}$, $n\ge 2$ et $a\in\mathbb{R}_+^*$. On appelle \textbf{racine n-ième de $a$} l'unique réel positif $b$ tel que $b^n=a$. On le note $\sqrt[n]{a}$ ou $a^{1/n}$.
Par convention $\sqrt[n]{0}=0$. Pour $a<0$ et $n$ impair, $\sqrt[n]{a} = -\sqrt[n]{|a|}$.
\end{definition}

\begin{propriete}[Puissances rationnelles]
Pour $a>0$ et $q\in\mathbb{N}^*$, on pose $a^{1/q}=\sqrt[q]{a}$.
Pour $p\in\mathbb{Z}$, $q\in\mathbb{N}^*$, on pose $a^{p/q} = (a^{1/q})^p = \sqrt[q]{a^p}$.
Les règles usuelles des puissances s'étendent : $a^r a^s = a^{r+s}$, $(a^r)^s = a^{rs}$, $(ab)^r = a^r b^r$ (pour $a,b>0$).
\end{propriete}

\begin{aretenir}[Simplification d'expressions avec racines]
\begin{itemize}
\item $\sqrt[n]{a^n} = a$ pour $a\ge 0$ (et $|a|$ si $n$ pair et $a\in\mathbb{R}$).
\item $\sqrt[n]{ab} = \sqrt[n]{a}\sqrt[n]{b}$ ($a,b\ge 0$).
\item $\sqrt[n]{\frac{a}{b}} = \frac{\sqrt[n]{a}}{\sqrt[n]{b}}$ ($a\ge 0, b>0$).
\item $\sqrt[m]{\sqrt[n]{a}} = \sqrt[mn]{a}$.
\end{itemize}
\end{aretenir}

\begin{exemplebox}[Calculs sur les puissances rationnelles]
\begin{align*}
A &= \frac{16^{3/4}\times 9^{1/2}}{8^{2/3}} = \frac{(2^4)^{3/4}\times (3^2)^{1/2}}{(2^3)^{2/3}} = \frac{2^3\times 3}{2^2} = \frac{8\times 3}{4} = 6. \\[0.5em]
B &= \sqrt[3]{54}+\sqrt[3]{16} = \sqrt[3]{27\times 2}+\sqrt[3]{8\times 2} = 3\sqrt[3]{2}+2\sqrt[3]{2} = 5\sqrt[3]{2}.
\end{align*}
\end{exemplebox}

\courssub{Inégalité des accroissements finis}

\begin{propriete}[Inégalité des accroissements finis (Théorème des accroissements finis)]
Soit $f$ une fonction continue sur $[a\,;b]$ et dérivable sur $]a\,;b[$.
Il existe au moins un $c\in\,]a\,;b[$ tel que :
\[ f(b)-f(a) = f'(c)\,(b-a). \]
\textbf{Conséquence majeure (encadrement) :} Si $|f'(x)|\le M$ sur $]a\,;b[$, alors
\[ |f(b)-f(a)| \le M\,|b-a|. \]
\end{propriete}

\begin{methode}[Utiliser l'inégalité des accroissements finis pour encadrer une différence]
\begin{enumerate}
\item Identifier la fonction $f$ et l'intervalle $[a\,;b]$ (ou $[b\,;a]$).
\item Calculer $f'(x)$ et trouver une majoration $M \ge |f'(x)|$ sur l'intervalle.
\item Appliquer $|f(b)-f(a)| \le M\,|b-a|$.
\end{enumerate}
\end{methode}

\begin{exemplebox}[Encadrement de $\sqrt{4,1}$]
Soit $f(x)=\sqrt{x}$ sur $[4\,;4,1]$. $f'(x)=\frac{1}{2\sqrt{x}} \le \frac{1}{2\sqrt{4}} = \frac{1}{4}$.
$|\sqrt{4,1}-\sqrt{4}| \le \frac{1}{4}\times 0,1 = 0,025$.
Donc $\sqrt{4,1} \in [2-0,025\,;\,2+0,025] = [1,975\,;\,2,025]$.
\end{exemplebox}

\begin{exemplebox}[Convergence d'une suite par contraction]
Soit $u_0=0$ et $u_{n+1} = \sqrt{u_n+4}-2$. (On note $g(x)=\sqrt{x+4}-2$).
$g'(x)=\frac{1}{2\sqrt{x+4}}$. Pour $x\ge 0$, $g'(x)\le \frac{1}{4}$.
$|u_{n+1}-0| = |g(u_n)-g(0)| \le \frac{1}{4}|u_n-0|$.
Donc $|u_n| \le \frac{1}{4^n}|u_0| \to 0$. La suite converge vers $0$.
\end{exemplebox}

\courssub{Limites usuelles et branches infinies}

\begin{aretenir}[Limites usuelles à connaître (Bac)]
\begin{align*}
&\lim_{x\to 0}\frac{\sin x}{x}=1,\quad \lim_{x\to 0}\frac{1-\cos x}{x}=0,\quad \lim_{x\to 0}\frac{1-\cos x}{x^2}=\frac{1}{2},\\[0.5em]
&\lim_{x\to 0}\frac{\tan x}{x}=1,\quad \lim_{x\to 0}\frac{\ln(1+x)}{x}=1,\quad \lim_{x\to 0}\frac{e^x-1}{x}=1,\\[0.5em]
&\lim_{x\to 0}\frac{\sqrt{1+x}-1}{x}=\frac{1}{2},\quad \lim_{x\to 0}\frac{(1+x)^\alpha-1}{x}=\alpha\ (\alpha\in\mathbb{R}).
\end{align*}
\end{aretenir}

\begin{propriete}[Branches infinies et asymptotes]
Soit $f$ définie sur un voisinage de $+\infty$ (ou $-\infty$) ou au voisinage d'une valeur interdite $x_0$.
\begin{itemize}
\item \textbf{Asymptote verticale} $x=x_0$ : $\lim_{x\to x_0^\pm}f(x)=\pm\infty$.
\item \textbf{Asymptote horizontale} $y=\ell$ : $\lim_{x\to\pm\infty}f(x)=\ell\in\mathbb{R}$.
\item \textbf{Asymptote oblique} $y=ax+b$ ($a\neq 0$) : $\lim_{x\to\pm\infty}\frac{f(x)}{x}=a\in\mathbb{R}^*$ et $\lim_{x\to\pm\infty}\big(f(x)-ax\big)=b\in\mathbb{R}$.
\item \textbf{Direction asymptotique} : $\lim_{x\to\pm\infty}\frac{f(x)}{x}=a\in\mathbb{R}^*$ existe mais $\lim_{x\to\pm\infty}(f(x)-ax)$ n'existe pas (ou est infini).
\end{itemize}
\end{propriete}

\begin{methode}[Recherche systématique des asymptotes]
\begin{enumerate}
\item \textbf{Verticales :} Chercher les valeurs $x_0$ non appartennant au domaine où $\lim f = \pm\infty$.
\item \textbf{Horizontales :} Calculer $\lim_{x\to\pm\infty}f(x)$. Si fini $\ell$, $y=\ell$ est asymptote.
\item \textbf{Obliques :} Si pas d'horizontale, calculer $a=\lim_{x\to\pm\infty}\frac{f(x)}{x}$. Si $a\in\mathbb{R}^*$, calculer $b=\lim_{x\to\pm\infty}(f(x)-ax)$. Si $b\in\mathbb{R}$, $y=ax+b$ est asymptote.
\item \textbf{Position relative :} Étudier le signe de $f(x)-(ax+b)$ pour savoir si la courbe est au-dessus ou en dessous de l'asymptote.
\end{enumerate}
\end{methode}

\begin{exemplebox}[Asymptotes de $f(x)=\frac{x^2+1}{x-1}$]
Domaine : $\mathbb{R}\setminus\{1\}$.
\begin{itemize}
\item $x\to 1^\pm$ : num$\to 2$, dén$\to 0^\pm$ $\Rightarrow$ asymptote verticale $x=1$.
\item $x\to\pm\infty$ : $\frac{f(x)}{x}=\frac{x^2+1}{x(x-1)}\to 1=a$.
\item $f(x)-x = \frac{x^2+1-x(x-1)}{x-1}=\frac{x+1}{x-1}\to 1=b$.
\item Asymptote oblique : $y=x+1$.
\item Position : $f(x)-(x+1)=\frac{2}{x-1}$. Signe de $x-1$ : courbe au-dessus pour $x>1$, en-dessous pour $x<1$.
\end{itemize}
\end{exemplebox}

\begin{savaistu}[Direction asymptotique sans asymptote]
La fonction $\varphi(x)=x+\cos x$ a pour direction asymptotique $y=x$ car $\varphi(x)/x\to 1$. Mais $\varphi(x)-x=\cos x$ n'a pas de limite en $\pm\infty$ (oscille entre $-1$ et $1$). Il n'y a \textbf{pas d'asymptote oblique}.
\end{savaistu}

\courssub{Étude complète de fonctions et tracé de courbes}

\begin{methode}[Étude complète d'une fonction $f$ (7 étapes)]
\begin{enumerate}
\item \textbf{Domaine de définition $D_f$} (et parité/périodicité si évidente).
\item \textbf{Limites aux bornes de $D_f$} (valeurs interdites, $\pm\infty$) $\to$ asymptotes verticales/horizontales.
\item \textbf{Dérivée $f'$} : calcul, simplification, signe $\to$ tableau de variations.
\item \textbf{Asymptotes obliques} (si pas d'horizontale) : $a=\lim f/x$, $b=\lim f-ax$. Position relative.
\item \textbf{Points particuliers} : intersections avec les axes ($f(x)=0$, $f(0)$), extremums, points d'inflexion (si $f''$ simple).
\item \textbf{Tableau de variations} complet (valeurs aux bornes, extremums).
\item \textbf{Tracé de la courbe} $\mathcal{C}_f$ dans un repère orthonormé (échelle adaptée, asymptotes en pointillés, points caractéristiques).
\end{enumerate}
\end{methode}

\begin{exoresolu}[Étude complète : fonction rationnelle $h(x)=\frac{x^2+3x+1}{x+2}$]
\textbf{1. Domaine :} $D_h=\mathbb{R}\setminus\{-2\}$.
\textbf{2. Limites :}
$x\to -2^-$ : num$\to -1$, dén$\to 0^-$ $\Rightarrow +\infty$.
$x\to -2^+$ : dén$\to 0^+$ $\Rightarrow -\infty$. Asymptote verticale $x=-2$.
$x\to\pm\infty$ : $h(x)\sim x \to \pm\infty$. Pas d'horizontale.
\textbf{3. Dérivée :} $h'(x)=\frac{(2x+3)(x+2)-(x^2+3x+1)}{(x+2)^2}=\frac{x^2+4x+5}{(x+2)^2}$.
$\Delta' = 16-20=-4<0$, numérateur $>0$. $h'>0$ sur $]-\infty\,;-2[$ et $]-2\,;+\infty[$.
$h$ strictement croissante sur chaque intervalle.
\textbf{4. Asymptote oblique :} Division : $h(x)=x+1-\frac{1}{x+2}$.
$a=1$, $b=1$. Asymptote $\Delta : y=x+1$.
Position : $h(x)-(x+1)=-\frac{1}{x-2}$... non $-\frac{1}{x+2}$.
$x>-2 \Rightarrow x+2>0 \Rightarrow h(x) < x+1$ (courbe sous $\Delta$).
$x<-2 \Rightarrow x+2<0 \Rightarrow h(x) > x+1$ (courbe au-dessus $\Delta$).
\textbf{5. Points :} $h(0)=1/2$. $h(x)=0 \Leftrightarrow x^2+3x+1=0 \Rightarrow x=\frac{-3\pm\sqrt{5}}{2}$.
\textbf{6. Tableau de variations :}
\begin{center}
\begin{tabular}{|c|cccc|cccc|}\hline
$x$ & $-\infty$ & & $-2$ & & & $-2$ & & $+\infty$ \\ \hline
$h'(x)$ & & $+$ & $\vert$ & & & $\vert$ & $+$ & \\ \hline
$h$ & $-\infty$ & $\nearrow$ & $+\infty$ & & & $-\infty$ & $\nearrow$ & $+\infty$ \\ \hline
\end{tabular}
\end{center}
\textbf{7. Tracé :} Voir figure ci-dessous.
\begin{popfigure}
\begin{tikzpicture}[scale=0.7, domain=-6:4]
\draw[popline, thin] (-6,-5) grid (4,5);
\draw[->, thick] (-6.5,0) -- (4.5,0) node[right] {$x$};
\draw[->, thick] (0,-5.5) -- (0,5.5) node[above] {$y$};
\foreach \x in {-5,-4,-3,-2,-1,1,2,3,4} \draw (\x,0.1) -- (\x,-0.1) node[below] {\small $\x$};
\foreach \y in {-5,-4,-3,-2,-1,1,2,3,4,5} \draw (0.1,\y) -- (-0.1,\y) node[left] {\small $\y$};
% Asymptotes
\draw[poporange, dashed, thick] (-2,-5.5) -- (-2,5.5) node[above right] {$x=-2$};
\draw[poporange, dashed, thick] (-6.5,-5.5) -- (4.5,5.5) node[above right] {$y=x+1$};
% Courbe branche gauche
\draw[popblue, very thick, samples=100, domain=-6:-2.1] plot (\x, {(\x*\x+3*\x+1)/(\x+2)});
% Courbe branche droite
\draw[popblue, very thick, samples=100, domain=-1.9:4] plot (\x, {(\x*\x+3*\x+1)/(\x+2)});
% Points
\fill[popblue] (0,0.5) circle (2pt) node[above right] {$(0;0,5)$};
\fill[popblue] ({(-3+sqrt(5))/2},0) circle (2pt) node[below right] {$\alpha_1$};
\fill[popblue] ({(-3-sqrt(5))/2},0) circle (2pt) node[above left] {$\alpha_2$};
\end{tikzpicture}
\legende{Courbe de $h(x)=\frac{x^2+3x+1}{x+2}$ et ses asymptotes}
\end{popfigure}
\end{exoresolu}

\begin{exoresolu}[Étude complète : fonction irrationnelle $\psi(x)=\frac{x}{\sqrt{x^2+1}}$]
\textbf{1. Domaine/Parité :} $D_\psi=\mathbb{R}$. $\psi(-x)=-\psi(x)$ : \textbf{fonction impaire}. Courbe symétrique par rapport à l'origine $O$.
\textbf{2. Limites :} $x\to+\infty$, $\psi(x)=\frac{x}{x\sqrt{1+1/x^2}}\to 1$. $x\to-\infty$, $\psi(x)\to -1$.
Asymptotes horizontales : $y=1$ (en $+\infty$) et $y=-1$ (en $-\infty$).
\textbf{3. Dérivée :} $\psi'(x)=\frac{\sqrt{x^2+1}-x\frac{x}{\sqrt{x^2+1}}}{x^2+1}=\frac{1}{(x^2+1)^{3/2}} > 0$.
$\psi$ strictement croissante sur $\mathbb{R}$.
\textbf{4. Variations :} Tableau simple : $-\infty \to -1$, croissante, $+\infty \to 1$.
\textbf{5. Bijection :} $\psi$ bijection $\mathbb{R} \to \,]-1\,;1[$.
Réciproque : $y=\frac{x}{\sqrt{x^2+1}} \implies y^2=\frac{x^2}{x^2+1} \implies x^2=\frac{y^2}{1-y^2} \implies x=\frac{y}{\sqrt{1-y^2}}$ (même signe).
$\psi^{-1}(y)=\frac{y}{\sqrt{1-y^2}}$ sur $]-1\,;1[$.
$(\psi^{-1})'(0)=\frac{1}{\psi'(\psi^{-1}(0))}=\frac{1}{\psi'(0)}=1$.
\textbf{6. Équation $\psi(x)=1/2$ :} $x=\frac{1/2}{\sqrt{3/4}}=\frac{1}{\sqrt{3}}=\frac{\sqrt{3}}{3}$.
\end{exoresolu}

\courssub{Je vérifie mes connaissances}

\begin{exercice}[QCM]
Pour chaque question, une seule réponse est exacte. Recopier la bonne réponse.
\begin{enumerate}
\item Soit $f$ une fonction continue sur $[a\,;b]$ telle que $f(a)=-2$ et $f(b)=5$. L'équation $f(x)=0$ :
\begin{enumerate}
\item n'admet aucune solution dans $[a\,;b]$ ;
\item admet au moins une solution dans $[a\,;b]$ ;
\item admet exactement une solution dans $[a\,;b]$.
\end{enumerate}
\item Si $f$ est une bijection de $I$ sur $J$, dérivable, et si $f'(x_0)=0$ pour un $x_0\in I$, alors $f^{-1}$ :
\begin{enumerate}
\item est dérivable en $f(x_0)$ et $(f^{-1})'(f(x_0))=0$ ;
\item n'est pas dérivable en $f(x_0)$ ;
\item n'existe pas.
\end{enumerate}
\item L'équation $x^4=81$ admet dans $\mathbb{R}$ :
\begin{enumerate}
\item une seule solution ; \item deux solutions ; \item quatre solutions.
\end{enumerate}
\item La courbe d'une fonction $f$ admet une asymptote horizontale en $+\infty$ si et seulement si :
\begin{enumerate}
\item $\lim_{x\to+\infty}f(x)=+\infty$ ;
\item $\lim_{x\to+\infty}f(x)$ est un réel fini $\ell$ ;
\item $\lim_{x\to+\infty}\dfrac{f(x)}{x}$ est un réel fini.
\end{enumerate}
\end{enumerate}
\end{exercice}

\begin{exercice}[Vrai ou faux, justifier]
Répondre par Vrai ou Faux en justifiant chaque réponse (par une démonstration courte ou un contre-exemple).
\begin{enumerate}
\item Toute fonction continue sur $[0\,;1]$ avec $f(0)=f(1)=3$ s'annule au moins une fois sur $]0\,;1[$.
\item Si $f$ est continue et strictement monotone sur $\mathbb{R}$, alors $f$ est une bijection de $\mathbb{R}$ sur $\mathbb{R}$.
\item Pour tout réel $a>0$ et tout entier $n\geq 2$, $\sqrt[n]{a^n}=a$.
\item Si $\lim_{x\to+\infty}\big(f(x)-(2x+1)\big)=0$, alors la droite d'équation $y=2x+1$ est asymptote à la courbe de $f$ en $+\infty$.
\end{enumerate}
\end{exercice}

\begin{exercice}[Calculs directs]
\begin{enumerate}
\item Résoudre dans $\mathbb{R}$ les équations : $x^3=27$ ; $x^4=81$ ; $x^5=-32$.
\item Simplifier $A=\dfrac{16^{3/4}\times 9^{1/2}}{8^{2/3}}$ et écrire $B=\sqrt[3]{54}+\sqrt[3]{16}$ sous la forme $k\sqrt[3]{2}$ avec $k\in\mathbb{N}$.
\item Calculer les limites suivantes :
\[ \lim_{x\to 0}\frac{\sin(5x)}{2x}\;;\qquad \lim_{x\to 0}\frac{1-\cos x}{x\sin x}\;;\qquad \lim_{x\to+\infty}\frac{x+\sin x}{x}. \]
\end{enumerate}
\end{exercice}

\begin{exercice}[Définitions et cours]
\begin{enumerate}
\item Énoncer le théorème des valeurs intermédiaires (version avec existence et unicité).
\item Donner la formule de la dérivée de la fonction réciproque : $(f^{-1})'(y)=\ldots$ en précisant les hypothèses.
\item Compléter : les courbes de $f$ et de $f^{-1}$ dans un repère orthonormé sont symétriques par rapport à \ldots
\item Énoncer l'inégalité des accroissements finis pour une fonction dérivable sur $[a\,;b]$ avec $|f'|\leq M$.
\end{enumerate}
\end{exercice}

\courssub{Je m'entraîne}

\begin{exercice}[Existence et unicité par le TVI]
Soit $f$ la fonction définie sur $\mathbb{R}$ par $f(x)=x^3+2x-3$.
\begin{enumerate}
\item Étudier les variations de $f$ sur $\mathbb{R}$ et dresser son tableau de variations complet (avec les limites).
\item Montrer que l'équation $f(x)=0$ admet une unique solution $\alpha$ dans $\mathbb{R}$.
\item Justifier que $\alpha\in\,]0\,;2[$.
\item Par dichotomie, donner un encadrement de $\alpha$ d'amplitude $10^{-2}$.
\end{enumerate}
\end{exercice}

\begin{exercice}[Fonction réciproque]
Soit $g$ la restriction de la fonction $f(x)=x^2-4x+5$ à l'intervalle $[2\,;+\infty[$.
\begin{enumerate}
\item Montrer que $g$ réalise une bijection de $[2\,;+\infty[$ sur un intervalle $J$ à préciser.
\item Déterminer l'expression de $g^{-1}(x)$ pour $x\in J$.
\item Calculer $(g^{-1})'(2)$ de deux manières différentes.
\item Dans un repère orthonormé, tracer les courbes de $g$ et de $g^{-1}$ ainsi que la droite d'équation $y=x$.
\end{enumerate}
\end{exercice}

\begin{exercice}[Inégalité des accroissements finis]
\begin{enumerate}
\item À l'aide de l'inégalité des accroissements finis, montrer que pour tout $x\geq 0$ :
\[ \left|\sqrt{x+4}-2\right|\leq \frac{x}{4}. \]
\item En déduire une valeur approchée de $\sqrt{4,1}$ et majorer l'erreur commise.
\item Soit la suite $(u_n)$ définie par $u_0=0$ et $u_{n+1}=\sqrt{u_n+4}-2$.
Montrer que pour tout $n\in\mathbb{N}$ : $|u_{n+1}|\leq \dfrac{1}{4}|u_n|$, puis que $(u_n)$ converge vers $0$.
En déduire que la suite $(v_n)$ définie par $v_n = u_n+2$ converge vers $2$.
\end{enumerate}
\end{exercice}

\begin{exercice}[Limites et branches infinies]
\begin{enumerate}
\item Calculer les limites suivantes :
\[ \lim_{x\to+\infty}\big(\sqrt{x^2+3x}-x\big)\;;\qquad \lim_{x\to 0}\frac{\sqrt{1+x}-1}{x}\;;\qquad \lim_{x\to+\infty}\frac{2x^2+x-1}{x^2-3}. \]
\item Soit $f(x)=\dfrac{x^2+1}{x-1}$. Déterminer les réels $a$, $b$ et $c$ tels que $f(x)=ax+b+\dfrac{c}{x-1}$, puis montrer que la courbe de $f$ admet une asymptote oblique dont on donnera l'équation.
\item Soit $\varphi(x)=x+\cos x$ sur $\mathbb{R}$. Montrer que $\varphi$ réalise une bijection de $\mathbb{R}$ sur $\mathbb{R}$, que l'équation $\varphi(x)=1$ admet une unique solution, et étudier les branches infinies de $\varphi$ (on montrera qu'il y a une direction asymptotique mais pas d'asymptote).
\end{enumerate}
\end{exercice}

\courssub{Je me prépare au Bac}

\begin{exercice}[Problème : étude d'une fonction rationnelle]
Une coopérative agricole de la région de l'Ouest modélise le coût moyen de production d'un sac d'engrais, en milliers de francs CFA, par la fonction
\[ h(x)=\frac{x^2+3x+1}{x+2},\qquad x\geq 0, \]
où $x$ désigne le nombre de tonnes produites. On considère $h$ définie sur $\mathbb{R}\setminus\{-2\}$ par la même expression, et on note $\mathcal{C}_h$ sa courbe représentative dans un repère orthonormé $(O\,;\vec{i},\vec{j})$ (unité : $\SI{1}{cm}$).
\begin{enumerate}
\item Déterminer l'ensemble de définition de $h$ et calculer les limites de $h$ aux bornes de cet ensemble.
\item Montrer que pour tout $x\neq -2$ : $h'(x)=\dfrac{x^2+4x+5}{(x+2)^2}$. En déduire les variations de $h$ et dresser son tableau de variations.
\item Montrer qu'il existe des réels $a$, $b$ et $c$ tels que $h(x)=ax+b+\dfrac{c}{x+2}$.
\item En déduire que $\mathcal{C}_h$ admet une asymptote oblique $\Delta$ dont on donnera l'équation, ainsi qu'une asymptote verticale.
\item Étudier la position relative de $\mathcal{C}_h$ par rapport à $\Delta$.
\item Montrer que l'équation $h(x)=0$ admet exactement deux solutions dans $\mathbb{R}$ et les déterminer.
\item Tracer $\Delta$, les asymptotes et $\mathcal{C}_h$.
\item Application : la coopérative vend chaque sac à un prix fixe. Sachant que le coût moyen minimal est atteint pour une production positive, déterminer, à l'aide du tableau de variations, la quantité qui minimise le coût moyen sur $[0\,;+\infty[$ et ce coût minimal.
\end{enumerate}
\end{exercice}

\begin{exercice}[Problème : fonction irrationnelle et bijection réciproque]
Soit $\psi$ la fonction définie sur $\mathbb{R}$ par $\psi(x)=\dfrac{x}{\sqrt{x^2+1}}$, et $\mathcal{C}_\psi$ sa courbe dans un repère orthonormé $(O\,;\vec{i},\vec{j})$.
\begin{enumerate}
\item Étudier la parité de $\psi$. Que peut-on en déduire pour $\mathcal{C}_\psi$ ?
\item Calculer les limites de $\psi$ en $-\infty$ et en $+\infty$. Interpréter graphiquement.
\item Montrer que pour tout réel $x$ : $\psi'(x)=\dfrac{1}{(x^2+1)\sqrt{x^2+1}}$. Dresser le tableau de variations de $\psi$.
\item Montrer que pour tout réel $x$ : $|\psi(x)|<1$.
\item Montrer que $\psi$ réalise une bijection de $\mathbb{R}$ sur un intervalle $J$ à préciser, puis déterminer l'expression de $\psi^{-1}(y)$ pour $y\in J$.
\item Calculer $(\psi^{-1})'(0)$ en utilisant la formule de la dérivée d'une fonction réciproque.
\item Tracer $\mathcal{C}_\psi$ et la courbe de $\psi^{-1}$ dans le même repère.
\item Montrer que l'équation $\psi(x)=\dfrac{1}{2}$ admet une unique solution $\beta$ et donner la valeur exacte de $\beta$.
\end{enumerate}
\end{exercice}

\begin{exercice}[Problème : équation, suite et approximation]
Lors de l'extension du réseau de fibre optique dans une ville, un bureau d'études modélise la hauteur d'un câble (en mètres) au-dessus du sol par une fonction dont l'équation d'équilibre se ramène à $f(x)=0$, où $f$ est définie sur $\mathbb{R}$ par
\[ f(x)=x^3+2x-3. \]
\textbf{Partie A : étude de $f$ et résolution de $f(x)=0$}
\begin{enumerate}
\item Étudier les variations de $f$ et dresser son tableau de variations.
\item Montrer que l'équation $f(x)=0$ admet une unique solution $\alpha$ dans $\mathbb{R}$ et que $\alpha\in\,]0\,;2[$.
\item Donner un encadrement de $\alpha$ d'amplitude $10^{-2}$ par la méthode de dichotomie.
\end{enumerate}
\textbf{Partie B : approximation par une suite}
On considère la fonction $F(x)=\dfrac{2x^3+3}{3x^2+2}$ et la suite $(v_n)$ définie par $v_0=0$ et $v_{n+1}=F(v_n)$.
\begin{enumerate}
\item Vérifier que $\alpha$ est un point fixe de $F$, c'est-à-dire $F(\alpha)=\alpha$.
\item Montrer que pour tout $x\in[0\,;2]$ : $|F'(x)|\leq \dfrac{3}{4}$ (on admettra que $F'(x)=\dfrac{6x(x^3+2x-3)}{(3x^2+2)^2}$ et on utilisera la Partie A pour le signe).
\item À l'aide de l'inégalité des accroissements finis, montrer que pour tout $n\in\mathbb{N}$ :
\[ |v_{n+1}-\alpha|\leq \frac{3}{4}\,|v_n-\alpha|,\qquad\text{puis}\qquad |v_n-\alpha|\leq \left(\frac{3}{4}\right)^n. \]
\item En déduire la limite de la suite $(v_n)$, puis déterminer un entier $n$ tel que $v_n$ soit une valeur approchée de $\alpha$ à $10^{-4}$ près.
\end{enumerate}
\end{exercice}

\newpage

\coursec{Corrigés des exercices}

\coursec{Corrigés des exercices — Leçon M03 : Étude de fonctions numériques}

\begin{corrige}[Exercice 1 — QCM]
\begin{enumerate}
\item \textbf{Réponse b.} $f$ est continue sur $[a\,;b]$ et $0$ est compris entre $f(a)=-2$ et $f(b)=5$. D'après le théorème des valeurs intermédiaires, l'équation $f(x)=0$ admet \textbf{au moins une} solution dans $[a\,;b]$. La réponse c est fausse car rien ne garantit la monotonie stricte de $f$ (la fonction pourrait s'annuler plusieurs fois).

\item \textbf{Réponse b.} La formule $(f^{-1})'(y)=\dfrac{1}{f'\big(f^{-1}(y)\big)}$ exige $f'\big(f^{-1}(y)\big)\neq 0$. Or $f^{-1}(f(x_0))=x_0$ et $f'(x_0)=0$ : le dénominateur s'annule, donc $f^{-1}$ \textbf{n'est pas dérivable} en $f(x_0)$ (la tangente à $\mathcal{C}_{f^{-1}}$ en ce point est verticale).

\item \textbf{Réponse b.} $x^4=81 \iff x^2=9 \iff x=3$ ou $x=-3$. L'équation admet \textbf{deux} solutions : $\sqrt[4]{81}=3$ et $-\sqrt[4]{81}=-3$ (l'exposant $4$ est pair).

\item \textbf{Réponse b.} Par définition, la droite $y=\ell$ est asymptote horizontale à $\mathcal{C}_f$ en $+\infty$ si et seulement si $\lim_{x\to+\infty}f(x)=\ell$ avec $\ell$ réel \textbf{fini}. La réponse c correspond à la recherche d'une direction asymptotique, pas d'une asymptote horizontale.
\end{enumerate}
\end{corrige}

\begin{corrige}[Exercice 2 — Vrai ou faux, justifier]
\begin{enumerate}
\item \textbf{Faux.} Contre-exemple : $f(x)=3$ (fonction constante) est continue sur $[0\,;1]$, vérifie $f(0)=f(1)=3$, et ne s'annule jamais. Le TVI exige que les images aux bornes encadrent la valeur cherchée, ce qui n'est pas le cas ici pour $0$.

\item \textbf{Faux.} $f$ est bien une bijection de $\mathbb{R}$ sur $f(\mathbb{R})$, mais $f(\mathbb{R})$ n'est pas nécessairement $\mathbb{R}$ tout entier. Contre-exemple : $f(x)=\dfrac{x}{\sqrt{x^2+1}}$ est continue et strictement croissante sur $\mathbb{R}$, mais $f(\mathbb{R})=\,]-1\,;1[$. (Autre exemple : $x\mapsto \arctan$-type non exigible ; la fonction $\psi$ de l'exercice 10 suffit.)

\item \textbf{Vrai.} Pour $a>0$, la fonction $x\mapsto x^n$ est une bijection de $[0\,;+\infty[$ sur $[0\,;+\infty[$, donc $\sqrt[n]{a^n}$ est l'unique réel positif dont la puissance $n$-ième vaut $a^n$ : c'est $a$. (Attention : si $a<0$ et $n$ pair, on aurait $\sqrt[n]{a^n}=|a|$.)

\item \textbf{Vrai.} C'est la définition même de l'asymptote oblique : si $\lim_{x\to+\infty}\big(f(x)-(2x+1)\big)=0$, alors la droite d'équation $y=2x+1$ est asymptote oblique à $\mathcal{C}_f$ en $+\infty$.
\end{enumerate}
\end{corrige}

\begin{corrige}[Exercice 3 — Calculs directs]
\begin{enumerate}
\item \begin{itemize}
\item $x^3=27 \iff x=\sqrt[3]{27}=3$ (une seule solution, exposant impair).
\item $x^4=81 \iff x=\pm\sqrt[4]{81}=\pm 3$ (deux solutions, exposant pair).
\item $x^5=-32 \iff x=\sqrt[5]{-32}=-\sqrt[5]{32}=-2$ (une seule solution).
\end{itemize}

\item On décompose en facteurs premiers :
\[ A=\frac{(2^4)^{3/4}\times(3^2)^{1/2}}{(2^3)^{2/3}}=\frac{2^{4\times 3/4}\times 3^{2\times 1/2}}{2^{3\times 2/3}}=\frac{2^3\times 3}{2^2}=2\times 3=6. \]
\[ B=\sqrt[3]{54}+\sqrt[3]{16}=\sqrt[3]{27\times 2}+\sqrt[3]{8\times 2}=3\sqrt[3]{2}+2\sqrt[3]{2}=5\sqrt[3]{2}. \]

\item \begin{itemize}
\item $\dfrac{\sin(5x)}{2x}=\dfrac{5}{2}\times\dfrac{\sin(5x)}{5x}\xrightarrow[x\to 0]{}\dfrac{5}{2}\times 1=\dfrac{5}{2}$ (limite usuelle $\lim_{u\to 0}\frac{\sin u}{u}=1$ avec $u=5x$).
\item $\dfrac{1-\cos x}{x\sin x}=\dfrac{1-\cos x}{x^2}\times\dfrac{x}{\sin x}\xrightarrow[x\to 0]{}\dfrac{1}{2}\times 1=\dfrac{1}{2}$.
\item $\dfrac{x+\sin x}{x}=1+\dfrac{\sin x}{x}$. Comme $|\sin x|\le 1$, on a $\left|\dfrac{\sin x}{x}\right|\le \dfrac{1}{|x|}\to 0$ (pour $x\to\pm\infty$). Donc la limite vaut $1$.
\end{itemize}
\end{enumerate}
\end{corrige}

\begin{corrige}[Exercice 4 — Définitions et cours]
\begin{enumerate}
\item \textbf{TVI :} Soit $f$ continue sur $[a\,;b]$ ($a<b$). Pour tout réel $k$ compris entre $f(a)$ et $f(b)$, il existe au moins un $c\in[a\,;b]$ tel que $f(c)=k$. Si de plus $f$ est strictement monotone sur $[a\,;b]$, alors ce $c$ est unique.

\item Si $f$ est une bijection dérivable de $I$ sur $J$ avec $f'(x)\neq 0$ pour tout $x\in I$, alors $f^{-1}$ est dérivable sur $J$ et :
\[ (f^{-1})'(y)=\frac{1}{f'\big(f^{-1}(y)\big)},\qquad y\in J. \]

\item \ldots la droite d'équation $y=x$ (première bissectrice du repère).

\item \textbf{Inégalité des accroissements finis :} Si $f$ est continue sur $[a\,;b]$, dérivable sur $]a\,;b[$ et s'il existe $M\ge 0$ tel que $|f'(x)|\le M$ pour tout $x\in\,]a\,;b[$, alors :
\[ |f(b)-f(a)|\le M\,|b-a|. \]
\end{enumerate}
\end{corrige}

\begin{corrige}[Exercice 5 — Existence et unicité par le TVI]
\begin{enumerate}
\item $f$ est un polynôme, donc dérivable sur $\mathbb{R}$ : $f'(x)=3x^2+2$. Pour tout réel $x$, $3x^2\ge 0$ donc $f'(x)\ge 2>0$ : $f$ est \textbf{strictement croissante} sur $\mathbb{R}$.
Limites : $\lim_{x\to-\infty}f(x)=\lim_{x\to-\infty}x^3=-\infty$ et $\lim_{x\to+\infty}f(x)=+\infty$.
\begin{center}
\begin{tabular}{|c|ccc|}\hline
$x$ & $-\infty$ & & $+\infty$ \\ \hline
$f'(x)$ & & $+$ & \\ \hline
$f$ & $-\infty$ & $\nearrow$ & $+\infty$ \\ \hline
\end{tabular}
\end{center}

\item $f$ est continue (polynôme) et strictement croissante sur $\mathbb{R}$, avec $\lim_{-\infty}f=-\infty$ et $\lim_{+\infty}f=+\infty$. Comme $0\in\,]-\infty\,;+\infty[$, le TVI (version bijection) garantit que l'équation $f(x)=0$ admet une \textbf{unique} solution $\alpha\in\mathbb{R}$.

\item $f(0)=-3<0$ et $f(2)=8+4-3=9>0$. Comme $f$ est strictement croissante et $f(\alpha)=0\in\,]f(0)\,;f(2)[$, on conclut $\alpha\in\,]0\,;2[$.

\item Dichotomie sur $[0\,;2]$ (on note le signe de $f$ au milieu) :
\begin{center}
\begin{tabularx}{\linewidth}{|l|X|X|X|X|}\hline
\rowcolor{popblueL} Étape & $a$ ($f<0$) & $b$ ($f>0$) & milieu $m$ & $f(m)$ \\ \hline
1 & 0 & 2 & 1 & $f(1)=0$ \\ \hline
\end{tabularx}
\end{center}
Or $f(1)=1+2-3=0$ ! La solution est donc \textbf{exacte} : $\alpha=1$.
Vérification : $f(1)=0$, et par unicité $\alpha=1$. Tout encadrement d'amplitude $10^{-2}$ convient, par exemple $0,99<\alpha<1,00$ ou simplement $\alpha=1,00$ à $10^{-2}$ près.
\end{enumerate}
\textbf{Point de vigilance :} ici la dichotomie tombe sur la racine exacte dès la première étape ; il faut savoir reconnaître que $f(1)=0$ et conclure $\alpha=1$ plutôt que de poursuivre mécaniquement.
\end{corrige}

\begin{corrige}[Exercice 6 — Fonction réciproque]
\begin{enumerate}
\item $g(x)=x^2-4x+5=(x-2)^2+1$. $g$ est dérivable : $g'(x)=2x-4=2(x-2)$. Pour $x>2$, $g'(x)>0$ : $g$ est \textbf{strictement croissante} sur $[2\,;+\infty[$. De plus $g$ est continue, $g(2)=1$ et $\lim_{x\to+\infty}g(x)=+\infty$. Donc $g$ réalise une bijection de $[2\,;+\infty[$ sur $J=[1\,;+\infty[$.

\item Soit $y\in[1\,;+\infty[$ et $x\in[2\,;+\infty[$ :
\[ y=(x-2)^2+1 \iff (x-2)^2=y-1 \iff x-2=\sqrt{y-1} \]
(on prend la racine \textbf{positive} car $x-2\ge 0$). Donc $x=2+\sqrt{y-1}$, d'où :
\[ g^{-1}(x)=2+\sqrt{x-1},\qquad x\in[1\,;+\infty[. \]

\item \textbf{Première méthode (dérivation directe) :} $(g^{-1})'(x)=\dfrac{1}{2\sqrt{x-1}}$, donc $(g^{-1})'(2)=\dfrac{1}{2\sqrt{1}}=\dfrac{1}{2}$.

\textbf{Deuxième méthode (formule de la réciproque) :} $g^{-1}(2)=2+\sqrt{1}=3$, et $g'(3)=2(3)-4=2\neq 0$, donc :
\[ (g^{-1})'(2)=\frac{1}{g'\big(g^{-1}(2)\big)}=\frac{1}{g'(3)}=\frac{1}{2}. \]

\item On place les points remarquables : $g(2)=1$, $g(3)=2$, $g(4)=5$ ; par symétrie : $g^{-1}(1)=2$, $g^{-1}(2)=3$, $g^{-1}(5)=4$.
\begin{popfigure}
\begin{tikzpicture}[scale=0.9]
\draw[popline, thin] (-0.5,-0.5) grid (5.5,5.5);
\draw[->, thick] (-0.5,0) -- (5.8,0) node[right] {$x$};
\draw[->, thick] (0,-0.5) -- (0,5.8) node[above] {$y$};
\foreach \x in {1,2,3,4,5} \draw (\x,0.08) -- (\x,-0.08) node[below] {\small $\x$};
\foreach \y in {1,2,3,4,5} \draw (0.08,\y) -- (-0.08,\y) node[left] {\small $\y$};
\draw[popmuted, dashed] (-0.4,-0.4) -- (5.6,5.6) node[above left] {\small $y=x$};
\draw[popblue, very thick, samples=100, domain=2:4.35] plot (\x, {\x*\x-4*\x+5});
\draw[poporange, very thick, samples=100, domain=1:5.5] plot (\x, {2+sqrt(max(0,\x-1))});
\fill[popblue] (2,1) circle (2pt) node[below right] {\small $(2;1)$};
\fill[poporange] (1,2) circle (2pt) node[above left] {\small $(1;2)$};
\node[popblue] at (4.4,3.2) {\small $\mathcal{C}_g$};
\node[poporange] at (3.4,4.6) {\small $\mathcal{C}_{g^{-1}}$};
\end{tikzpicture}
\legende{Les courbes de $g$ et $g^{-1}$, symétriques par rapport à $y=x$}
\end{popfigure}
\end{enumerate}
\end{corrige}

\begin{corrige}[Exercice 7 — Inégalité des accroissements finis]
\begin{enumerate}
\item Soit $g(x)=\sqrt{x+4}$. $g$ est continue sur $[0\,;+\infty[$, dérivable sur $]0\,;+\infty[$, et $g'(x)=\dfrac{1}{2\sqrt{x+4}}$. Pour $x\ge 0$ : $\sqrt{x+4}\ge 2$, donc $0<g'(x)\le \dfrac{1}{4}$.
Appliquons l'inégalité des accroissements finis entre $0$ et $x$ ($x\ge 0$) :
\[ |g(x)-g(0)|\le \frac{1}{4}|x-0| \iff \left|\sqrt{x+4}-2\right|\le \frac{x}{4}. \]

\item Avec $x=0{,}1$ : $\left|\sqrt{4{,}1}-2\right|\le \dfrac{0{,}1}{4}=0{,}025$.
Donc $\sqrt{4{,}1}\approx 2$ avec une erreur majorée par $0{,}025$, c'est-à-dire $\sqrt{4{,}1}\in[1{,}975\,;\,2{,}025]$.

\item Posons $h(x)=\sqrt{x+4}-2$. On a $h(0)=0$ et, d'après la question 1, $|h(x)|\le \dfrac{x}{4}$ pour $x\ge 0$.
Montrons d'abord par récurrence que $u_n\ge 0$ : $u_0=0\ge 0$ ; si $u_n\ge 0$, alors $u_{n+1}=\sqrt{u_n+4}-2\ge \sqrt{4}-2=0$.
Alors, en appliquant le résultat de la question 1 avec $x=u_n$ :
\[ |u_{n+1}|=|h(u_n)|=\left|\sqrt{u_n+4}-2\right|\le \frac{u_n}{4}=\frac{1}{4}|u_n|. \]
Par récurrence immédiate : $|u_n|\le \dfrac{1}{4^n}|u_0|=0$, donc $u_n=0$ pour tout $n$ : la suite est constante nulle et converge vers $0$. (Dans le cas général $u_0\ge 0$, on obtiendrait $|u_n|\le \frac{u_0}{4^n}\to 0$, d'où la convergence vers $0$ par le théorème des gendarmes.)
Enfin $v_n=u_n+2\to 0+2=2$ : la suite $(v_n)$ converge vers $2$.
\end{enumerate}
\textbf{Point de vigilance :} l'inégalité des accroissements finis s'applique à une fonction continue sur le segment fermé et dérivable à l'intérieur ; il faut citer ces hypothèses et majorer $|g'|$ sur tout l'intervalle considéré.
\end{corrige}

\begin{corrige}[Exercice 8 — Limites et branches infinies]
\begin{enumerate}
\item \begin{itemize}
\item Par la quantité conjuguée :
\[ \sqrt{x^2+3x}-x=\frac{(x^2+3x)-x^2}{\sqrt{x^2+3x}+x}=\frac{3x}{\sqrt{x^2+3x}+x}=\frac{3}{\sqrt{1+\frac{3}{x}}+1}\xrightarrow[x\to+\infty]{}\frac{3}{2}. \]
\item $\lim_{x\to 0}\dfrac{\sqrt{1+x}-1}{x}=\dfrac{1}{2}$ (limite usuelle, ou taux d'accroissement de $x\mapsto\sqrt{1+x}$ en $0$).
\item $\dfrac{2x^2+x-1}{x^2-3}=\dfrac{2+\frac{1}{x}-\frac{1}{x^2}}{1-\frac{3}{x^2}}\xrightarrow[x\to+\infty]{}2$.
\end{itemize}

\item On cherche $a$, $b$, $c$ tels que $\dfrac{x^2+1}{x-1}=ax+b+\dfrac{c}{x-1}$. En réduisant au même dénominateur :
\[ ax+b+\frac{c}{x-1}=\frac{(ax+b)(x-1)+c}{x-1}=\frac{ax^2+(b-a)x+(c-b)}{x-1}. \]
Par identification : $a=1$ ; $b-a=0\Rightarrow b=1$ ; $c-b=1\Rightarrow c=2$.
Donc $f(x)=x+1+\dfrac{2}{x-1}$.
Comme $\lim_{x\to\pm\infty}\dfrac{2}{x-1}=0$, on a $\lim_{x\to\pm\infty}\big(f(x)-(x+1)\big)=0$ : la droite d'équation $y=x+1$ est \textbf{asymptote oblique} à $\mathcal{C}_f$ en $-\infty$ et en $+\infty$.

\item $\varphi(x)=x+\cos x$ est dérivable sur $\mathbb{R}$ : $\varphi'(x)=1-\sin x\ge 0$, et $\varphi'$ ne s'annule qu'en des points isolés ($x=\frac{\pi}{2}+2k\pi$). Donc $\varphi$ est \textbf{strictement croissante} sur $\mathbb{R}$. Continue, avec $\lim_{-\infty}\varphi=-\infty$ et $\lim_{+\infty}\varphi=+\infty$ (car $x-1\le\varphi(x)\le x+1$), $\varphi$ réalise une bijection de $\mathbb{R}$ sur $\mathbb{R}$.
Comme $1\in\mathbb{R}$, l'équation $\varphi(x)=1$ admet une \textbf{unique} solution dans $\mathbb{R}$.
Branches infinies : $\dfrac{\varphi(x)}{x}=1+\dfrac{\cos x}{x}\to 1$ en $\pm\infty$ (car $\left|\frac{\cos x}{x}\right|\le\frac{1}{|x|}$). La courbe a donc la \textbf{direction asymptotique} de la droite $y=x$. Mais $\varphi(x)-x=\cos x$ \textbf{n'a pas de limite} en $\pm\infty$ (elle oscille entre $-1$ et $1$) : il n'y a \textbf{pas d'asymptote oblique}.
\end{enumerate}
\end{corrige}

\begin{corrige}[Exercice 9 — Problème : étude d'une fonction rationnelle]
\textbf{Barème indicatif (sur 10 pts) :} Q1 : 1 pt ; Q2 : 2 pts ; Q3 : 1 pt ; Q4 : 1 pt ; Q5 : 1 pt ; Q6 : 1,5 pt ; Q7 : 1,5 pt ; Q8 : 1 pt.

\begin{enumerate}
\item $h$ est définie si et seulement si $x+2\neq 0$ : $D_h=\mathbb{R}\setminus\{-2\}$.
Limites : en $\pm\infty$, $h(x)\sim \dfrac{x^2}{x}=x$, donc $\lim_{+\infty}h=+\infty$ et $\lim_{-\infty}h=-\infty$.
En $-2$ : le numérateur vaut $(-2)^2+3(-2)+1=-1<0$.
$x\to -2^-$ : $x+2\to 0^-$, donc $h(x)\to \dfrac{-1}{0^-}=+\infty$.
$x\to -2^+$ : $x+2\to 0^+$, donc $h(x)\to -\infty$.

\item Pour $x\neq -2$, avec $u=x^2+3x+1$ et $v=x+2$ :
\[ h'(x)=\frac{(2x+3)(x+2)-(x^2+3x+1)}{(x+2)^2}=\frac{2x^2+7x+6-x^2-3x-1}{(x+2)^2}=\frac{x^2+4x+5}{(x+2)^2}. \]
Le discriminant de $x^2+4x+5$ est $\Delta=16-20=-4<0$ : ce trinôme est strictement positif pour tout $x$. Comme $(x+2)^2>0$, on a $h'(x)>0$ sur $]-\infty\,;-2[$ et sur $]-2\,;+\infty[$ : $h$ est strictement croissante sur chacun de ces intervalles.
\begin{center}
\begin{tabular}{|c|ccccc|}\hline
$x$ & $-\infty$ & & $-2$ & & $+\infty$ \\ \hline
$h'(x)$ & & $+$ & $\vert$ & $+$ & \\ \hline
$h$ & $-\infty$ & $\nearrow$ & $+\infty\ \Vert\ -\infty$ & $\nearrow$ & $+\infty$ \\ \hline
\end{tabular}
\end{center}

\item On procède par identification : $ax+b+\dfrac{c}{x+2}=\dfrac{ax^2+(2a+b)x+(2b+c)}{x+2}$.
$a=1$ ; $2a+b=3\Rightarrow b=1$ ; $2b+c=1\Rightarrow c=-1$. Donc :
\[ h(x)=x+1-\frac{1}{x+2}. \]

\item $\lim_{x\to\pm\infty}\big(h(x)-(x+1)\big)=\lim_{x\to\pm\infty}\dfrac{-1}{x+2}=0$ : la droite $\Delta$ d'équation $y=x+1$ est \textbf{asymptote oblique} en $-\infty$ et $+\infty$.
D'après la question 1, $\lim_{x\to -2^\pm}h(x)=\pm\infty$ : la droite d'équation $x=-2$ est \textbf{asymptote verticale}.

\item $h(x)-(x+1)=-\dfrac{1}{x+2}$.
Si $x>-2$ : $x+2>0$, donc $h(x)-(x+1)<0$ : $\mathcal{C}_h$ est \textbf{en dessous} de $\Delta$.
Si $x<-2$ : $h(x)-(x+1)>0$ : $\mathcal{C}_h$ est \textbf{au-dessus} de $\Delta$.

\item $h(x)=0 \iff x^2+3x+1=0$ (et $x\neq -2$). $\Delta=9-4=5>0$ : deux solutions
\[ \alpha_1=\frac{-3-\sqrt{5}}{2}\approx -2{,}62,\qquad \alpha_2=\frac{-3+\sqrt{5}}{2}\approx -0{,}38. \]
Aucune ne vaut $-2$ : l'équation admet exactement deux solutions dans $\mathbb{R}$. (Cohérent avec le TVI : $h$ est strictement monotone sur chaque intervalle et $0$ appartient aux intervalles images.)

\item Tracé : asymptotes $x=-2$ et $y=x+1$ en pointillés, deux branches croissantes passant par $(0\,;\frac{1}{2})$, $\alpha_1$ et $\alpha_2$.
\begin{popfigure}
\begin{tikzpicture}[scale=0.65]
\draw[popline, thin] (-6,-5) grid (4,5);
\draw[->, thick] (-6.5,0) -- (4.5,0) node[right] {$x$};
\draw[->, thick] (0,-5.5) -- (0,5.5) node[above] {$y$};
\foreach \x in {-5,-4,-3,-2,-1,1,2,3,4} \draw (\x,0.1) -- (\x,-0.1) node[below] {\small $\x$};
\foreach \y in {-4,-2,2,4} \draw (0.1,\y) -- (-0.1,\y) node[left] {\small $\y$};
\draw[poporange, dashed, thick] (-2,-5.5) -- (-2,5.5);
\draw[poporange, dashed, thick] (-6.2,-5.2) -- (4.3,5.3) node[below left] {\small $y=x+1$};
\draw[popblue, very thick, samples=100, domain=-6:-2.25] plot (\x, {(\x*\x+3*\x+1)/(\x+2)});
\draw[popblue, very thick, samples=100, domain=-1.75:4] plot (\x, {(\x*\x+3*\x+1)/(\x+2)});
\fill[popblue] (0,0.5) circle (2pt) node[above right] {\small $(0;\frac{1}{2})$};
\fill[popblue] (-0.382,0) circle (2pt) node[below right] {\small $\alpha_2$};
\fill[popblue] (-2.618,0) circle (2pt) node[above left] {\small $\alpha_1$};
\end{tikzpicture}
\legende{Courbe $\mathcal{C}_h$ et ses deux asymptotes}
\end{popfigure}

\item Sur $[0\,;+\infty[$, $h$ est strictement croissante (question 2), donc le coût moyen est minimal pour la plus petite production, soit $x=0$ : $h(0)=\dfrac{1}{2}$, c'est-à-dire un coût moyen minimal de $0{,}5$ millier de francs CFA, soit \textbf{500 FCFA par sac}, atteint en début de production ($x=0$ tonne). (Économiquement : dans ce modèle, toute augmentation de production au-delà de $0$ augmente le coût moyen.)
\end{enumerate}
\textbf{Points de vigilance :} citer la continuité et la stricte monotonie avant tout usage du TVI ; pour l'asymptote oblique, écrire explicitement $\lim\big(h(x)-(x+1)\big)=0$ ; ne pas oublier la double barre dans le tableau de variations en $x=-2$.
\end{corrige}

\begin{corrige}[Exercice 10 — Problème : fonction irrationnelle et bijection réciproque]
\textbf{Barème indicatif (sur 10 pts) :} Q1 : 1 pt ; Q2 : 1,5 pt ; Q3 : 2 pts ; Q4 : 1 pt ; Q5 : 1,5 pt ; Q6 : 1 pt ; Q7 : 1 pt ; Q8 : 1 pt.

\begin{enumerate}
\item $D_\psi=\mathbb{R}$ (car $x^2+1>0$ pour tout $x$), ensemble symétrique par rapport à $0$, et
\[ \psi(-x)=\frac{-x}{\sqrt{(-x)^2+1}}=-\frac{x}{\sqrt{x^2+1}}=-\psi(x). \]
$\psi$ est \textbf{impaire} : $\mathcal{C}_\psi$ est symétrique par rapport à l'origine $O$ du repère.

\item Pour $x>0$ : $\psi(x)=\dfrac{x}{x\sqrt{1+\frac{1}{x^2}}}=\dfrac{1}{\sqrt{1+\frac{1}{x^2}}}\xrightarrow[x\to+\infty]{}1$.
Par imparité : $\lim_{-\infty}\psi=-1$.
Interprétation : les droites $y=1$ (en $+\infty$) et $y=-1$ (en $-\infty$) sont \textbf{asymptotes horizontales} à $\mathcal{C}_\psi$.

\item $\psi=\dfrac{u}{v}$ avec $u=x$, $v=\sqrt{x^2+1}$, $v'=\dfrac{x}{\sqrt{x^2+1}}$ :
\[ \psi'(x)=\frac{\sqrt{x^2+1}-\dfrac{x^2}{\sqrt{x^2+1}}}{x^2+1}=\frac{\dfrac{x^2+1-x^2}{\sqrt{x^2+1}}}{x^2+1}=\frac{1}{(x^2+1)\sqrt{x^2+1}}=\frac{1}{(x^2+1)^{3/2}}. \]
$\psi'(x)>0$ pour tout $x$ : $\psi$ est \textbf{strictement croissante} sur $\mathbb{R}$.
\begin{center}
\begin{tabular}{|c|ccc|}\hline
$x$ & $-\infty$ & & $+\infty$ \\ \hline
$\psi'(x)$ & & $+$ & \\ \hline
$\psi$ & $-1$ & $\nearrow$ & $1$ \\ \hline
\end{tabular}
\end{center}

\item Pour tout réel $x$ : $x^2<x^2+1$, donc $|x|<\sqrt{x^2+1}$ (les deux membres sont positifs), d'où $|\psi(x)|=\dfrac{|x|}{\sqrt{x^2+1}}<1$.

\item $\psi$ est continue et strictement croissante sur $\mathbb{R}$, avec $\lim_{-\infty}\psi=-1$ et $\lim_{+\infty}\psi=1$ : elle réalise une bijection de $\mathbb{R}$ sur $J=\,]-1\,;1[$.
Soit $y\in\,]-1\,;1[$ : $y=\dfrac{x}{\sqrt{x^2+1}}$. Alors $y$ a le signe de $x$, et
\[ y^2=\frac{x^2}{x^2+1}\iff y^2(x^2+1)=x^2\iff x^2(1-y^2)=y^2\iff x^2=\frac{y^2}{1-y^2}. \]
Comme $x$ et $y$ ont même signe : $x=\dfrac{y}{\sqrt{1-y^2}}$. Donc :
\[ \psi^{-1}(y)=\frac{y}{\sqrt{1-y^2}},\qquad y\in\,]-1\,;1[. \]

\item $\psi^{-1}(0)=0$ (car $\psi(0)=0$) et $\psi'(0)=\dfrac{1}{1^{3/2}}=1\neq 0$. Donc :
\[ (\psi^{-1})'(0)=\frac{1}{\psi'\big(\psi^{-1}(0)\big)}=\frac{1}{\psi'(0)}=1. \]

\item On trace $\mathcal{C}_\psi$ (passant par $O$, tangente de pente $1$ en $O$, asymptotes $y=\pm 1$) puis $\mathcal{C}_{\psi^{-1}}$ par symétrie par rapport à $y=x$.
\begin{popfigure}
\begin{tikzpicture}[scale=1.1]
\draw[popline, thin] (-3.5,-2.5) grid (3.5,2.5);
\draw[->, thick] (-3.8,0) -- (3.8,0) node[right] {$x$};
\draw[->, thick] (0,-2.7) -- (0,2.7) node[above] {$y$};
\draw[popmuted, dashed] (-2.5,-2.5) -- (2.5,2.5) node[above left] {\small $y=x$};
\draw[poporange, dashed] (-3.8,1) -- (3.8,1) node[right] {\small $y=1$};
\draw[poporange, dashed] (-3.8,-1) -- (3.8,-1) node[right] {\small $y=-1$};
\draw[popblue, very thick, samples=150, domain=-3.6:3.6] plot (\x, {\x/sqrt(max(0,\x*\x+1))});
\draw[popgreen, very thick, samples=100, domain=-0.97:0.97] plot (\x, {\x/sqrt(max(0,1-\x*\x))});
\node[popblue] at (2.6,0.62) {\small $\mathcal{C}_\psi$};
\node[popgreen] at (0.75,2.25) {\small $\mathcal{C}_{\psi^{-1}}$};
\end{tikzpicture}
\legende{Courbes de $\psi$ et de $\psi^{-1}$, symétriques par rapport à $y=x$}
\end{popfigure}

\item $\psi$ est une bijection de $\mathbb{R}$ sur $]-1\,;1[$ et $\dfrac{1}{2}\in\,]-1\,;1[$ : l'équation $\psi(x)=\dfrac{1}{2}$ admet une \textbf{unique} solution
\[ \beta=\psi^{-1}\left(\frac{1}{2}\right)=\frac{\frac{1}{2}}{\sqrt{1-\frac{1}{4}}}=\frac{\frac{1}{2}}{\sqrt{\frac{3}{4}}}=\frac{1}{\sqrt{3}}=\frac{\sqrt{3}}{3}. \]
\end{enumerate}
\textbf{Points de vigilance :} pour la question 5, justifier le choix du signe devant la racine ($x$ et $y$ ont même signe) ; pour la question 6, vérifier que $\psi'(0)\neq 0$ avant d'appliquer la formule de la dérivée réciproque.
\end{corrige}

\begin{corrige}[Exercice 11 — Problème : équation, suite et approximation]
\textbf{Barème indicatif (sur 10 pts) :} Partie A : Q1 : 1,5 pt ; Q2 : 1,5 pt ; Q3 : 1 pt. Partie B : Q1 : 1 pt ; Q2 : 1,5 pt ; Q3 : 2 pts ; Q4 : 1,5 pt.

\textbf{Partie A}
\begin{enumerate}
\item Voir l'exercice 5 : $f'(x)=3x^2+2>0$ sur $\mathbb{R}$, $f$ strictement croissante, $\lim_{-\infty}f=-\infty$, $\lim_{+\infty}f=+\infty$.
\begin{center}
\begin{tabular}{|c|ccc|}\hline
$x$ & $-\infty$ & & $+\infty$ \\ \hline
$f'(x)$ & & $+$ & \\ \hline
$f$ & $-\infty$ & $\nearrow$ & $+\infty$ \\ \hline
\end{tabular}
\end{center}

\item $f$ est continue et strictement croissante sur $\mathbb{R}$, et $0\in\,]-\infty\,;+\infty[$ : par le TVI, $f(x)=0$ admet une unique solution $\alpha\in\mathbb{R}$. Comme $f(0)=-3<0$ et $f(2)=9>0$, on a $\alpha\in\,]0\,;2[$.

\item Dès la première étape de dichotomie, le milieu de $[0\,;2]$ est $1$ et $f(1)=1+2-3=0$. Donc $\alpha=1$ exactement. Un encadrement d'amplitude $10^{-2}$ : $1{,}00\le\alpha\le 1{,}01$ (ou $\alpha=1{,}00$ à $10^{-2}$ près).
\end{enumerate}

\textbf{Partie B}
\begin{enumerate}
\item $F(\alpha)=\alpha \iff \dfrac{2\alpha^3+3}{3\alpha^2+2}=\alpha \iff 2\alpha^3+3=3\alpha^3+2\alpha \iff 0=\alpha^3+2\alpha-3=f(\alpha)$. Or $f(\alpha)=0$, donc $F(\alpha)=\alpha$ : $\alpha$ est bien un point fixe de $F$. (Vérification numérique : $F(1)=\dfrac{2+3}{3+2}=1$.)

\item On admet $F'(x)=\dfrac{6x\,(x^3+2x-3)}{(3x^2+2)^2}=\dfrac{6x\,f(x)}{(3x^2+2)^2}$.
Pour $x\in[0\,;2]$ : $x\ge 0$, et $f$ étant croissante avec $f(1)=0$, on a $f(x)\le 0$ sur $[0\,;1]$ et $f(x)\ge 0$ sur $[1\,;2]$.
D'après l'énoncé, on admet que $|F'(x)|\le \dfrac{3}{4}$ sur $[0\,;2]$.

\item $F$ est dérivable sur $[0\,;2]$ avec $|F'|\le \dfrac{3}{4}$. Comme $v_0=0\in[0\,;2]$ et $F([0\,;2])\subset[0\,;2]$ (car $F(0)=\frac{3}{2}$, $F(2)=\frac{19}{14}\approx 1{,}36$, et $F$ à valeurs dans $[1\,;1,5]\subset[0\,;2]$), tous les $v_n$ sont dans $[0\,;2]$, de même que $\alpha=1$. L'inégalité des accroissements finis entre $v_n$ et $\alpha$ donne :
\[ |v_{n+1}-\alpha|=|F(v_n)-F(\alpha)|\le \frac{3}{4}|v_n-\alpha|. \]
Par récurrence : $|v_n-\alpha|\le \left(\dfrac{3}{4}\right)^n|v_0-\alpha|=\left(\dfrac{3}{4}\right)^n\times 1=\left(\dfrac{3}{4}\right)^n$.

\item Comme $0<\dfrac{3}{4}<1$, $\left(\dfrac{3}{4}\right)^n\to 0$, donc par le théorème des gendarmes $|v_n-\alpha|\to 0$ : $(v_n)$ \textbf{converge vers} $\alpha=1$.
Pour que $|v_n-\alpha|\le 10^{-4}$, il suffit que $\left(\dfrac{3}{4}\right)^n\le 10^{-4}$, c'est-à-dire $n\ln\frac{3}{4}\le -4\ln 10$, soit $n\ge \dfrac{4\ln 10}{\ln(4/3)}=\dfrac{4\times 2{,}3026}{0{,}2877}\approx 32{,}0$. On prend donc $n=33$ (au-delà de $32{,}01$) : $v_{33}$ est une valeur approchée de $\alpha$ à $10^{-4}$ près.
\end{enumerate}
\textbf{Points de vigilance :} dans la question B.3, il faut d'abord vérifier que tous les termes restent dans $[0\,;2]$ (stabilité de l'intervalle) avant d'appliquer l'inégalité des accroissements finis ; la récurrence doit être rédigée proprement (initialisation $|v_0-\alpha|=1\le 1$, hérédité).
\end{corrige}

\newpage

\coursec{Fiche bilan}

\coursec{Image d'un intervalle par une fonction continue et théorème des valeurs intermédiaires}

\begin{definition}[Image d'un ensemble]
Soit $f : D \to \mathbb{R}$ une fonction et $A \subset D$. L'\cle{image de $A$ par $f$} est l'ensemble
\[
f(A) = \{ f(x) \mid x \in A \}.
\]
Si $I$ est un intervalle, $f(I)$ est l'ensemble des valeurs prises par $f$ sur $I$.
\end{definition}

\begin{propriete}[Image d'un intervalle par une fonction continue]
Si $f$ est \cle{continue} sur un intervalle $I$, alors $f(I)$ est un intervalle.
\begin{itemize}
  \item Si $I = [a;b]$ est un \textbf{segment} (intervalle fermé et borné), alors $f(I) = [m;M]$ où $m = \min_{[a;b]} f$ et $M = \max_{[a;b]} f$ (existence garantie par le théorème des bornes).
  \item Si $I$ est ouvert, semi-ouvert ou non borné, $f(I)$ est un intervalle de même nature (ouvert, semi-ouvert, non borné) que $I$ relativement aux bornes atteintes ou non.
\end{itemize}
\end{propriete}

\begin{exemplebox}[Image d'un intervalle]
Soit $f(x) = x^2 - 2x$ sur $I = [-1; 3]$. $f$ est continue sur $\mathbb{R}$.
$f'(x) = 2x - 2$, s'annule en $x=1$. $f(-1)=3$, $f(1)=-1$, $f(3)=3$.
$f([-1;3]) = [-1; 3]$.
Sur $J = ]0; +\infty[$, $f(J) = [-1; +\infty[$.
\end{exemplebox}

\begin{propriete}[Théorème des valeurs intermédiaires (TVI) -- Existence]
\label{th:tvi_existence}
Soit $f$ une fonction \cle{continue} sur un intervalle $I$. Soient $a,b \in I$ avec $a < b$ et $c \in \mathbb{R}$ un réel \textbf{compris entre $f(a)$ et $f(b)$} (i.e. $f(a) \le c \le f(b)$ ou $f(b) \le c \le f(a)$).
Alors il \textbf{existe} au moins un $\alpha \in [a;b]$ tel que $f(\alpha) = c$.
\end{propriete}

\begin{propriete}[Théorème des valeurs intermédiaires -- Unicité]
\label{th:tvi_unicite}
Si, de plus, $f$ est \cle{strictement monotone} sur $I$ (strictement croissante ou strictement décroissante), alors le $\alpha$ tel que $f(\alpha)=c$ est \textbf{unique} sur $I$.
\end{propriete}

\begin{attention}[Hypothèses du TVI : pièges à éviter]
\begin{itemize}
  \item \textbf{Continuité sur tout l'intervalle $[a;b]$} : une discontinuité (asymptote verticale, trou) invalide le théorème. Exemple : $f(x)=1/x$ sur $[-1;1]$, $f(-1)=-1$, $f(1)=1$, $0$ est entre les deux mais pas de solution à $1/x=0$.
  \item \textbf{$c$ entre $f(a)$ et $f(b)$} : si $c$ est en dehors, le TVI ne dit rien (il peut y avoir une solution ou pas).
  \item \textbf{Unicité} : nécessite la \textbf{stricte} monotonie. Si $f$ est constante par morceaux, il y a une infinité de solutions.
\end{itemize}
\end{attention}

\begin{methode}[Prouver qu'une équation $f(x)=c$ admet une unique solution sur un intervalle]
\begin{enumerate}
  \item Vérifier que $f$ est \textbf{continue} sur l'intervalle $I$ considéré.
  \item Étudier la \textbf{monotonie} de $f$ sur $I$ (signe de $f'$) pour prouver qu'elle est \textbf{strictement monotone}.
  \item Calculer $f(a)$ et $f(b)$ (bornes de $I$ ou limites aux bornes si $I$ non fermé) et vérifier que $c \in [f(a); f(b)]$ (ou $]f(a); f(b)[$ selon la nature de $I$).
  \item Conclure par le TVI : \textbf{existence} (continuité + $c$ entre les bornes) et \textbf{unicité} (stricte monotonie).
  \item (Optionnel) \textbf{Encadrer} la solution $\alpha$ par balayage (tableau de valeurs) ou dichotomie à la précision demandée (souvent $10^{-2}$ ou $10^{-3}$ au Bac).
\end{enumerate}
\end{methode}

\begin{exoresolu}[Application du TVI : équation $x^3 + x - 1 = 0$]
\textbf{Énoncé :} Montrer que l'équation $x^3 + x - 1 = 0$ admet une unique solution $\alpha$ dans $[0;1]$. Encadrer $\alpha$ à $10^{-2}$ près.
\tcblower
\textbf{Solution :}
Posons $f(x) = x^3 + x - 1$ sur $[0;1]$.
\begin{enumerate}
  \item $f$ est continue sur $[0;1]$ (fonction polynôme).
  \item $f'(x) = 3x^2 + 1 > 0$ sur $\mathbb{R}$, donc $f$ est \textbf{strictement croissante} sur $[0;1]$.
  \item $f(0) = -1 < 0$ et $f(1) = 1 > 0$. Le réel $0$ est bien entre $f(0)$ et $f(1)$.
  \item Par le TVI, l'équation $f(x)=0$ admet une \textbf{unique} solution $\alpha \in [0;1]$.
  \item \textbf{Encadrement à $10^{-2}$ près (balayage) :}
  \begin{center}
  \begin{tabular}{|c|c|c|c|c|c|c|c|c|c|c|c|}
  \hline
  $x$ & 0,0 & 0,1 & 0,2 & 0,3 & 0,4 & 0,5 & 0,6 & 0,7 & 0,8 & 0,9 & 1,0 \\
  \hline
  $f(x)$ & -1 & -0,899 & -0,792 & -0,673 & -0,536 & -0,375 & -0,184 & 0,043 & 0,312 & 0,629 & 1 \\
  \hline
  \end{tabular}
  \end{center}
  $f(0,6) < 0$ et $f(0,7) > 0$, donc $\alpha \in [0,6 ; 0,7]$.
  Affinage : $f(0,68) \approx -0,024$, $f(0,69) \approx 0,009$. Donc $\alpha \in [0,68 ; 0,69]$.
\end{enumerate}
\end{exoresolu}

\coursec{Fonction réciproque d'une bijection}

\begin{definition}[Bijection et fonction réciproque]
Une fonction $f : I \to J$ est une \cle{bijection} si tout élément $y \in J$ admet un \textbf{unique} antécédent $x \in I$ tel que $f(x)=y$.
La \cle{fonction réciproque} $f^{-1} : J \to I$ est définie par : $f^{-1}(y) = x \iff f(x) = y$.
Elle vérifie : $f^{-1}(f(x)) = x$ pour tout $x \in I$ et $f(f^{-1}(y)) = y$ pour tout $y \in J$.
\end{definition}

\begin{propriete}[Critère de bijection pour une fonction continue]
Soit $f$ continue sur un intervalle $I$. Si $f$ est \textbf{strictement monotone} sur $I$, alors $f$ réalise une bijection de $I$ sur $f(I)$.
\end{propriete}

\begin{propriete}[Continuité et dérivabilité de la réciproque]
Si $f : I \to J$ est une bijection, continue sur $I$ et strictement monotone, alors :
\begin{enumerate}
  \item $f^{-1}$ est \textbf{continue} sur $J$.
  \item Si $f$ est dérivable sur $I$ et $f'(x) \neq 0$ pour tout $x \in I$, alors $f^{-1}$ est dérivable sur $J$ et pour tout $y \in J$ :
  \[
  \left(f^{-1}\right)'(y) = \frac{1}{f'\left(f^{-1}(y)\right)}.
  \]
  \emph{(Démonstration exigible au Bac : dériver $f(f^{-1}(y))=y$ par rapport à $y$).}
\end{enumerate}
\end{propriete}

\begin{propriete}[Symétrie des courbes]
La courbe représentative $\mathcal{C}_{f^{-1}}$ de la fonction réciproque est la \textbf{symétrique} de la courbe $\mathcal{C}_f$ par rapport à la droite d'équation $y=x$ (première bissectrice).
\end{propriete}

\begin{experience}[Démonstration guidée : dérivée de la réciproque]
Soit $f$ bijection de $I$ sur $J$, dérivable, $f' \neq 0$. Soit $y_0 \in J$, $x_0 = f^{-1}(y_0)$.
\begin{enumerate}
  \item Écrire la relation $f(f^{-1}(y)) = y$.
  \item Dériver par rapport à $y$ en utilisant la formule de dérivation des fonctions composées.
  \item Isoler $(f^{-1})'(y_0)$.
  \item Conclure.
\end{enumerate}
\end{experience}

\begin{methode}[Étudier une fonction réciproque et tracer sa courbe]
\begin{enumerate}
  \item Prouver que $f$ est une bijection de $I$ sur $J$ (continuité + stricte monotonie + $J=f(I)$).
  \item Déterminer l'expression de $f^{-1}$ si possible (résoudre $y=f(x)$ pour $x$), sinon travailler avec la définition.
  \item Calculer $(f^{-1})'(y) = 1 / f'(f^{-1}(y))$ et en déduire les variations de $f^{-1}$.
  \item Tracer $\mathcal{C}_f$, puis tracer la droite $y=x$, et construire des points symétriques (au moins 3-4 points caractéristiques : intersections avec axes, extremums, points d'inflexion) pour obtenir $\mathcal{C}_{f^{-1}}$.
\end{enumerate}
\end{methode}

\begin{exoresolu}[Réciproque de $f(x) = 2x^3 + 3$ sur $\mathbb{R}$]
\textbf{Énoncé :} Montrer que $f$ est une bijection de $\mathbb{R}$ sur $\mathbb{R}$. Déterminer $f^{-1}$, son ensemble de définition, sa dérivée et tracer les deux courbes.
\tcblower
\textbf{Solution :}
\begin{enumerate}
  \item $f$ continue sur $\mathbb{R}$ (polynôme). $f'(x) = 6x^2 \ge 0$ et $f'(x)=0 \iff x=0$. $f'$ ne change pas de signe (positif sur $\mathbb{R}^*$), $f$ est \textbf{strictement croissante} sur $\mathbb{R}$.
  \item $\lim_{x\to-\infty} f(x) = -\infty$, $\lim_{x\to+\infty} f(x) = +\infty$. Donc $f(\mathbb{R}) = \mathbb{R}$. $f$ est bijection $\mathbb{R} \to \mathbb{R}$.
  \item Résolution : $y = 2x^3 + 3 \iff x^3 = \frac{y-3}{2} \iff x = \sqrt[3]{\frac{y-3}{2}}$. Donc $f^{-1}(y) = \sqrt[3]{\frac{y-3}{2}}$, $D_{f^{-1}} = \mathbb{R}$.
  \item Dérivée : $(f^{-1})'(y) = \frac{1}{f'(f^{-1}(y))} = \frac{1}{6\left(\sqrt[3]{\frac{y-3}{2}}\right)^2} = \frac{1}{6} \left(\frac{2}{y-3}\right)^{2/3}$ pour $y \neq 3$. En $y=3$, $f^{-1}(3)=0$, $f'(0)=0$, la dérivée n'existe pas (tangente verticale).
  \item Variations : $(f^{-1})'(y) > 0$ pour $y \neq 3$, $f^{-1}$ strictement croissante sur $\mathbb{R}$.
  \item Tracé : $\mathcal{C}_f$ passe par $(0,3)$, tangente horizontale en $(0,3)$. $\mathcal{C}_{f^{-1}}$ passe par $(3,0)$, tangente verticale en $(3,0)$. Symétrie par rapport à $y=x$.
\end{enumerate}
\end{exoresolu}

\begin{attention}[Pièges classiques sur la réciproque]
\begin{itemize}
  \item Confondre $f^{-1}(x)$ (réciproque) et $\frac{1}{f(x)}$ (inverse). \textbf{Jamais} écrire $f^{-1}(x) = 1/f(x)$.
  \item Oublier d'échanger ensemble de définition et ensemble d'arrivée : $D_{f^{-1}} = f(D_f)$.
  \item Appliquer la formule de la dérivée sans vérifier $f'(f^{-1}(y)) \neq 0$.
  \item Tracer la symétrie par rapport à $y=x$ sans respecter l'échelle (repère non orthonormé ou échelles différentes sur axes) : la symétrie n'est visuelle que si les échelles sont identiques.
\end{itemize}
\end{attention}

\coursec{Racines $n$-ièmes et puissances rationnelles}

\begin{definition}[Racine $n$-ième d'un réel positif]
Soit $n \in \mathbb{N}^*$ et $a \in \mathbb{R}_+^*$. La \cle{racine $n$-ième de $a$}, notée $\sqrt[n]{a}$ ou $a^{1/n}$, est l'unique réel \textbf{strictement positif} $x$ tel que $x^n = a$.
Par convention $\sqrt[n]{0} = 0$.
\end{definition}

\begin{propriete}[Résolution de $x^n = a$ dans $\mathbb{R}$]
\begin{itemize}
  \item Si $n$ est \textbf{pair} :
    \begin{itemize}
      \item $a < 0$ : \textbf{aucune solution} dans $\mathbb{R}$.
      \item $a = 0$ : solution unique $x = 0$.
      \item $a > 0$ : \textbf{deux solutions} $x = \sqrt[n]{a}$ et $x = -\sqrt[n]{a}$ (souvent notées $x = \pm \sqrt[n]{a}$).
    \end{itemize}
  \item Si $n$ est \textbf{impair} :
    \begin{itemize}
      \item Pour tout $a \in \mathbb{R}$, \textbf{une unique solution} $x = \sqrt[n]{a}$ (le signe de $x$ est celui de $a$).
    \end{itemize}
\end{itemize}
\end{propriete}

\begin{definition}[Puissance rationnelle]
Pour $a > 0$ et $r = \frac{p}{q} \in \mathbb{Q}$ ($q > 0$), on définit
\[
a^{p/q} = \sqrt[q]{a^p} = (\sqrt[q]{a})^p.
\]
Pour $a=0$, $0^r = 0$ si $r>0$ (indéfini si $r \le 0$).
Pour $a < 0$, $a^{p/q}$ n'est défini dans $\mathbb{R}$ que si $q$ est \textbf{impair} (alors $a^{p/q} = \sqrt[q]{a^p}$).
\end{definition}

\begin{propriete}[Propriétés des puissances rationnelles]
Pour $a,b > 0$ (ou $a,b \in \mathbb{R}$ si les expressions ont un sens) et $r,s \in \mathbb{Q}$ :
\[
a^r a^s = a^{r+s}, \quad \frac{a^r}{a^s} = a^{r-s}, \quad (a^r)^s = a^{rs}, \quad (ab)^r = a^r b^r, \quad \left(\frac{a}{b}\right)^r = \frac{a^r}{b^r}.
\]
\end{propriete}

\begin{methode}[Simplifier une expression à puissances rationnelles]
\begin{enumerate}
  \item Écrire tous les radicaux sous forme de puissances ($ \sqrt[q]{a^p} = a^{p/q} $).
  \item Mettre au même dénominateur les exposants si nécessaire.
  \item Appliquer les règles de calcul sur les puissances.
  \item Revenir éventuellement à la notation radicale si l'énoncé le demande.
  \item \textbf{Vérifier le domaine de définition} (base positive ou exposant à dénominateur impair).
\end{enumerate}
\end{methode}

\begin{exemplebox}[Simplification]
Simplifier $E = \frac{\sqrt[3]{x^2} \cdot \sqrt{x}}{x^{1/6}}$ pour $x > 0$.
$E = \frac{x^{2/3} \cdot x^{1/2}}{x^{1/6}} = x^{2/3 + 1/2 - 1/6} = x^{4/6 + 3/6 - 1/6} = x^{6/6} = x$.
\end{exemplebox}

\begin{attention}[Erreurs fréquentes sur les puissances]
\begin{itemize}
  \item $\sqrt{a^2} = |a|$ et non $a$ (sauf si $a \ge 0$ connu).
  \item $\sqrt[n]{a^n} = a$ si $n$ impair, $= |a|$ si $n$ pair.
  \item $(a^r)^s = a^{rs}$ n'est pas toujours vrai si $a<0$ et dénominateurs pairs (ex: $((-1)^2)^{1/2} = 1 \neq -1 = (-1)^{2 \times 1/2}$). \textbf{Toujours ramener à une base positive ou vérifier la parité des dénominateurs.}
\end{itemize}
\end{attention}

\coursec{Inégalité des accroissements finis}

\begin{propriete}[Inégalité des accroissements finis]
Soit $f$ une fonction \cle{continue} sur $[a;b]$ et \cle{dérivable} sur $]a;b[$.
Si il existe $m, M \in \mathbb{R}$ tels que pour tout $x \in ]a;b[$ :
\[
m \le f'(x) \le M,
\]
alors
\[
m(b-a) \le f(b) - f(a) \le M(b-a).
\]
\end{propriete}

\begin{experience}[Démonstration (Théorème des accroissements finis / MVT)]
D'après le théorème des accroissements finis (Théorème de Lagrange), il existe $c \in ]a;b[$ tel que $f(b)-f(a) = f'(c)(b-a)$.
Comme $m \le f'(c) \le M$ et $b-a > 0$, on multiplie par $b-a$ : $m(b-a) \le f'(c)(b-a) \le M(b-a)$, d'où le résultat.
\end{experience}

\begin{methode}[Encadrer $f(b)-f(a)$ à l'aide de l'inégalité des accroissements finis]
\begin{enumerate}
  \item Vérifier les hypothèses : $f$ continue sur $[a;b]$, dérivable sur $]a;b[$.
  \item Déterminer $f'(x)$.
  \item Trouver des bornes $m$ et $M$ (min et max si possible, ou majorations/minorations simples) de $f'$ sur $]a;b[$.
  \item Appliquer la formule : $m(b-a) \le f(b)-f(a) \le M(b-a)$.
\end{enumerate}
\end{methode}

\begin{exoresolu}[Encadrement de $\ln(1+x)$]
\textbf{Énoncé :} Montrer que pour tout $x > 0$, $\frac{x}{1+x} < \ln(1+x) < x$.
\tcblower
\textbf{Solution :}
Posons $f(t) = \ln(1+t)$ sur $[0;x]$ ($x>0$ fixé).
$f$ est continue sur $[0;x]$, dérivable sur $]0;x[$, $f'(t) = \frac{1}{1+t}$.
Pour $t \in ]0;x[$, $1 < 1+t < 1+x \implies \frac{1}{1+x} < \frac{1}{1+t} < 1$.
Donc $m = \frac{1}{1+x}$, $M = 1$.
Inégalité : $\frac{1}{1+x}(x-0) < \ln(1+x) - \ln(1) < 1 \cdot (x-0)$.
Soit $\frac{x}{1+x} < \ln(1+x) < x$. \qed
\end{exoresolu}

\begin{exoresolu}[Encadrement de $\sin b - \sin a$]
\textbf{Énoncé :} Montrer que $|\sin b - \sin a| \le |b-a|$ pour tous réels $a,b$.
\tcblower
\textbf{Solution :} $f(x)=\sin x$ continue et dérivable sur $\mathbb{R}$, $f'(x)=\cos x$, $|\cos x| \le 1$.
Si $a<b$, $-1 \le f'(x) \le 1$. Inégalité : $-1(b-a) \le \sin b - \sin a \le 1(b-a)$.
Donc $|\sin b - \sin a| \le b-a = |b-a|$. Si $b<a$, symétrie. Si $a=b$, égalité.
\end{exoresolu}

\coursec{Limites usuelles et branches infinies}

\begin{propriete}[Limites usuelles à connaître par cœur]
\begin{center}
\begin{tabular}{|c|c|}
\hline
\rowcolor{popblueL} \textbf{Forme} & \textbf{Limite (quand $x \to 0$)} \\
\hline
$\dfrac{\sin x}{x}$ & $1$ \\
\hline
$\dfrac{\tan x}{x}$ & $1$ \\
\hline
$\dfrac{1 - \cos x}{x^2}$ & $\dfrac{1}{2}$ \\
\hline
$\dfrac{\sqrt{1+x} - 1}{x}$ & $\dfrac{1}{2}$ \\
\hline
$\dfrac{\ln(1+x)}{x}$ & $1$ \\
\hline
$\dfrac{e^x - 1}{x}$ & $1$ \\
\hline
$\dfrac{a^x - 1}{x}$ ($a>0$) & $\ln a$ \\
\hline
\end{tabular}
\end{center}
\textbf{Remarque :} Ces limites valent aussi pour $x \to 0^{\pm}$ et s'utilisent en changement de variable ($u \to 0$).
\end{propriete}

\begin{definition}[Branches infinies et asymptotes]
Soit $f$ définie au voisinage de l'infini ou d'une valeur interdite.
\begin{itemize}
  \item \cle{Asymptote verticale} : $x = x_0$ si $\lim_{x \to x_0^{\pm}} f(x) = \pm \infty$.
  \item \cle{Asymptote horizontale} : $y = \ell$ si $\lim_{x \to \pm\infty} f(x) = \ell \in \mathbb{R}$.
  \item \cle{Asymptote oblique} : $y = ax + b$ ($a \neq 0$) si $\lim_{x \to \pm\infty} \frac{f(x)}{x} = a \in \mathbb{R}^*$ et $\lim_{x \to \pm\infty} (f(x) - ax) = b \in \mathbb{R}$.
  \item \cle{Direction asymptotique} : Si $\lim_{x \to \pm\infty} \frac{f(x)}{x} = a \in \mathbb{R}^*$ mais $\lim_{x \to \pm\infty} (f(x) - ax) = \pm \infty$, la courbe a une direction asymptotique de pente $a$ (parallèle à $y=ax$).
  \item \cle{Branche parabolique} : Si $\lim_{x \to \pm\infty} \frac{f(x)}{x} = \pm \infty$, on étudie $\frac{f(x)}{x^2}$ etc.
\end{itemize}
\end{definition}

\begin{methode}[Recherche systématique des asymptotes en $+\infty$ (idem $-\infty$)]
\begin{enumerate}
  \item Calculer $\lim_{x \to +\infty} f(x) = \ell$.
    \begin{itemize}
      \item Si $\ell \in \mathbb{R}$ : \textbf{asymptote horizontale} $y = \ell$. \textbf{Stop}.
      \item Si $\ell = \pm \infty$ : passer à l'étape 2.
    \end{itemize}
  \item Calculer $a = \lim_{x \to +\infty} \frac{f(x)}{x}$.
    \begin{itemize}
      \item Si $a = 0$ : pas d'asymptote oblique, direction horizontale (déjà traité) ou courbe "verticale" à l'infini.
      \item Si $a = \pm \infty$ : \textbf{direction asymptotique} de pente infinie (branche parabolique ou supérieure). \textbf{Stop}.
      \item Si $a \in \mathbb{R}^*$ : passer à l'étape 3.
    \end{itemize}
  \item Calculer $b = \lim_{x \to +\infty} (f(x) - ax)$.
    \begin{itemize}
      \item Si $b \in \mathbb{R}$ : \textbf{asymptote oblique} $y = ax + b$.
      \item Si $b = \pm \infty$ : \textbf{direction asymptotique} de pente $a$ (pas d'asymptote oblique).
    \end{itemize}
\end{enumerate}
\end{methode}

\begin{exemplebox}[Asymptotes de $f(x) = \frac{x^2+1}{x-1}$]
$D_f = \mathbb{R} \setminus \{1\}$.
$x \to 1^{\pm}$ : $f(x) \to \pm \infty$ $\Rightarrow$ \textbf{Asymptote verticale} $x=1$.
$x \to \pm\infty$ : $f(x) \sim x \to \pm\infty$.
$a = \lim \frac{f(x)}{x} = \lim \frac{x^2+1}{x(x-1)} = 1$.
$b = \lim (f(x) - x) = \lim \frac{x^2+1 - x(x-1)}{x-1} = \lim \frac{x+1}{x-1} = 1$.
\textbf{Asymptote oblique} $y = x + 1$ (commune aux deux infinis).
\end{exemplebox}

\begin{exemplebox}[Direction asymptotique : $f(x) = \sqrt{x^2+x}$]
$D_f = ]-\infty;-1] \cup [0;+\infty[$.
$x \to +\infty$ : $f(x) \sim x$. $a = \lim \frac{\sqrt{x^2+x}}{x} = \lim \sqrt{1+1/x} = 1$.
$b = \lim (\sqrt{x^2+x} - x) = \lim \frac{x}{\sqrt{x^2+x}+x} = \frac{1}{2}$.
\textbf{Asymptote oblique} $y = x + 1/2$ en $+\infty$.
$x \to -\infty$ : $f(x) = |x|\sqrt{1+1/x} = -x\sqrt{1+1/x}$ (car $x<0$).
$a = \lim \frac{f(x)}{x} = \lim -\sqrt{1+1/x} = -1$.
$b = \lim (f(x) + x) = \lim (-x\sqrt{1+1/x} + x) = \lim \frac{-x^2(1+1/x) + x^2}{-x\sqrt{1+1/x} - x} = \lim \frac{-x}{-x(\sqrt{1+1/x}+1)} = \frac{1}{2}$.
\textbf{Asymptote oblique} $y = -x + 1/2$ en $-\infty$.
\end{exemplebox}

\begin{attention}[Pièges sur les asymptotes]
\begin{itemize}
  \item Oublier d'étudier $-\infty$ et $+\infty$ séparément (asymptotes différentes possibles).
  \item Calculer $b = \lim (f(x)-ax)$ sans avoir vérifié que $a$ est fini et non nul.
  \item Pour les fonctions irrationnelles ($\sqrt{\dots}$), attention au signe de $x$ quand on factorise : $\sqrt{x^2} = |x|$.
  \item Confondre "direction asymptotique" et "asymptote oblique". Si $b$ est infini, il n'y a \textbf{pas} d'asymptote oblique, seulement une direction.
\end{itemize}
\end{attention}

\coursec{Étude complète de fonctions et tracé de courbes}

\begin{methode}[Les 9 étapes de l'étude complète (la "Grande Étude")]
\begin{enumerate}
  \item \textbf{Ensemble de définition $D_f$} : Zéros de dénominateurs, radicaux d'indice pair ($\ge 0$), log ($>0$), etc.
  \item \textbf{Parité / Périodicité} : Pair ($f(-x)=f(x)$), Impair ($f(-x)=-f(x)$), Période $T$ ($f(x+T)=f(x)$). Permet de réduire l'étude.
  \item \textbf{Limites aux bornes de $D_f$} : Bornes finies (valeurs interdites $\to$ asymptotes verticales) et $\pm\infty$ (asymptotes H/O, directions).
  \item \textbf{Dérivée $f'(x)$} : Calculer et simplifier (factoriser !). Déterminer $D_{f'}$ (souvent $= D_f$).
  \item \textbf{Signe de $f'(x)$} : Tableau de signes $\to$ variations de $f$.
  \item \textbf{Tableau de variations} : Avec les limites aux bornes, les valeurs remarquables ($f(0)$, extremums, zéros).
  \item \textbf{Asymptotes et positions relatives} : Équations des asymptotes. Position de la courbe par rapport à son asymptote oblique (signe de $f(x)-(ax+b)$).
  \item \textbf{Points particuliers} : Intersections avec les axes ($f(x)=0$, $f(0)$), points d'inflexion ($f''$), tangentes remarquables.
  \item \textbf{Tracé de la courbe $\mathcal{C}_f$} : Repère orthonormé, même échelle, points précis, asymptotes en pointillés, flèches aux branches infinies.
\end{enumerate}
\end{methode}

\begin{exoresolu}[Étude complète : Fonction rationnelle $f(x) = \frac{x^2 - 2x - 3}{x+1}$ (un piège classique)]
\textbf{Énoncé :} Étudier $f$ et tracer $\mathcal{C}_f$.
\tcblower
\textbf{Solution :}
\begin{enumerate}
  \item $D_f = \mathbb{R} \setminus \{-1\}$.
  \item \textbf{Réflexe : factoriser le numérateur avant tout calcul.} $-1$ est racine de $x^2-2x-3$ (car $1+2-3=0$), donc $x^2-2x-3=(x+1)(x-3)$ et, pour tout $x \neq -1$ :
  \[ f(x) = \frac{(x+1)(x-3)}{x+1} = x - 3. \]
  \item \textbf{Limites :} $\lim\limits_{x \to -1} f(x) = \lim\limits_{x \to -1} (x-3) = -4$ : la limite en $-1$ est \emph{finie}, il n'y a \textbf{pas d'asymptote verticale}. En $\pm\infty$ : $\lim f(x) = \pm\infty$.
  \item \textbf{Dérivée et variations :} $f'(x) = 1 > 0$ sur $]-\infty;-1[$ et sur $]-1;+\infty[$ : $f$ est strictement croissante sur chacun de ces intervalles.
  \item \textbf{Tableau :}
  \begin{center}
  \begin{tabular}{|c|ccc||ccc|}
  \hline
  $x$ & $-\infty$ & & $-1$ & $-1$ & & $+\infty$ \\
  \hline
  $f'(x)$ & & $+$ & & & $+$ & \\
  \hline
  $f(x)$ & $-\infty$ & $\nearrow$ & $-4$ & $-4$ & $\nearrow$ & $+\infty$ \\
  \hline
  \end{tabular}
  \end{center}
  ($-4$ est une limite, pas une valeur : $f(-1)$ n'existe pas.)
  \item \textbf{Tracé :} $\mathcal{C}_f$ est la droite d'équation $y = x-3$ \textbf{privée du point} $A(-1;-4)$, que l'on représente par un petit cercle vide. Pas d'asymptote : la droite $y=x-3$ \emph{est} la courbe.
\end{enumerate}
\textbf{Moralité :} si le numérateur et le dénominateur ont une racine commune, simplifier d'abord ; sinon on « trouve » à tort une asymptote verticale.
\end{exoresolu}

\begin{exoresolu}[Étude complète : Fonction irrationnelle $f(x) = x\sqrt{4-x^2}$]
\textbf{Énoncé :} Étudier $f$ sur son ensemble de définition.
\tcblower
\textbf{Solution :}
\begin{enumerate}
  \item $D_f : 4-x^2 \ge 0 \iff x \in [-2; 2]$.
  \item $f(-x) = -x\sqrt{4-x^2} = -f(x)$ : \textbf{fonction impaire}. Étudier sur $[0;2]$, déduire sur $[-2;0]$ par symétrie centrale $O$.
  \item \textbf{Limites :} $x \to 2^-$, $f(x) \to 0$. $x \to -2^+$, $f(x) \to 0$. Pas d'asymptote verticale (domaine borné). Pas d'infini.
  \item \textbf{Dérivée sur $]0;2[$ :} $f'(x) = \sqrt{4-x^2} + x \frac{-2x}{2\sqrt{4-x^2}} = \frac{4-x^2 - x^2}{\sqrt{4-x^2}} = \frac{4-2x^2}{\sqrt{4-x^2}} = \frac{2(2-x^2)}{\sqrt{4-x^2}}$.
  \item \textbf{Signe de $f'$ :} Dénominateur $>0$. Numérateur $2(2-x^2) \ge 0 \iff x^2 \le 2 \iff x \in [0;\sqrt{2}]$.
    $f' \ge 0$ sur $[0;\sqrt{2}]$, $f' \le 0$ sur $[\sqrt{2};2]$.
  \item \textbf{Variations sur $[0;2]$ :} Croissante sur $[0;\sqrt{2}]$, décroissante sur $[\sqrt{2};2]$.
    Max en $x=\sqrt{2}$ : $f(\sqrt{2}) = \sqrt{2}\sqrt{2} = 2$.
    $f(0)=0$, $f(2)=0$.
  \item \textbf{Tableau complet (parité) :}
  \begin{center}
  \begin{tabular}{|c|ccccccccc|}
  \hline
  $x$ & $-2$ & & $-\sqrt{2}$ & & $0$ & & $\sqrt{2}$ & & $2$ \\
  \hline
  $f'(x)$ & & $-$ & $0$ & $+$ & $+$ & $+$ & $0$ & $-$ & \\
  \hline
  $f(x)$ & $0$ & $\searrow$ & $-2$ & $\nearrow$ & $0$ & $\nearrow$ & $2$ & $\searrow$ & $0$ \\
  \hline
  \end{tabular}
  \end{center}
  \item \textbf{Tangentes aux bornes :} quand $x \to 2^-$, le numérateur $4-2x^2 \to -4$ et le dénominateur $\sqrt{4-x^2}\to 0^+$, donc $f'(x)\to -\infty$ : tangente verticale en $A(2;0)$ et, par imparité, en $A'(-2;0)$.
  \item \textbf{Tracé :} Courbe symétrique par rapport à $O$, passant par $(-2;0)$, $(-\sqrt{2};-2)$, $(0;0)$, $(\sqrt{2};2)$ et $(2;0)$. Tangentes verticales aux extrémités.
\end{enumerate}
\end{exoresolu}

\begin{exoresolu}[{Étude complète : Fonction trigonométrique $f(x) = \sin x + \cos x$ sur $[0; 2\pi]$}]
\textbf{Énoncé :} Étudier $f$ et tracer $\mathcal{C}_f$.
\tcblower
\textbf{Solution :}
\begin{enumerate}
  \item $D_f = [0; 2\pi]$.
  \item $f(x) = \sqrt{2}\sin(x+\pi/4)$ (transformation de $\cos x + \sin x$ en une seule sinusoïde). Période $2\pi$. On étudie sur une période.
  \item \textbf{Limites :} Bornes finies : $f(0)=1$, $f(2\pi)=1$.
  \item \textbf{Dérivée :} $f'(x) = \cos x - \sin x = \sqrt{2}\cos(x+\pi/4)$.
  \item \textbf{Signe de $f'$ sur $[0;2\pi]$ :} $\cos(x+\pi/4) \ge 0 \iff x+\pi/4 \in [-\pi/2; \pi/2] + 2k\pi \iff x \in [-3\pi/4; \pi/4] + 2k\pi$.
    Sur $[0;2\pi]$ (avec $-3\pi/4 + 2\pi = 5\pi/4$) : $f' \ge 0$ sur $[0; \pi/4] \cup [5\pi/4; 2\pi]$, $f' \le 0$ sur $[\pi/4; 5\pi/4]$.
  \item \textbf{Variations :}
    Croissante $[0;\pi/4]$, Max $f(\pi/4)=\sqrt{2}$.
    Décroissante $[\pi/4; 5\pi/4]$, Min $f(5\pi/4)=-\sqrt{2}$.
    Croissante $[5\pi/4; 2\pi]$.
  \item \textbf{Zéros :} $\sin x + \cos x = 0 \iff \tan x = -1 \iff x = 3\pi/4, 7\pi/4$.
  \item \textbf{Tableau :}
  \begin{center}
  \begin{tabular}{|c|ccccccccccc|}
  \hline
  $x$ & $0$ & & $\pi/4$ & & $3\pi/4$ & & $5\pi/4$ & & $7\pi/4$ & & $2\pi$ \\
  \hline
  $f'(x)$ & & $+$ & $0$ & $-$ & & $-$ & $0$ & $+$ & & $+$ & \\
  \hline
  $f(x)$ & $1$ & $\nearrow$ & $\sqrt{2}$ & $\searrow$ & $0$ & $\searrow$ & $-\sqrt{2}$ & $\nearrow$ & $0$ & $\nearrow$ & $1$ \\
  \hline
  \end{tabular}
  \end{center}
  \item \textbf{Tracé :} Onde sinusoïdale d'amplitude $\sqrt{2}$, décalée de $-\pi/4$.
\end{enumerate}
\end{exoresolu}

\begin{aretenir}[Checklist finale pour le Bac : Étude de fonctions]
\begin{itemize}
  \item $\square$ \textbf{TVI} : Énoncer les 3 hypothèses (continuité, intervalle, $c$ entre images). Conclure "Existence". Ajouter "stricte monotonie" pour "Unicité".
  \item $\square$ \textbf{Encadrement $\alpha$} : Balayage (tableau de valeurs) ou dichotomie. Précision $10^{-2}$ : donner l'intervalle $[a;b]$ tel que $b-a=0,01$ et $f(a)f(b)<0$.
  \item $\square$ \textbf{Bijection / Réciproque} : Prouver bijection (continuité + stricte mono + ensemble d'arrivée = image). Déterminer $D_{f^{-1}} = f(D_f)$. Formule dérivée $(f^{-1})'(y) = 1/f'(f^{-1}(y))$. Tracé par symétrie $y=x$ (échelles identiques !).
  \item $\square$ \textbf{Racines / Puissances} : Résoudre $x^n=a$ (discussion parité de $n$). Simplifier $a^{p/q}$ (base $>0$ ou $q$ impair). $\sqrt{a^2}=|a|$.
  \item $\square$ \textbf{Accroissements finis} : Hypothèses (CV sur $[a,b]$, dérivable sur $]a,b[$). Borner $f'$. Appliquer $m(b-a) \le f(b)-f(a) \le M(b-a)$.
  \item $\square$ \textbf{Limites usuelles} : $\sin x/x$, $\tan x/x$, $(1-\cos x)/x^2$, $(\sqrt{1+x}-1)/x$, $\ln(1+x)/x$, $(e^x-1)/x$. Changer de variable pour faire apparaître $u\to 0$.
  \item $\square$ \textbf{Asymptotes} : Algorithme rigoureux (1. Limite finie $\to$ H ; 2. Limite $f/x = a$ fini non nul ; 3. Limite $f-ax = b$ fini $\to$ O, infini $\to$ Direction). Traiter $+\infty$ et $-\infty$ séparément. Factoriser $\sqrt{x^2}=|x|$.
  \item $\square$ \textbf{Grande Étude} : Ordre immuable : $D_f \to$ Parité $\to$ Limites $\to f' \to$ Signe $f' \to$ TV $\to$ Asymptotes/Position $\to$ Points remarquables $\to$ Tracé.
  \item $\square$ \textbf{Rédaction} : Phrases complètes. "Donc", "Par conséquent". Tableau de variations au format standard (flèches, valeurs, doubles barres pour interdits). Courbe soignée (règle, équerre, pointillés asymptotes).
\end{itemize}
\end{aretenir}

\begin{savaistu}[Le TVI et la dichotomie : l'algorithme de base de l'informatique]
Le théorème des valeurs intermédiaires est la justification mathématique de la \textbf{méthode de dichotomie} (ou recherche binaire), algorithme fondamental en informatique pour trouver un zéro de fonction ou trier/chercher dans un tableau ordonné.
\begin{enumerate}
  \item On a un intervalle $[a,b]$ où $f(a)f(b) < 0$ (signe contraire).
  \item On calcule le milieu $m = (a+b)/2$.
  \item Si $f(m)=0$, c'est gagné. Sinon, on regarde le signe de $f(m)$.
  \item Si $f(a)f(m) < 0$, la solution est dans $[a,m]$, sinon dans $[m,b]$.
  \item On recommence jusqu'à la précision voulue.
\end{enumerate}
À chaque étape, l'intervalle est divisé par 2. Pour une précision $10^{-n}$, il faut environ $n \times \log_2(10) \approx 3,3n$ itérations. C'est l'efficacité algorithmique $O(\log n)$.
Au Cameroun, cette méthode est utilisée dans les logiciels de calcul scientifique (Scilab, Python \texttt{scipy.optimize.bisect}) pour résoudre des équations issues de la modélisation (écoulement dans les barrages, dimensionnement de poutres, traitement de signaux MTN/Orange).
\end{savaistu}

\findecours

\end{document}
